\chapter{Fluctuation integration and normalization on block regions}\label{ch17:chapter}

In \cref{ch15:chapter,ch16:chapter} we set up the averaging steps of one renormalization group step on a block region. There, the coordinates attached to a union $\Omega$ of coarse blocks are analytic variables, and the coordinates on the fixed complement of $\Omega$ are frozen at real values chosen earlier. Those sections constructed the conditional constrained minimizers of the classical problem on $\Omega$, proved a lower bound on the constrained Hessian independent of the shape of $\Omega$, and controlled the effect of a small perturbation of the classical objective. The goal of this section is to integrate out the normalized fluctuation variables of such a step. We prove that the resulting effective interaction differs from the corresponding torus quantity only by terms that decay exponentially with the distance to the fixed complement, and we keep track of the normalization constant $\Lambda$ by a separate argument. All estimates are uniform in the shape and topology of $\Omega$, the number of its components, the preceding depth, the number of renormalization group steps already performed, and the values of the frozen coordinates.

\Cref{ch17:sec:setting} fixes the notation and the coefficient calculus. We then compare the conditional Hessians and Haar densities with their torus counterparts (\cref{ch17:sec:hessians}) and derive the conditional fluctuation potential and its concentration bounds (\cref{ch17:sec:potential}). We obtain the effective interaction by a cluster expansion (\cref{ch17:sec:expansion}) and control $\Lambda$ by a separate comparison of real integrals (\cref{ch17:sec:normalization}). Next, we choose the parameters and prove the finite iteration (\cref{ch17:thm:controlled-kernel}). Finally, we sum the truncation errors for couplings in $(0,\Gamma]$ with any finite $\Gamma$ (\cref{ch17:sec:endpoint}).

We use the norms and the dependence of constants of \cref{ch01:chapter}, the concentration estimates of \cref{ch02:thm:reflecting-concentration}, and the triangular change of variables and the finite graph inequality of \cref{ch09:thm:sequential-density,ch09:thm:hard-core-log}. From \cref{ch15:chapter} we use the norm with boundary decay, the network estimates, and the substitution identities, and from \cref{ch16:chapter} the conditional domains, image bounds, perturbation theorem, and column bounds. The results are used in \cref{ch18:chapter}, which compares the integrals over the regions of the large-field decomposition with the analytic integrals constructed here, in \cref{ch19:chapter}, which proves the exponential-moment estimates for blocked actions, and in \cref{ch20:chapter}, which chooses the parameters.

All statements of this section concern analytic functions on stated domains and real integrals over stated products. We prove in \cref{ch18:chapter} that the configurations of the large-field decomposition lie in these domains.

\paragraph{Heuristics.}
We first explain the orders of magnitude. We write the free logarithmic coordinates of one intermediate scale as $z=z^*+gK^{-1/2}w$, where $z^*$ is the conditional constrained minimizer, $K$ the constrained Hessian, $g$ the coupling, and $w$ the normalized fluctuation variable. Since the classical action enters the density divided by $g^2$, and its expansion around the minimizer has no linear term and a quadratic term equal to $\abs w^2/2$, the remaining classical terms are
\[
 g^{-2}\,O\bigl((z-z^*)^3\bigr)=O(g\abs w^3).
\]
The Haar disintegration density, the logarithmic determinants of the higher scales, and the incoming effective interaction are not divided by $g^2$, and their dependence on $w$ starts with a term of order $g\abs w$. On fluctuation balls of radius $p\ge1$, the interaction is therefore small when $g[p^3+(1+H)p]$ is small, where $H$ bounds the incoming effective interaction; this is the smallness condition of \cref{ch17:thm:carried-networks}. The terms of first order in $g$ are odd under $w\mapsto-w$. Their expectations under a symmetric reference measure vanish, and the fluctuation correction to the effective interaction is of order $g^2$. The conditional measure with the frozen coordinates held fixed is interacting and need not be symmetric. For this reason, we bound the means of the fluctuation variables separately in \cref{ch17:lem:ambient-white-comparison}.

We control the normalization constant by two identities. First, a missing scalar Gaussian coordinate at an intermediate scale with coupling $g_i$ removes one factor $g_i$ from the density. This contributes $\log(1/g_i)$ to the difference between the regional and the torus normalization. Second, fixing coordinates at values $c$ uses
\[
 e^{\Lambda-R(v,c)}=e^{\Lambda-R(0,c)}e^{-[R(v,c)-R(0,c)]},
\]
which moves the value at $v=0$ into $\Lambda$ and leaves a centered function. Since the norms of the effective interaction vanish on constants, they cannot control $\Lambda$, and we estimate $\Lambda$ by comparing real integrals.

Finally, fixing the coordinates of a boundary layer affects the boundary part of the effective interaction in two ways. First, the norm of the old boundary discrepancy is bounded by $2e^{-\awall(s_{\mathrm{er}}-1)}$ times its previous norm, where $s_{\mathrm{er}}$ is the width of the frozen layer. Second, fixing coordinates at nonzero values creates a new mixed term from the torus interaction, of size at most a constant times $E\beta_{\mathrm{rec}}/\rho$, where $E$ bounds the torus interaction and $\beta_{\mathrm{rec}}$ bounds the frozen values. The first estimate needs a wide layer, the second needs small frozen values, and neither replaces the other.

\paragraph{Attribution.}
The architecture used here, namely constrained minimization in a gauge-fixed chart, separation of a background field from normalized fluctuations, and integration of the fluctuations by a cluster expansion, is Ba\l aban's renormalization group for lattice gauge theories \cite{Balaban1988RGII,Balaban1988Convergent}. We attribute the individual tools where they are used.

\paragraph{Order of the proofs.}
No conclusion is used before it is proved. In particular, \cref{ch17:lem:ambient-white-comparison} uses only the coefficient, product, and curvature bounds of \cref{ch17:lem:complete-white-potential}, the density sum \eqref{ch17:eq:density-sum}, and \cref{ch02:thm:reflecting-concentration}. It determines the constant $C_1$ in a hypothesis of \cref{ch17:lem:complete-white-potential}, and the logarithmic radii of \cref{ch17:thm:controlled-kernel} enter it only as a sufficient condition for an inclusion.

\section{Setting and coefficient calculus}\label{ch17:sec:setting}

\paragraph{Geometry and labels.}
Fix an odd integer $M\ge3$, an integer $J\ge1$, and $B=M^J$. We consider a composition of $s$ averaging steps, $1\le s\le J$, from a lattice of side $M^sm$ to a lattice of side $m\ge4$, together with an arbitrary union $\Omega$ of blocks of the output lattice. When $s=J$, we speak of a full step; it is one renormalization group step. The intermediate scales are numbered $i=0,\ldots,s-1$. A label is a scalar coordinate: it specifies a direction, a color, an intermediate scale (or the output lattice), and the lattice on which the coordinate was introduced, and it has an assigned position. A fluctuation block is the three-component vector $w_{i,e}\in\R^3$ of normalized fluctuation variables attached to one free link $e$ at one intermediate scale $i$; it is a single variable of the reference measures below, not three independent ones. Scalar labels and fluctuation blocks are counted separately.

The labels attached to $\Omega$ are unfixed: their coordinates are analytic variables. These are the coarse background variable $v=\Log c$ on the output lattice and the free logarithmic coordinates $z_i$ at the intermediate scales. We write $\Omega^c$ for the fixed complement of the unfixed positions, $W$ for the boundary (the interface between the unfixed positions and $\Omega^c$), and $\Omega^c_{L_J}$ for the fixed enlargement of $\Omega^c$ determined by the construction stencils, as in \cref{ch15:lem:wall-networks}. The fixed complement is denoted $H$ in \cref{ch15:chapter} and \cref{ch16:chapter}: in their notation, $\Omega^c=H$ and $\Omega^c_{L_J}=H_{L_J}$. In this section the letter $H$ denotes instead the bound $E+\mathcal X^\rho(f)$ on the incoming effective interaction introduced in \cref{ch17:lem:decorated-white-primitives}. All distances are graph distances on the refined position graph of \cref{ch15:sec:norms}, in units of the output lattice, and network lengths use its periodic $\ell^1$ metric, with internal vertices allowed; positions of labels introduced on finer lattices are embedded in these units. We write $d_g$ and $\Tree_g$ for this distance and network length, and $d_E$, $\Tree_E$ for the periodic Euclidean distance and network length of \cref{ch01:def:analytic-norms}. Then
\[
 d_E\le d_g\le\sqrt3\,d_E,\qquad\Tree_E\le\Tree_g\le\sqrt3\,\Tree_E,
\]
since $\abs y_2\le\abs y_1\le\sqrt3\abs y_2$ for $y\in\R^3$, applied to each periodic image and to each segment $[p,q]$ of a network, with $y=p-q$. Consequently, a coefficient bound stated with Euclidean distances and network lengths, at spatial exponent $c$ and boundary exponent $c'$, holds with the graph quantities at exponents $c/\sqrt3$ and $c'/\sqrt3$. Conversely, a bound stated with the graph quantities implies the Euclidean bound at the same exponents. Every imported Euclidean coefficient bound is converted in this way before the rates of this section are chosen, and the fixed factor $\sqrt3$ is included in the geometric constants. Throughout this section, $\dist$ and $\Tree$ denote $d_g$ and $\Tree_g$, as in \cref{ch15:sec:norms}. The torus seminorm $\mathcal B_{\rho,a_{\mathbb T}}$ keeps the Euclidean convention of \cref{ch01:def:analytic-norms}. For an unfixed label $i$ we write $d_\Omega(i)$ for its distance to $\Omega^c$. We write $\Omega'$ for the set of unfixed scalar output labels. Separately, let $\partial_{\mathrm{blk}}\Omega$ be the set of anchors of retained final blocks that share a face with a final block of the fixed complement. This is a geometric set on the output lattice, counted once per anchor, regardless of directions, colors, or the presence of an output link. In particular it includes an isolated retained block, even though that block has no internal coarse bond and contributes no label to $\Omega'$. It is empty when $\Omega$ or its fixed complement is empty. We assume the label-density bound
\[
 \#\{y:R\le\dist(x,y)<R+1\}\le v_J(R+2)^3,
\]
where $v_J$ accounts for directions, colors, intermediate scales, and the lattices on which labels are introduced.

\paragraph{Frozen coordinates.}
The frozen coordinates are written $b=(\xi,\Rec)$. Here $\xi$ consists of the coordinates frozen at the end of the preceding step, written in the lattice on which they were frozen, and $\Rec$ consists of the older frozen coordinates, each on the lattice on which it was frozen. Throughout this section the frozen coordinates are real and fixed, and only the unfixed background and free variables are differentiated or continued to complex values. The supremum over the frozen values is taken after every coefficient sum and every degree sum. We claim no derivative of a perturbed quantity in $\xi$ or $\Rec$.

\paragraph{Norms.}
For a function $f(v;\Rec)=\sum_\nu f_\nu(\Rec)v^\nu$ of the unfixed labels, a radius $\rho>0$, a positive profile $\phi\le\rho$, a spatial exponent $\asp>0$, and a boundary exponent $\awall>0$, we use the coefficient norm with boundary decay of \cref{ch15:def:regional-norm}, which is the graph-metric variant of \eqref{ch01:eq:regional-norm}:
\[
 \mathcal X_\Omega^\phi(f)=\sup_{\Rec\ \mathrm{real}}\rho^{-2}
 \sum_{k\ge1}\sup_{i\in\Omega}e^{\awall d_\Omega(i)}
 \sum_{\abs\nu=k}\frac{\nu_i}{k}\abs{f_\nu(\Rec)}e^{\asp\Tree(\supp\nu)}\phi^\nu .
\]
Here $\Tree(S)=\Tree_g(S)$ is the $\ell^1$ length of a shortest connected network containing the positions in $S$, internal vertices allowed, and $d_\Omega$ is the graph distance, as in \cref{ch15:def:regional-norm}. The supremum over the root $i$ precedes the degree sum, and the supremum over the frozen values follows it. Constants and functions on an empty set of unfixed labels have norm zero, and the degree-one term is included. When $\Omega^c$ is empty, a boundary difference has zero nonconstant part and norm zero. We write $\mathcal X^\rho$ for the norm with the uniform profile $\phi\equiv\rho$, and $\mathcal X^{\mathrm{cl}}$ for the classical version, which has its own radius $\rho_{\mathrm{cl}}$, profiles, and exponents. For a torus function $r$ with $Dr(0)=0$ we use
\[
 \Anorm{\rho}{a}{r}=\sum_{k\ge2}\frac{\rho^{k-2}}{k!}\norm{D^kr(0)}_{a,\mathrm{root}}
\]
from \cref{ch01:def:analytic-norms}. This seminorm does not control the value $r(0)$, which is unrestricted. The three profiles are
\[
 \phi(i)=\begin{cases}\rho_\partial,&d_\Omega(i)<w_{\mathrm{bd}},\\ \rho,&d_\Omega(i)\ge w_{\mathrm{bd}},\end{cases}
 \qquad \WideProfile=\sqrt{K_*}\phi,\qquad \InflatedProfile=K_*\phi,
\]
where $K_*\ge2$ is a fixed analytic-radius enlargement factor, $0<\rho_\partial\le\rho$, and $w_{\mathrm{bd}}$ is the width of the boundary layer with the smaller radius. We call $\InflatedProfile$ the outer profile, $\WideProfile$ the intermediate profile, and $\phi$ the working profile. A bound on a coefficient tensor is rooted at an external index if it holds with that index fixed and all other indices summed, where for an input index the volume-normalized column factor $a_{\mathrm{out}}^3/a_{\mathrm{in}}^3$ of \cref{ch01:def:analytic-norms} is included. A coefficient bound with each external index fixed in turn is a bound rooted at every external index. All coefficient bounds concern absolute Taylor coefficients; we never substitute for them a bound on values of a holomorphic function or on an evaluated Hessian.

\paragraph{Admissible data.}
We work on the products of \cref{ch16:lem:admitted-products,ch16:lem:strict-two-radius-images}: the conditional solution neighborhoods of the finitely many stages of the conditional minimization, the smaller neighborhoods of the higher parameters, and the product of sequential free coordinates contained in all of them. On these products, the images under the absolute reconstruction majorant satisfy strict bounds for the classical and quantum radii. A coarse background, a frozen value, or a product of fluctuation balls is called admissible if it lies in these products and satisfies these strict image bounds and the inequality on the conditional forces in \cref{ch16:lem:admitted-products}.

\paragraph{Dependence of constants.}
Unless stated otherwise, every constant in this section ($C_J$, $c_J$, $\lambda_J$, radii, and bounds) has the dependence described in item (3) of \cref{ch01:conv:constants}. It depends only on $M$, $J$, the fixed larger analytic domains and their radius ratios, the spatial and boundary exponents, the label density $v_J$, and the fixed torus and perturbation bounds. It does not depend on $L$, $\gamma$, $m$, $N$, $K$, the preceding depth $d_0$, the number of renormalization group steps already performed, the shape or topology of $\Omega$, the number of components of $\Omega$ or $\Omega^c$, or the values of the frozen coordinates. It is also independent of the widths chosen later (the width of the boundary layer, the width $s_{\mathrm{er}}$ of the frozen layer, the widths of protected neighborhoods, and the size of the cells of the large-field decomposition): every constant is evaluated on the fixed larger domains before these widths are chosen. We assert no bound uniform as $J\to\infty$. Constants in statements about summed errors have the dependence described in item (4) of \cref{ch01:conv:constants}.

\paragraph{Density sums and boundary transport.}
Spatial sums use the label density in units of the output lattice. For $b>0$ we use
\begin{equation}\label{ch17:eq:density-sum}
 S_J(b)=Cv_J\sum_{l\ge0}(l+1)^3e^{-bl},\qquad
 \sup_x\sum_y e^{-b\dist(x,y)}\le S_J(b).
\end{equation}
The constant $C$ accounts for directions, colors, and the finitely many intermediate scales. This is the periodic subset estimate of \cref{ch15:lem:wall-networks}; restricting any of the sets of labels only decreases the sum. Kernels between different meshes keep the volume-normalized column factors of \cref{ch01:def:analytic-norms}. For an unfixed root $x$ and a position $u$ in $\Omega^c_{L_J}$ we have
\begin{equation}\label{ch17:eq:wall-transport}
 d_\Omega(x)\le d_\Omega(u)+\dist(x,u)+C_J,
\end{equation}
where $d_\Omega(u)=0$ if $u\in \Omega^c$ and $C_J$ accounts for the fixed enlargement. The weights for network length and for distance to $\Omega^c$ are both bounded by one fixed part of the available spatial exponent. The fixed stencil offset in \eqref{ch17:eq:wall-transport} contributes a factor $e^{bC_J}$ to the coefficient bound and no factor depending on a later width. We always use the version of \cref{ch15:lem:wall-networks} with every external index rooted. It performs all contractions with this single allocation of the exponent, independently of the number of vertices and of the Taylor order, instead of reducing an exponent at each contraction.

For a proper nonempty block region, we also use the boundary sum
\begin{equation}\label{ch17:eq:boundary-sum}
 \sum_{u\ \mathrm{unfixed}}e^{-c\,d_\Omega(u)}\le C_J(c)\abs{\partial_{\mathrm{blk}}\Omega},\qquad c>0.
\end{equation}
Suppose first that $\Omega$ and its fixed complement are nonempty. A shortest route from an unfixed position to the fixed complement crosses an exposed face of a retained final block. The anchor of that block belongs to $\partial_{\mathrm{blk}}\Omega$ and is within a fixed distance $C_{\mathrm{face}}$ of the crossing point, in output units. The constant includes the fixed convention for link midpoints and endpoints. Thus each unfixed label $u$ can be assigned to an anchor $a(u)\in\partial_{\mathrm{blk}}\Omega$ with
\[
 \dist(u,a(u))\le d_\Omega(u)+C_{\mathrm{face}}.
\]
The density sum \eqref{ch17:eq:density-sum}, applied also at these anchor positions, gives
\[
 \sum_u e^{-c d_\Omega(u)}
 \le e^{cC_{\mathrm{face}}}
       \sum_{a\in\partial_{\mathrm{blk}}\Omega}\sum_u e^{-c\dist(u,a)}
 \le e^{cC_{\mathrm{face}}}S_J(c)\abs{\partial_{\mathrm{blk}}\Omega}.
\]
An anchor need not itself carry a variable; the same density bound applies at every position of the refined graph, after a fixed enlargement of $v_J$ if needed. This proves \eqref{ch17:eq:boundary-sum}, without a regularity or component assumption on the boundary. If $\Omega$ is empty there is no label to sum. If its fixed complement is empty, boundary differences vanish and the bound is used only for those zero terms.

We also record the comparison with the interface labels used in the normalization count. For a fixed construction width $w_{\mathrm{band}}\ge1$, let $N_{\mathrm{band}}$ count the scalar labels of the regional step whose positions are within that width of $W$, and set it to zero when $W$ is empty. Then
\begin{equation}\label{ch17:eq:block-boundary-count}
 \abs{\partial_{\mathrm{blk}}\Omega}\le N_{\mathrm{band}}
       \le C_J\abs{\partial_{\mathrm{blk}}\Omega}.
\end{equation}
For the upper bound, a position within $w_{\mathrm{band}}$ of $W$ is within $w_{\mathrm{band}}+C_{\mathrm{face}}$ of a boundary-block anchor. The label-density bound on this fixed-radius ball proves the claim. For the lower bound, choose one exposed face of each boundary block and a finest-scale $M$-block touching that face. There is a finest-scale free edge on the exposed face of this smaller block. Indeed, for the coordinate-ordered tree of \cref{ch15:def:internal-coordinates}, edges in direction $1$ are tree edges only on the line where coordinates $2$ and $3$ equal those of the root, and edges in direction $2$ are tree edges only in the plane where coordinate $3$ equals that of the root. On a face normal to direction $1$, take an edge in direction $2$ with coordinate $3$ different from that of the root; on a face normal to direction $2$ or $3$, take an edge in direction $1$ off the root line. Such edges exist since $M\ge3$. Each lies inside a finest-scale block, so it is not an eliminated edge, which joins two such blocks. Its midpoint is within one input mesh unit, hence within $1/B_s\le1$, of a link of $W$. One fixed color component of this edge is therefore counted by $N_{\mathrm{band}}$. The choices for distinct final blocks are distinct scalar labels. This gives the lower bound, including isolated final blocks. All constants depend only on $M,J,v_J$ and the fixed construction width, not on any later width of an excluded or frozen neighborhood.

\paragraph{Two series.}
For $0\le q<1$ and integers $j\ge0$ we have
\begin{equation}\label{ch17:eq:factorial-series}
 \sum_{k\ge0}(k+1)(k+2)q^k=\frac2{(1-q)^3},\qquad
 \sum_{k\ge j}(k)_jq^{k-j}=\frac{j!}{(1-q)^{j+1}},
\end{equation}
where $(k)_j=k(k-1)\cdots(k-j+1)$. They follow by differentiating $\sum_{k\ge0}q^k=(1-q)^{-1}$ twice and $j$ times, respectively, on compact subintervals of $[0,1)$. In particular a tensor of order $k+2$ with factorial bound $C^{k+2}(k+2)!$ becomes, after division by $k!$, a summable series of degree-$k$ Hessian coefficients. Additional fixed external indices contribute fixed powers of $k+1$, which are absorbed by passing to one strictly smaller auxiliary product. No smaller radius is chosen at a later renormalization group step.

\paragraph{Coefficient resolution.}
The following lemma makes precise how the connected networks of coefficients behave under the operations used below. It combines the calculus of coefficientwise majorants of \cref{ch02:lem:coefficient-algebra} with \cref{ch15:def:regional-norm}. For a tensor $V$ with output index $o(V)$ we write $v_i$ for the cube of the mesh of the lattice of an index $i$, and $\abs V_b$ for the maximum of its output-row norm and its volume-normalized input-column norms at spatial exponent $b$, as in \cref{ch01:def:analytic-norms}.

\begin{lemma}\label{ch17:lem:coefficient-resolution}
Consider a finite family of typed analytic coefficient tensors. Each tensor is supplied with a coefficient resolution whose total variation, weighted by the exponentials of the lengths of its assigned connected networks, satisfies the bounds of \cref{ch01:def:analytic-norms} with each external index fixed in turn on one common larger product. Consider a sum of externally rooted connected contraction diagrams built from these tensors by substitution, contraction at shared indices (represented by equality tensors), finitely many Schur complements of positive constrained free matrices, and positive implicit or Neumann series. Products of factors with disjoint fluctuation supports and no shared index are not connected; they form separate connected components, each with its own network. Assume that every substitution $F\circ G$ is made on a product whose absolute image majorant for $G$ lies strictly inside the coefficient domain of $F$, that all diagram majorants, including the mesh conversions of part (4), converge absolutely there, and that the inverses and Schur operations satisfy the imported uniform gap and coefficient Neumann conditions on that product. Then the following hold.
\begin{enumerate}
\item The sum has an exact coefficient resolution indexed by the occurrences of its constituent tensors, its contracted indices, and its derivative partitions. In each connected diagram component, the network of a resolved occurrence is the union of the assigned networks of its constituent occurrences, joined at the shared indices. This union is connected, contains every external position of that component, and
\[
 \Tree(\text{external positions of the component})\le\sum_{\text{constituent occurrences}}\ell(\text{assigned network}).
\]
\item The monomial of a resolved occurrence has an exact multi-index. Its fluctuation support is the set of whole fluctuation blocks on which its multi-index is positive. Background positions need not belong to this support, and no enlargement of the support to a connected geometric set is made.
\item If one constituent occurrence of a connected diagram component has a network that visits $\Omega^c_{L_J}$, then the union visits $\Omega^c_{L_J}$, and the visit lies in the same connected network as every external position of that component.
\item Fix an external root of a connected diagram component. For a vector factor $V$ attached to its parent at an input index $j$ use the counting bound $A_V(j)=(v_j/v_{o(V)})\abs V_b$, and if it is attached at its output use $A_V(o(V))=\abs V_b$; for a scalar factor use its counting-root bound. Then the sum over the resolution of the absolute coefficients, weighted by the exponentials of the network lengths, is bounded by the product of these bounds, multiplied at the end by the external normalization $v_{o}/v_{\mathrm{root}}$ when the component has an output $o$. For a substitution tree the volume ratios cancel exactly and the bound is the product of the norms $\abs\cdot_b$, in the row norm and in every volume-normalized column norm. For other diagrams the remaining ratios are kept in the scalar majorant. In particular, a counting scalar outer tensor on an old mesh, rooted at a new input $l$, carries the factor $v_l/v_{\mathrm{old}}$.
\item If a connected network of length $\ell$ contains the external positions $S$ and the unfixed index $i$, and visits $\Omega^c_{L_J}$, then $\Tree(S)\le\ell$ and $d_\Omega(i)\le\ell+L_0$, where $L_0$ is the fixed constant of the enlargement, and consequently
\[
 e^{a\Tree(S)+\alpha d_\Omega(i)}\le e^{\alpha L_0}e^{(a+\alpha)\ell}.
\]
\end{enumerate}
\end{lemma}

\begin{proof}
We first work with finite Taylor polynomials and finite diagrams. For a substitution $U(o;\mathbf l_1,\ldots,\mathbf l_k)=\sum_{j_1,\ldots,j_k}T(o;j_1,\ldots,j_k)\prod_rS_r(j_r;\mathbf l_r)$, the union of the network of $T$ with the networks of the $S_r$, joined at the indices $j_r$, connects every external position, so
\[
 \Tree(o,\mathbf l_1,\ldots,\mathbf l_k)\le\Tree(o,\mathbf j)+\sum_r\Tree(j_r,\mathbf l_r).
\]
We take absolute values before summing. For an output row we sum the inner factors by their rows and then the outer tensor by its row. If a fixed external input belongs to the factor $S_r$, we sum the other inner factors by their rows, use the input column of $S_r$, and use the $j_r$-column of $T$. Along this path the volume factors multiply to $(v_o/v_{j_r})(v_{j_r}/v_l)=v_o/v_l$, which is exactly the normalization of the column of the composite at $l$. This proves the cancellation for substitution trees. Products with a common output are treated in the same way: the networks meet at the output, and an input-rooted estimate uses the column of the factor carrying the input and the rows of the others.

For a general connected diagram, we replace every edge outside a spanning tree by two independently summed copies of its index, which only increases the nonnegative sums. We orient the tree from the factor containing the fixed external root and, starting at the leaves, sum each factor over all its indices except the one attached to its parent. A factor attached at an input index $j$ contributes at most its volume-normalized column bound times $v_j/v_{o(V)}$, since the column norm carries the factor $v_{o(V)}/v_j$; this is the counting bound $A_V(j)$. A factor attached at its output contributes its row bound. Finally summing the root factor over all its indices except the fixed root, we obtain the product of the counting bounds, with no count of the internal indices. The external normalization is applied only at the end. In a substitution tree the ratios on the path from the root input to the output multiply to $v_l/v_o$ and cancel against this normalization; in a general diagram they are kept. Connectedness is essential: factors with disjoint supports and no shared index have no tree, and their networks remain separate. A cycle without an external index, such as an undifferentiated trace, is not bounded in this way.

For compositions we use the ordered derivative formula
\[
 D^N(F\circ G)(0)[u_1,\ldots,u_N]=\sum_{n=0}^N\frac1{n!}\sum_{(B_1,\ldots,B_n)}F_n\bigl[G_{\abs{B_1}}[u_{B_1}],\ldots,G_{\abs{B_n}}[u_{B_n}]\bigr],
\]
where $G(0)=0$, the inner sum is over ordered partitions of $\{1,\ldots,N\}$ into nonempty blocks, and $F_n=D^nF(0)$, $G_k=D^kG(0)$. If $G(0)\ne0$, we write $G=G(0)+\widetilde G$ and apply this formula to $\widetilde G$, with $F_n$ replaced by $D^nF(G(0))$; the strict image condition justifies this translation and the infinite sums over the constant terms. Dividing by $N!$, we obtain the coefficient of $s^N$ in $\widehat F(\widehat G(s))$, where $\widehat F,\widehat G$ are the absolute majorants. Fixing an external index chooses its branch and then its position in that branch, as in $(\widehat F\circ\widehat G)'=(\widehat F'\circ\widehat G)\widehat G'$; no second Taylor factorial occurs. We define the resolution in part (1) by retaining a separate signed occurrence for every summand of this formula, every tuple of contracted indices, and every network refinement of the constituent tensors; grouping occurrences can only decrease the total variation. Its monomial is the product of the constituent monomials, which proves part (2). In a connected component the union of networks contains the visit of a constituent occurrence and is connected, which proves part (3). The diagram bound above, applied to every resolved occurrence, proves part (4) for finite truncations; the gap and Neumann conditions cover the Schur complements, Neumann series, and implicit series. Finally, by the image condition the nonnegative majorant sums converge on the common product. Tonelli's theorem then justifies every grouping, contraction, substitution, and derivative, and the resulting analytic functions agree because their coefficients agree.

For part (5), the network contains every external position, so $\Tree(S)\le\ell$. A route in the network from $i$ to the visited position $h\in \Omega^c_{L_J}$ has length at most $\ell$, and at most $L_0$ more length reaches $\Omega^c$. We conclude that $d_\Omega(i)\le\ell+L_0$, and adding the two inequalities with weights $a$ and $\alpha$ completes the proof.
\end{proof}

When we apply \cref{ch17:lem:coefficient-resolution}, we sum an internal index by keeping a positive part $b$ of the connecting exponent and using \eqref{ch17:eq:density-sum}. We convert a visit to $\Omega^c_{L_J}$ by \eqref{ch17:eq:wall-transport} into decay in the distance of every unfixed root to $\Omega^c$. In the applications below, the mixed-mesh contractions are substitutions, chains of matrices on a single lattice, and the composition with the incoming interaction, where the remaining ratio is the fixed factor $\Gamma_J\le M^{3J}$ of Step 4 of the proof of \cref{ch17:lem:decorated-white-primitives}; these factors are included in the common majorants.

\section{Conditional minimizers, constrained Hessians, and Haar densities}\label{ch17:sec:hessians}

We order the intermediate scales from the finest, $i=0$, to the coarsest, $i=s-1$. A stage of the conditional minimization is given by an initial segment $I$ of the conditional integrations in this order. At that stage the free coordinates of the scales in $I$ are minimized, with the coordinates $y=z_{I^c}$ of the remaining scales held fixed. The fluctuation increments in the following theorem are the coupling-scaled fluctuation inputs, $\zeta_i=q_i=g_iw_i$ in the notation of \cref{ch16:chapter}.

\begin{theorem}\label{ch17:thm:perturbed-hessian-coarea}
Fix an odd integer $M\ge3$, an integer $J\ge1$, $B=M^J$, the larger analytic domains and their strict margins, the label-density bound $v_J$, and the bounds for the classical and quantum torus norms. Use the conditional solution neighborhoods, the smaller neighborhoods of the higher parameters, and the product of sequential free coordinates of \cref{ch16:lem:admitted-products}, and the perturbed objective of \cref{ch16:thm:controlled-perturbation}. In joint coordinates this objective is
\begin{equation}\label{ch17:eq:objective}
 \mathcal F(v,u;\Rec)=F_0\bigl(v,z_0(v;\xi)+S_0(v;\xi)u;\xi\bigr)+\bar P\bigl(\Phi_0(v,u;\xi);\xi,\Rec\bigr),\qquad S_0=K_0^{-1/2}.
\end{equation}
Here $v=\Log c$ is the coarse background variable, $z_0$ is the unperturbed joint constrained minimizer, $K_0$ is its constrained Hessian, $u$ is the joint normalized free variable, and $\Phi_0$ is the unperturbed reconstruction followed by the gauge transformation $G^\Omega(c)$ of \cref{ch16:lem:strict-hybrid-gauge}, which depends only on $c$ and is the identity at all protected endpoints. The function $\bar P=P-S^{\mathrm{cl}}_\xi$ is the incoming perturbation of the classical action with the classical boundary-completion difference $S^{\mathrm{cl}}_\xi$ subtracted, including its classical constant, as in \cref{ch16:thm:controlled-perturbation}. Assume
\[
 \mathcal X^{\mathrm{cl}}(\bar P)\le2\sigma_*,\qquad0<\sigma_*\le\sigma_0,
\]
where the preliminary perturbation bound $\sigma_0>0$ satisfies the margins of \cref{ch16:lem:admitted-products} for the preconditioned conditional forces, the Hessians, and the images. Here $\mathcal X^{\mathrm{cl}}$ is the incoming classical norm on the incoming set $D_{\mathrm{in}}$ of unfixed labels, with the uniform profile $\rho_{\mathrm{cl}}$, as in \cref{ch16:thm:controlled-perturbation}. Assume the inequality on the conditional forces of \cref{ch16:lem:admitted-products}, and assume that the incoming fixed complement lies on the current fixed side, in the sense that for every current unfixed position $e$ and every incoming label $i$,
\begin{equation}\label{ch17:eq:incoming-wall-gate}
 d_{D_{\mathrm{in}}}(i)\ge B_sd_{\mathrm{current}}(e)-B_s\dist(e,i/B_s)-C_{\mathrm{buf},J}.
\end{equation}
Here $B_s=M^s$, $d_{D_{\mathrm{in}}}$ is measured in input units, $d_{\mathrm{current}}=d_\Omega$, $i/B_s$ is the position of $i$ in output units, and $C_{\mathrm{buf},J}$ is a fixed buffer that includes the chart stencils. Assume finally that both images of the incoming functions under the absolute reconstruction majorant, for the classical radius and for the quantum radius, are strictly smaller than half of that radius, as in \cref{ch16:lem:strict-two-radius-images}. Derivatives are taken at fixed real $(\xi,\Rec)$.

Then the following statements hold for every preceding depth $d_0\ge0$, every $1\le s\le J$, every output side $m\ge4$, every union $\Omega$ of output blocks, and every admissible real value of the frozen coordinates.
\begin{enumerate}
\item Every stage of the conditional minimization has a holomorphic conditional constrained minimizer, and every scale has a positive constrained Hessian $K_i^{\Omega,P}$, $0\le i<s$, which includes the exact factor $M^i$ of its scale. Writing the unnormalized fluctuation increments as $\zeta_i=g_iw_i$, the free logarithmic coordinates at scale $i$ are
\[
 z_i=z_i^*(c_{i+1};\Rec)+K_i(c_{i+1};\Rec)^{-1/2}\zeta_i,
\]
where $c_{i+1}$ is the reconstruction from the coarser scales, with their fluctuation increments inserted. There are $Q,a_1,\alpha_1,\lambda_J>0$ and $C_J,\mathcal L_{K,J}<\infty$ such that, on the outer profile $\InflatedProfile$ and for $\abs{\zeta_{i,e}}\le Q$, the real matrices satisfy
\begin{equation}\label{ch17:eq:hessian-bounds}
 \tfrac12\lambda_JI\le K_i^{\Omega,P}\le\mathcal L_{K,J}I,\qquad
 \norm{(K_i^{\Omega,P}+t)^{-1}}_{\mathrm{row,col},a_1}\le\frac{C_J}{\lambda_J+t}\quad(t\ge0).
\end{equation}
The complex continuation satisfies the same resolvent bound with a fixed larger $C_J$.
\item Adjoin to the regional problem, as identical independent auxiliary variables with unit Hessian, the free labels of the matched torus problem that are absent from it, so that both problems have the same sets of free labels. Let $K_i^{\mathbb T,\mathrm{cut}}$ be the restriction of the matched torus Hessian to the unfixed free labels, extended by the identity on the auxiliary labels and by zero on the mixed blocks. If
\[
 \Delta_i(v,\zeta;\Rec)=K_i^{\Omega,P}-K_i^{\mathbb T,\mathrm{cut}}=\sum_{\nu,\kappa}\Delta_{i;\nu\kappa}(\Rec)v^\nu\zeta^\kappa,
\]
then
\begin{equation}\label{ch17:eq:hessian-defect}
 \sup_{u\ \mathrm{unfixed}}e^{\alpha_1d_\Omega(u)}\sum_{v',\nu,\kappa}\abs{\Delta_{i;uv',\nu\kappa}(\Rec)}\,\InflatedProfile^\nu Q^{\abs\kappa}e^{a_1\Tree(u,v',\supp\nu,\supp\kappa)}\le C_J.
\end{equation}
The transposed bound holds, and so do the bounds rooted at any unfixed background or free external index, with the volume-normalized column factors of \cref{ch01:def:analytic-norms}. The underlying mixed tensors with $k$ additional background or free derivatives have bounds $C_J^{k+2}(k+2)!$ at every external index, on a fixed larger product. The auxiliary blocks have zero defect. If $\Omega^c$ is empty, the defect has zero nonconstant part.
\item The same rooted comparisons with boundary decay hold for the conditional minimizers of every stage, the inverse square roots, the logarithms of the Haar disintegration densities, and the diagonal densities $(\log K_i)_{xx}/2$, after composition with the reconstruction from every value of the coarser fluctuation variables. In particular, let $L_i$ be the negative logarithm of the Haar disintegration density at scale $i$ after its field-independent normalization has been moved into $\Lambda$, and put
\[
 L_{0,i}(c_{i+1})=L_i\bigl(c_{i+1},z_i^*(c_{i+1})\bigr),\qquad B_i=L_{0,i}+\tfrac12\log\det K_i.
\]
The fluctuation-dependent part of the density contains
\begin{equation}\label{ch17:eq:measure-fluctuation}
 \sum_{i=0}^{s-1}\bigl[L_i(c_{i+1},z_i)-L_{0,i}(c_{i+1})+B_i(c_{i+1})-B_i(\widehat c_{i+1})\bigr],
\end{equation}
where $\widehat c_{i+1}$ is the reconstruction with every fluctuation increment equal to zero. These differences start at degree one in $\zeta$, whereas the classical remainders start at degree three. The comparison includes the correction $\log(T_{II})-(\log T)_{II}$ that arises when the logarithm of a torus matrix $T$ is restricted to the list $I$ of unfixed labels. A torus diagonal density that is absent from the regional problem and has an unfixed external index has a network visiting $\Omega^c$ and is summable from that index; its constant part remains an explicit sum over omitted labels.
\item All constants have the dependence described in item (3) of \cref{ch01:conv:constants}, including independence of the later widths. They are evaluated on the fixed larger products before the final radius for the free coordinates is decreased. In particular they do not change at the renormalization group steps following the creation of a new component of $\Omega^c$.
\end{enumerate}
\end{theorem}

\begin{proof}
Steps 1 and 2 verify that the incoming perturbation satisfies the hypotheses of \cref{ch16:thm:controlled-perturbation} on all parameter domains. Steps 3 to 5 construct the conditional minimizers by a contraction argument that is uniform in the higher parameters, and check that every sequential substitution stays in the admissible domains. Steps 6 to 9 compare the regional and torus quantities, first without and then with the perturbation. Step 10 treats the matrix functions and the Haar densities, Step 11 the triangular determinant, and Step 12 the dependence of the constants on the frozen values.

We fix the real frozen value $b=(\xi,\Rec)$ and write $\mathcal Q_\theta=F_0+\theta V$, where $V=\bar P\circ\Phi_0$ and $0\le\theta\le1$; here $\theta$ is an interpolation parameter. Returning from $u$ to $z=z_0+S_0u$ is an invertible affine change of variables at fixed $v$. In the free logarithmic coordinates, the objective is
\[
 \mathcal Q_\theta(v,z)=F_0(v,z;\xi)+\theta\bar P\bigl(\Log[(\exp\Psi_R(v,z;\xi))^{G^\Omega(v)}];\xi,\Rec\bigr),
\]
where $\Psi_R$ is the reconstruction before the gauge transformation. We use only this change of variables; we do not need the incoming perturbation to take the same value at the reconstruction before and after the gauge transformation.

\paragraph{Step 1: the incoming hypotheses.}
The incoming analytic labels form the set $D_{\mathrm{in}}$, with their positions on the lattices where they were introduced and their volume factors. On this set the incoming profile is uniformly $\rho_{\mathrm{cl}}$; the coordinates $\xi$ are fixed parameters, and the spatially varying profiles are output profiles. The strict bounds on the image under the reconstruction majorant place $\Phi_0$, on every admissible stage, inside this uniform input ball with a margin of half the radius. The same holds for every segment between two compared images, since coefficientwise a convex combination is bounded by the larger of the two absolute image bounds. In particular, the input jet bounds of \cref{ch16:thm:controlled-perturbation} apply before any conditional minimization, on all its parameter neighborhoods.

\paragraph{Step 2: the incoming boundary comparison.}
For example, the inclusion of fixed complements used in \cref{ch16:thm:controlled-perturbation} implies \eqref{ch17:eq:incoming-wall-gate} as follows. Let $h$ be an incoming fixed position nearest to $i$. A path from $e$ to the current fixed complement can first go to $i/B_s$, then follow the scaled image of a shortest path from $i$ to $h$, of length $d_{D_{\mathrm{in}}}(i)/B_s$, and finally cross the fixed buffer between the incoming and the current fixed complements, whose length is included in $C_{\mathrm{buf},J}/B_s$. Multiplying by $B_s$, we obtain \eqref{ch17:eq:incoming-wall-gate}. This argument uses no regularity or component assumption on the boundary. We keep \eqref{ch17:eq:incoming-wall-gate} as a hypothesis on the geometry; choosing small radii does not establish it.

In particular, if $\alpha_{\mathrm{out}}\le B_s\alpha_{\mathrm{cl}}$ and a chart branch has spatial rate $b>\alpha_{\mathrm{out}}$, then
\begin{align*}
 &e^{\alpha_{\mathrm{out}}d_{\mathrm{current}}(e)}e^{-\alpha_{\mathrm{cl}}d_{D_{\mathrm{in}}}(i)}e^{-b\dist(e,i/B_s)}\\
 &\quad\le e^{\alpha_{\mathrm{out}}C_{\mathrm{buf},J}/B_s}e^{-(\alpha_{\mathrm{cl}}-\alpha_{\mathrm{out}}/B_s)d_{D_{\mathrm{in}}}(i)}e^{-(b-\alpha_{\mathrm{out}})\dist(e,i/B_s)}.
\end{align*}
This is the transport of the input boundary weight used in the perturbation bound. Input network lengths are divided by $B_s$ when expressed in output units. Column normalizations remain the ratios of the cubed mesh sizes of their output and input labels. Local rates may be chosen below the finite minimum of these bounds over $1\le s\le J$; the rates returned after a full step use $s=J$ and $B_s=B$. If there is no incoming boundary, the incoming nonconstant boundary function is zero.

\paragraph{Step 3: conditional minimizers.}
For each stage $I$ (an initial segment of the conditional integrations, with $y=z_{I^c}$ fixed), let $a_I(v,y)$ be the unperturbed conditional minimizer. By \cref{ch16:lem:admitted-products}, there exist a fixed solution neighborhood, a smaller neighborhood of the higher parameters, and an invertible matrix $A_I=F_{0,II}(v,a_I,y)$. Let $I_0$ be a common bound for the coefficient norms of the inverses $A_I^{-1}$ on the fixed larger neighborhoods. Using the uniform incoming jet bound with $j=1$ and the fixed rows of the chart derivatives, we obtain, at every admissible higher parameter $y$ and not only at the joint minimizer,
\begin{equation}\label{ch17:eq:conditional-force-bound}
 \norm{V_I(v,a_I,y;b)}_{\mathrm{coeff}}\le C_P\sigma_*\rho_{\mathrm{cl}},\qquad\delta_V\le I_0C_P\sigma_*\rho_{\mathrm{cl}},
\end{equation}
where $\delta_V$ is the bound on the preconditioned conditional forces of \cref{ch16:lem:admitted-products}. The constant $C_P$ is evaluated on the larger products at the preliminary bound $\sigma_0$ and includes the conversion from $z$ to $u$. The jet bound with $j=2$ and the chain rule, including the term $\bar P'\Phi_{0,zz}$, give a uniform bound $C_P\sigma_*$ on the mixed second derivatives. The real gap, the complex coefficient margins, the inequality $\delta_V\le r_I/4$, where $r_I$ is the radius of the solution neighborhood, and the derivative inequality below are then the hypotheses of \cref{ch16:lem:admitted-products}. On that neighborhood the map
\[
 h\longmapsto h-A_I^{-1}\mathcal Q_{\theta,I}(v,a_I+h,y)
\]
has derivative norm at most $1/2$ and value at zero of norm at most $r_I/4$. Hence it maps the closed ball of radius $r_I/2$ into itself, and its successive iterates converge: the difference of the iterates $n+1$ and $n$ is at most $2^{-n}$ times the first difference. The same bounds on absolute coefficients hold for the analytic iterates, which gives normal convergence on every smaller product and a holomorphic limit. Uniqueness follows by subtracting two fixed point equations. Since the objective is holomorphic and real at real arguments, the minimizer is real at real data. The real Hessian and complex resolvent margins are those of \cref{ch16:thm:controlled-perturbation}; its weighted resolvent bounds rest on the exponential conjugation method of Combes and Thomas \cite{CombesThomas1973}. Taking the finite minimum over $s\le J$ and over the stages fixes $\lambda_J$ and $Q$, and taking the finite maximum fixes $\mathcal L_{K,J}$ and the bound on the inverses. The factors $M^i$ belong to these matrices before the extrema are taken.

\paragraph{Step 4: sequential substitution.}
For the sequential construction we use the inequality
\begin{equation}\label{ch17:eq:sequential-gate}
 C_{\mathrm{seq},J}(2\delta_V+Q)<\tfrac12\varepsilon_{\mathrm{tube}},
\end{equation}
where $\varepsilon_{\mathrm{tube}}$ is the minimal margin of the parameter and solution neighborhoods of \cref{ch16:lem:admitted-products}. By the contraction bound, the displacement from $a_I(v,y)$ is at most $2\theta\delta_V$, for every admissible $y$. The uniform mixed coefficient bounds and the bound on the free inverses give a finite common Lipschitz constant $D_J$ of the conditional minimizers in $y$, and $\norm{K_i^{-1/2}}\le C_{S,J}$. Let $A_i$ be a bound for the deviations of the already reconstructed coarser coordinates from the unperturbed joint chain. The next insertion satisfies
\begin{equation}\label{ch17:eq:sequential-displacement}
 A_{i-1}\le(1+D_J)A_i+2\delta_V+C_{S,J}Q,\qquad A_s=0.
\end{equation}
Iterating, we find that the total deviation is at most $\max\{1,C_{S,J}\}\sum_{k=0}^{J-1}(1+D_J)^k(2\delta_V+Q)$. We choose $C_{\mathrm{seq},J}$ larger than the coefficient of this finite expression and than every fixed embedding constant of the stages. Then \eqref{ch17:eq:sequential-gate} keeps every substituted coarser argument in the admissible domains, even after the deviation is doubled. The mixed bounds and the constant $C_{\mathrm{seq},J}$ are fixed before the later widths. In particular the preliminary bound $\sigma_0$ must satisfy the same inequality with $\delta_V\le I_0C_P\sigma_0\rho_{\mathrm{cl}}$.

\paragraph{Step 5: dependence on the coarser coordinates.}
We explain why this construction controls coefficients when the coarser coordinates move. Differentiating the stationary equation at fixed $y$, we obtain
\begin{equation}\label{ch17:eq:prefix-response}
 \partial_\theta z_I=-\mathcal Q_{\theta,II}^{-1}V_I.
\end{equation}
If $y$ moves, the right side has the additional term $-\mathcal Q_{\theta,II}^{-1}\mathcal Q_{\theta,II^c}\partial_\theta y$. A sequential insertion also has the term $(\partial_\theta K_i^{-1/2})\zeta_i$. All these terms are included in \cref{ch16:thm:controlled-perturbation}, and \eqref{ch17:eq:sequential-displacement} allows their evaluation at every value of the coarser coordinates. At zero free input, conditional uniqueness identifies the sequential chain with the joint minimizer. Decreasing $Q$ therefore does not change the admissible profile of the coarse background.

\paragraph{Step 6: the unperturbed boundary comparison.}
We use matched reconstruction formulas on the unfixed labels. Away from the fixed construction enlargement of $\Omega^c$, the formulas defining the common unfixed outputs, the block axial gauge equations, and the choices of eliminated edges coincide. Their unique admissible solutions on the unfixed outputs therefore coincide there. This does not assert that the extended fields agree inside a component of $\Omega^c$. The prescribed values in $\Omega^c$ are separate bounded sources; a source deep inside a component of $\Omega^c$ need not be small at its own position. For the torus action composed with two reconstructions $\Psi_R=\Psi_{\mathbb T}+d$, the exact difference is
\begin{equation}\label{ch17:eq:action-word-difference}
 F_R-F_{\mathbb T}=\int_0^1DE_{d_0}(\Psi_{\mathbb T}+\tau d)[d]\,d\tau
\end{equation}
on the common unfixed variables. If auxiliary missing variables are kept, their separate auxiliary quadratic form is kept as a term rooted in $\Omega^c$. Every differentiated coefficient of the integral contains either a derivative of a changed formula or the factor $d$ attached to an action tensor. The action and chart tensors on the larger domain are the tensors of \cref{ch16:thm:controlled-perturbation}, and their contractions with a changed formula are networks covered by \cref{ch15:lem:wall-networks}. Rooting at any unfixed external index connects that index to the change. For example, the second derivative is
\[
 D_z^2(E_{d_0}\circ\Psi)=\Psi_z^{\mathsf T}E_{d_0}''(\Psi)\Psi_z+E_{d_0}'(\Psi)\Psi_{zz}.
\]
Both terms are present: stationarity in $z$ does not make $E_{d_0}'(\Psi)$ vanish. Differentiating \eqref{ch17:eq:action-word-difference} produces derivatives falling on $E_{d_0}$ as well as derivatives falling on $d$; the former remain connected to the inserted change. Summing a bounded source in $\Omega^c$ from an unfixed root is justified by \eqref{ch17:eq:density-sum} and \eqref{ch17:eq:wall-transport}. This proves factorial bounds with a visit to $\Omega^c$ for the difference before minimization, rooted at every external index.

\paragraph{Step 7: minimization, Schur complements, and restriction.}
We pass to the conditional minimizers through their stationary equations. A force carrying a visit to $\Omega^c$ is multiplied by the free inverse, as in \eqref{ch17:eq:prefix-response}. The coefficientwise majorant series of the implicit equation converges on the admissible products: if $Y$ denotes the majorant of the solution without the visit, the variation carrying the visit has the form $H(x,Y)/(1-A(x,Y))$, where $A\le3/4$ after the auxiliary radius has been fixed. This follows by substituting the majorant of the variation into the differentiated fixed point equation and summing $\sum_{n\ge0}A^n$. The visit and all fixed external indices are kept. For a partial minimum the exact second derivative is the Schur complement
\begin{equation}\label{ch17:eq:schur-hessian}
 \mathcal Q_{yy}-\mathcal Q_{yI}\mathcal Q_{II}^{-1}\mathcal Q_{Iy}.
\end{equation}
We subtract two such expressions by replacing the three factors of the product one at a time. An inverse difference is $A^{-1}-B^{-1}=-A^{-1}(A-B)B^{-1}$, so each summand keeps a factor with a visit to $\Omega^c$. Only constrained free blocks are inverted. Positivity passes to the Schur complement by minimizing the positive quadratic form over the coordinates in $I$.

A separate correction compares restriction with elimination. Let $A$ be a single torus free matrix, and let $R$ and $H$ be the lists of unfixed and of missing labels (in this paragraph $R$ and $H$ denote lists of labels). Block elimination gives
\begin{equation}\label{ch17:eq:restriction-elimination}
 (A^{-1})_{RR}-A_{RR}^{-1}=A_{RR}^{-1}A_{RH}\bigl(A_{HH}-A_{HR}A_{RR}^{-1}A_{RH}\bigr)^{-1}A_{HR}A_{RR}^{-1}.
\end{equation}
To check this identity, we solve $Ax=(f_{R},0)$, first for $x_{R}$ in terms of $x_{H}$ and then for $x_{H}$. Each term on the right of \eqref{ch17:eq:restriction-elimination} passes through a missing index. Its coefficient network from an unfixed root therefore reaches $\Omega^c$, and the part of the exponential decay not used for the network weight sums over that index. Adding this correction to the regional comparison, we obtain the unperturbed defect at each stage. This operation compares restrictions of one torus matrix; it does not identify a regional problem with a torus problem whose complement is set to zero.

\paragraph{Step 8: from factorial bounds to coefficient bounds.}
The bounds just obtained for the tensors of order $k+2$, at every external index, are $C_J^{k+2}(k+2)!$. Dividing by the Taylor factorials and using \eqref{ch17:eq:factorial-series}, we find that on a fixed smaller auxiliary radius $h$ the summed matrix coefficients are at most
\[
 C_J^2\sum_{k\ge0}(k+1)(k+2)(C_Jh)^k=2C_J^2(1-C_Jh)^{-3}.
\]
We choose the outer background and free products inside this radius and do not change them later. Since any index may be fixed before the sums are taken, we may fix a background or free input, which proves the column versions.

\paragraph{Step 9: adding the perturbation.}
We insert the unperturbed regional matrix at its own chain of coarser arguments. The defect in \eqref{ch17:eq:hessian-defect} is the sum of three terms: the change of the objective and of the minimizer at a common coarser argument, the change of that argument, and the unperturbed regional--torus defect. The first term has the rooted perturbation bound of \cref{ch16:thm:controlled-perturbation}, with its visit to $\Omega^c$. The second term is the integral of an unperturbed matrix derivative against the difference of the coarser reconstructions, which carries a visit to $\Omega^c$; both are defined on the same products by the statement on the sequential product in \cref{ch16:lem:admitted-products}, and the network joins the new index to the visit. The third term was treated in Steps 6 to 8. The sum satisfies \eqref{ch17:eq:hessian-defect} for fixed smaller $a_1,\alpha_1>0$. The real matrices are symmetric, which gives the transposed bound with the same constants. Rooting each contraction at the specified background or free index gives the other bounds, with its volume factor. The auxiliary blocks are identical throughout and have zero difference; no matrix element rooted purely at an auxiliary label is assigned a boundary decay.

\paragraph{Step 10: matrix functions and Haar densities.}
For positive matrices, integrating the scalar resolvent and applying the resolvent identity, we obtain the ordered identities of holomorphic functional calculus (see \cite{Bhatia1997})
\begin{align}
 A^{-1/2}-B^{-1/2}&=-\frac1\pi\int_0^\infty t^{-1/2}(A+t)^{-1}(A-B)(B+t)^{-1}\,dt,\label{ch17:eq:square-root}\\
 \log A-\log B&=\int_0^\infty(A+t)^{-1}(A-B)(B+t)^{-1}\,dt.\label{ch17:eq:log-difference}
\end{align}
Indeed $A^{-1/2}=\pi^{-1}\int_0^\infty t^{-1/2}(A+t)^{-1}\,dt$ and $\log A=\int_0^\infty[(1+t)^{-1}I-(A+t)^{-1}]\,dt$ hold for positive numbers, hence for real symmetric positive matrices by orthogonal diagonalization. The common complex resolvent margins extend these identities holomorphically. The coefficient majorants of the integrands are $C_Jt^{-1/2}(\lambda_J+t)^{-2}$ and $C_J(\lambda_J+t)^{-2}$, which are integrable at both endpoints. We may therefore integrate all contractions and derivatives termwise. Each difference keeps its factor $A-B$.

The logarithm of the Haar disintegration density is a sum of logarithms of the nonvanishing Haar factors of finitely many local formulas and of the determinants of the constraints at the eliminated edges, as in \cref{ch01:def:block-map}. On the fixed chart we write each factor as $a(1+h)$, with its nonzero value $a$ at the base point and a fixed strict coefficient margin $\abs h<1$; the series $\log(1+h)=\sum_{k\ge1}(-1)^{k+1}h^k/k$ gives its absolute derivatives on every smaller chart. The base factors are moved into $\Lambda$. Away from the boundary the local formulas agree. In a difference of compositions, a changed local formula, or a changed minimizer, inverse square root, or coarser argument, is a factor carrying a visit to $\Omega^c$. The preceding calculus therefore yields the comparison of the Haar densities at all orders.

A diagonal matrix-logarithm density keeps its index until an external unfixed index is chosen. Its nonconstant derivatives are chains of resolvents with a differentiated matrix inserted. Rooting at that insertion makes the trace sum summable by \eqref{ch17:eq:density-sum}. If a torus diagonal density is missing from the regional problem, its network with any unfixed external index must reach the missing index, and the same argument sums over the missing index from the unfixed index. A pure constant has no such index and remains its exact sum over omitted labels.

To compare with the restriction of the torus logarithm, we set $T=K_i^{\mathbb T}$ and $C=T_{II}\oplus T_{HH}$, where $I=I_i$ and $H=H_i$ are the lists of unfixed and missing free labels at scale $i$. Using the integral representation of the logarithm, we obtain
\begin{equation}\label{ch17:eq:restricted-log}
 \log T_{II}-(\log T)_{II}=\int_0^\infty\bigl[(T+t)^{-1}-(C+t)^{-1}\bigr]_{II}\,dt.
\end{equation}
The integrand equals $-[(T+t)^{-1}(T-C)(C+t)^{-1}]_{II}$. Its middle factor connects unfixed and missing labels, so it has a visit to $\Omega^c$, and it has the same integrable bound as \eqref{ch17:eq:log-difference}. Thus this correction, composed with every value of the coarser fluctuation variables, is included in the comparison.

\paragraph{Step 11: the triangular determinant.}
The determinant of the triangular change of variables of \cref{ch09:thm:sequential-density}, or of the identical finite triangular calculation for the present minimizers, is
\[
 \prod_ig_i^{3f_i}\det K_i(c_{i+1})^{-1/2},
\]
where $f_i$ is the number of free blocks at scale $i$. Adding and subtracting the logarithm of its value at zero fluctuation, we combine the Haar terms into \eqref{ch17:eq:measure-fluctuation}. Only one determinant occurs. For an analytic function $G$ on the common product,
\begin{equation}\label{ch17:eq:free-difference}
 G(v,\zeta)-G(v,0)=\int_0^1D_\zeta G(v,\tau\zeta)[\zeta]\,d\tau.
\end{equation}
This keeps every positive degree in $\zeta$. Stationarity and the exact normalization by $K_i^{-1/2}$ cancel the terms of degrees one and two of each classical remainder, so its degree is at least three. Composition with the coarser scales does not lower that degree.

\paragraph{Step 12: constants and measurability.}
The proof uses only margins on fixed domains, finitely many operations on the stages, \eqref{ch17:eq:density-sum}, normally convergent coefficientwise majorant series, and the resolvent integrals above. These give the dependence stated in part (4). A new component of $\Omega^c$ changes the set of indices carrying a visit and the frozen values, but no constant. If $\Omega^c$ is empty, no formula changes and there is no input carrying a visit. This proves the stated special cases and uniformities.

We also record the regularity in the frozen values needed for finite iterations. If the incoming coefficients and the charts depend on the real value $b$ in a Borel way, then so does every iterate of the contraction in Step 3: finite sums, products, and admissible inverses preserve this property, and a convergent coefficient composition is a limit of its polynomial truncations. The geometric contraction bound is uniform in $b$, so the coefficients of the limit are Borel. The resolvent integrals for square roots and logarithms are limits of finite Riemann sums with uniformly bounded tails, and preserve this property. With continuous incoming data on a fixed branch, the same uniform convergence gives continuity. None of these arguments differentiates in $b$. In the finite analytic recursion, the initial torus functions and the local formulas in the identity chart have this regularity, and substitution of a frozen value and the integrals over fixed products below preserve Borel dependence. Thus every coefficient family constructed by the recursion depends on the frozen values in a Borel way.
\end{proof}

\section{The conditional fluctuation potential}\label{ch17:sec:potential}

\begin{lemma}\label{ch17:lem:complete-white-potential}
Use the data and products of \cref{ch17:thm:perturbed-hessian-coarea}. Let the incoming effective interaction be
\[
 R=r_{d_0}(\,\cdot\oplus0)+f,\qquad Dr_{d_0}(0)=0,\quad\Anorm{\rho}{a_{\mathbb T}}{r_{d_0}}\le E,\quad\mathcal X^\rho(f)\le\sigma_*\le\sigma_0,
\]
where $r_{d_0}$ is the torus effective interaction at depth $d_0$, with its torus symmetries, evaluated with the fixed slots set to zero, and $f$ is its boundary part. The value of $r_{d_0}$ at zero is unrestricted. For $0\le i<s\le J$ set
\[
 t_{\max}=\max_ig_i=\sqrt{\gamma/(Mm)},\qquad g_i=t_{\max}\lambda_i,\qquad\lambda_i=M^{(i-s+1)/2}\in[\lambda_*,1],
\]
where $\lambda_*=M^{-(J-1)/2}$. Let $F_i$ be the partially minimized objective at scale $i$, without the coupling factor and with its exact normalization $M^i$, and define
\[
 T_i(c_{i+1},\zeta_i)=F_i\bigl(c_{i+1},z_i^*+K_i^{-1/2}\zeta_i\bigr)-F_i(c_{i+1},z_i^*)-\tfrac12\abs{\zeta_i}^2.
\]
For a fixed admissible real coarse background and fixed real frozen values, the conditional fluctuation potential, with its values at zero fluctuation moved into the prefactor, is exactly
\begin{equation}\label{ch17:eq:potential}
\begin{split}
 V(w)={}&\tfrac12\sum_i\abs{w_i}^2+\sum_ig_i^{-2}T_i(c_{i+1},g_iw_i)\\
 &+\sum_i\bigl[L_i(c_{i+1},z_i)-L_{0,i}(c_{i+1})+B_i(c_{i+1})-B_i(\widehat c_{i+1})\bigr]\\
 &+R\bigl(\Log C(v,(g_iw_i)_i;\Rec)\bigr)-R\bigl(\Log C(v,0;\Rec)\bigr),
\end{split}
\end{equation}
where $C(v,\zeta;\Rec)$ denotes the reconstructed configuration on the unfixed labels. Only logarithms of unfixed labels in the identity chart are taken. The Haar disintegration density, the degree-one terms of the incoming interaction, and all determinant terms of the coarser scales occur in \eqref{ch17:eq:potential}. The separate prefactor contains all constants of the block averaging constraint and of the Haar measure, the Gaussian powers of the couplings, the values at the minimizers, and the external compact factors.

There are $A\ge1$, $a'>0$, and $C_{V,J},C_0<\infty$ with $C_0\ge2A^3$, and a fixed radius $s_W>0$ for the increments $\zeta_i$, such that on the product $\abs{g_iw_{i,e}}\le s_W$
\begin{equation}\label{ch17:eq:potential-series}
 V(w)=\tfrac12\abs w^2+\sum_{k\ge3}t_{\max}^{k-2}P_k(w)+\sum_{k\ge1}t_{\max}^kQ_k(w),\qquad\norm{P_k}_{a',1}\le A^k,\quad\norm{Q_k}_{a',1}\le\mathcal L_VA^k,
\end{equation}
where $P_k,Q_k$ are homogeneous polynomials of degree $k$ in the scalar coordinates of $w$ and $\mathcal L_V=C_{V,J}(1+E+\sigma_0)\ge1$. For $T_k(w)=\sum_{\abs\eta=k}t_\eta w^\eta$ the norm is
\[
 \norm{T_k}_{a',1}=\sup_j\sum_{\eta_j>0}\eta_j\abs{t_\eta}e^{a'\Tree(\supp\eta)}.
\]
The bounds \eqref{ch17:eq:potential-series} also hold before the background is evaluated, in the coefficient algebra at the outer profile, with the background positions included in the network length.

The following inequalities suffice for a common enclosing product. All constants below are computed at the preliminary bound $\sigma_0$, before $\sigma_*$, $\rrecord$, $r$, and the later widths are fixed. Let $r_{q,U},r_{q,P},u_0$ be the radii of the unperturbed and perturbed sequential free products and of the joint product of \cref{ch16:lem:admitted-products}, and $c_{\mathrm{tube},J}$ its fixed conversion constant. For $\rho_\bullet\in\{\rho,\rho_{\mathrm{cl}}\}$ let $\Rraw$, $C_{\mathrm{lift},J}$, $s_{\mathrm{aux},\bullet}$, $C_{T,\bullet}$, $C_{\Delta,\bullet}$, and $A_{0,\bullet}$ be the radius of the larger construction domain, the bound on the lifts, the auxiliary radii, the constants of the image increments, and the bound on the image at zero fluctuation in \cref{ch16:lem:strict-two-radius-images}. Choose $\Rwhite>0$ strictly below
\[
\begin{split}
 \min\bigl\{&s_W/2,\ r_{q,U}/2,\ r_{q,P}/2,\ c_{\mathrm{tube},J}u_0,\ \Rraw/(64C_{\mathrm{lift},J}),\ 1/(4A),\ 1/(16C_0),\\
 &\min_{\rho_\bullet\in\{\rho,\rho_{\mathrm{cl}}\}}\{s_{\mathrm{aux},\bullet}/2,\ 1/(4C_{T,\bullet}),\ \rho_\bullet/[8(2C_{T,\bullet}+C_{\Delta,\bullet})]\}\bigr\}.
\end{split}
\]
Impose the margins of the larger construction domain and of the stages in \cref{ch16:lem:strict-two-radius-images} and, for both radii, $A_{0,\bullet}\le\rho_\bullet/4$ and $C_{T,\bullet}K_*\rho_\bullet\le1/4$. Let $C_{\mathrm{en},J}\ge1$ be a fixed constant (in the application it is the constant of the enclosure estimate \cref{ch18:thm:all-point-enclosure}; nothing else from that section is used), and require
\begin{equation}\label{ch17:eq:enclosure-budgets}
 r,\ \rrecord,\ \sigma_*\rho_{\mathrm{cl}}\le\frac{\lambda_*\Rwhite}{6C_{\mathrm{en},J}},\qquad
 \delta_{\mathrm{en}}=C_{\mathrm{en},J}(r+\rrecord+\sigma_*\rho_{\mathrm{cl}}),\qquad
 \Ren=2\delta_{\mathrm{en}}/\lambda_*,
\end{equation}
where $r$ is the radius of the small-link domain $\SmallChart r$ and $\rrecord$ is the bound on the frozen values. Then the whole product
\[
 \mathcal K_{\mathrm{en}}=\prod_{i,e}\{\abs{w_{i,e}}\le\Ren/t_{\max}\}
\]
is admissible. On it the absolute coefficient increment of either incoming image is bounded by
\[
 \frac{C_{T,\bullet}\Ren}{1-C_{T,\bullet}(K_*\rho_\bullet+\Ren)}+C_{\Delta,\bullet}\Ren\le\rho_\bullet/8,
\]
the denominator is at least $1/2$, and the image under the absolute reconstruction majorant at spatial exponent zero is at most $3\rho_\bullet/8$. If these bounds are imposed in a weighted coefficient row, the same bound holds in that row. In the alternative provided by \cref{ch16:lem:strict-two-radius-images}, the unperturbed image row $S_c^0\le C_J\rho_\bullet$ at an exponent $c>0$, larger than the working exponent, gives one fixed positive exponent $\vartheta c$ with
\[
 S_{\vartheta c}^0\le(3\rho_\bullet/8)^{1-\vartheta}(C_J\rho_\bullet)^\vartheta<(\sqrt3/4)\rho_\bullet.
\]
For the unperturbed construction, this exponent is fixed before the regional rates. For the perturbed reconstruction, impose the output-rate conditions of \cref{ch16:thm:controlled-perturbation} at the allocated working rates below this exponent and bound its coefficient difference from the unperturbed reconstruction there. Fix the larger-product difference-row constant $C_{P,\bullet}$ before the final perturbation bound and require $C_{P,\bullet}\sigma_*\rho_{\mathrm{cl}}\le\rho_\bullet/64$ for both radii. The total image row is then less than $(\sqrt3/4+1/64)\rho_\bullet<\rho_\bullet/2$ at those rates; no bound for the perturbed image at the original higher exponent is required. Thus the margin of half the incoming radius is an absolute weighted coefficient bound, not only a bound on values.

Let $C_1$ be the constant of \cref{ch17:lem:ambient-white-comparison}, which depends only on $A$, $a'$, and $v_J$. For $0<t_{\max}\le t_*<e^{-1}$ impose separately
\begin{equation}\label{ch17:eq:coupling-budgets}
 C_0t_*^2\mathcal L_V\le1/16,\qquad C_1t_*(\mathcal L_V+2)<\frac1{4\sqrt3}.
\end{equation}
Then on a neighborhood of the whole real product $\mathcal K_{\mathrm{en}}$,
\begin{equation}\label{ch17:eq:curvature}
 \norm{D_w^2(V-\abs w^2/2)}_{\mathrm{row,col},a'}\le\frac{C_0(\Ren+t_{\max}^2\mathcal L_V)}{(1-A\Ren)^3}\le\frac8{27}<\frac12.
\end{equation}
If $p(t)=A_{\mathrm w}\log(1/t)$, the products of balls with radii $p_i\in[p(t_{\max}),p(t_{\max})+A_{\mathrm w}(J-1)(\log M)/2]$ lie in $\mathcal K_{\mathrm{en}}$ whenever
\[
 t_{\max}\bigl[p(t_{\max})+A_{\mathrm w}(J-1)(\log M)/2\bigr]<\Ren.
\]
The cutoff of the small-field event \eqref{ch01:eq:literal-event} is the case $p_i=p(t_{\max})$ for all $i$; the cutoffs $p_i=p(g_i)$ at the individual scales are a separate case. All constants have the dependence described in item (3) of \cref{ch01:conv:constants}, and are also independent of the later widths, of the frozen values, and of the individual radii of the fluctuation balls.
\end{lemma}

\begin{proof}
We keep $b=(\xi,\Rec)$ fixed and write $t=t_{\max}$. The data of \cref{ch17:thm:perturbed-hessian-coarea} include the uniform incoming classical profile, \eqref{ch17:eq:incoming-wall-gate}, and \eqref{ch17:eq:conditional-force-bound} on all parameter neighborhoods. We choose the radius for the increments $\zeta$ below $Q$; then \eqref{ch17:eq:sequential-gate} holds at every stage. These are the hypotheses for the perturbed tensors used here, and no derivative in the frozen values is needed.

\paragraph{Step 1: the exact formula.}
At scale $i$, the minimum at fixed coarser arguments is $F_i(c_{i+1},z_i^*)$. With $S_i=K_i^{-1/2}$, Taylor's formula with integral remainder gives
\begin{equation}\label{ch17:eq:cubic-remainder}
 g_i^{-2}T_i(c_{i+1},g_iw_i)=\frac{g_i}2\int_0^1(1-\tau)^2D_z^3F_i(c_{i+1},z_i^*+\tau g_iS_iw_i)[S_iw_i,S_iw_i,S_iw_i]\,d\tau.
\end{equation}
The terms of orders zero, one, and two have been, respectively, subtracted, annihilated by stationarity, and normalized to $\abs{g_iw_i}^2/2$ by $S_i^{\mathsf T}K_iS_i=I$. The successive partial minima satisfy $F_{i+1}=M\min F_i$, and the couplings satisfy $g_{i+1}^2=Mg_i^2$, so the values at the minimizers telescope in the exponent. In decreasing order of the scales, the map from $w$ to $z$ is triangular with diagonal blocks $g_iS_i$. This proves the classical part of \eqref{ch17:eq:potential} and its Gaussian term, by the same finite algebra as in \cref{ch09:thm:sequential-density}. The Haar disintegration density and the single triangular determinant split as in part (3) of \cref{ch17:thm:perturbed-hessian-coarea}. Splitting $R(\Log C)$ at its value at zero fluctuation, we obtain the last two terms. All values removed at zero fluctuation, the Gaussian powers, and the constant Haar densities belong to the prefactor; none of them is set to one.

\paragraph{Step 2: the two series.}
We fix an auxiliary radius $h_0>0$ for the increments $\zeta$, strictly inside every admissible stage product and its joint neighborhood. \Cref{ch17:thm:perturbed-hessian-coarea} gives tensors with bounds at every external index on a larger product. Dividing by the Taylor factorials, composing coefficientwise majorants (\cref{ch17:lem:coefficient-resolution}), and using \eqref{ch17:eq:factorial-series}, we bound the degree-$k$ scalar coefficients, rooted at any free index, by $D_JH_J^k$ for fixed $D_J,H_J$. An additional multiplicity $\eta_j\le k$ is bounded by $2^k$. For an incoming function, the strict image bound gives the jet bound
\begin{equation}\label{ch17:eq:residual-jets}
 \norm{D^jR}_{\mathrm{root}}\le(E+\sigma_0)\rho^{2-j}\frac{j!}{(1-\kappa)^{j+1}},\qquad\kappa\le\tfrac12,\quad j\ge1,
\end{equation}
at a point whose image under the absolute reconstruction majorant is at most $\kappa\rho$. To verify it, we fix the first differentiated input index. At original degree $l$, the fixed-index weight in the input norm and the sum over the other $j-1$ differentiated positions give the factor $(l)_j$, with the remaining factor $\kappa^{l-j}\rho^{2-j}$, and the second identity in \eqref{ch17:eq:factorial-series} sums the result. The torus seminorm starts at degree two and the boundary norm at degree one; the same estimate bounds both. Constants of $R$ disappear only in the fluctuation difference. The input networks are joined to the reconstruction networks, and a prescribed differentiated index stays fixed. This proves \eqref{ch17:eq:residual-jets} in the coefficient algebra, before the background is evaluated.

Consider a monomial of $T_i$, and let $\eta_j$ be its degree in the free coordinates of scale $j$. It has at least three indices at scale $i$, so $\eta_i\ge3$ and $\sum_{j\ge i}\eta_j=k\ge3$. After division by $g_i^2$ its coupling factor is
\[
 t^{k-2}\lambda_i^{\eta_i-2}\prod_{j>i}\lambda_j^{\eta_j}.
\]
The product of ratios is at most one; the weaker bound $\lambda_*^{-2}$ would also be an admissible fixed constant. By \eqref{ch17:eq:free-difference}, every Haar, determinant, and incoming difference has at least one free index. A term of degree $k$ there has the factor $t^k\prod_j\lambda_j^{\eta_j}$, whose product of ratios is again at most one. The values $B_i$ in \eqref{ch17:eq:potential} are evaluated at the coarser reconstructions, and their coefficients in the coarser free variables belong to the same series. Summing over at most $J$ scales, increasing $A\ge\max\{1,2H_J,2D_JH_J\}$, and decreasing $h_0$ if necessary, we obtain
\[
 \norm{P_k}_{a',1}\le A^k,\qquad\norm{Q_k}_{a',1}\le C_{V,J}(1+E+\sigma_0)A^k.
\]
The first quantum coefficient $Q_1$ includes every degree-one term of the incoming interaction. Only powers of the fixed radii were absorbed in these constants. We choose $s_W<h_0$ with a strict margin in the larger domain. The series are then normally convergent on that radius, at the outer profile as well as after evaluation.

\paragraph{Step 3: the enclosing product is admissible.}
The three inequalities in \eqref{ch17:eq:enclosure-budgets} imply that
\[
 \delta_{\mathrm{en}}\le\lambda_*\Rwhite/2,\qquad\Ren\le\Rwhite.
\]
For $w\in\mathcal K_{\mathrm{en}}$ we have $\abs{\zeta_i}=t\lambda_i\abs{w_i}\le\Ren$. The minimum defining $\Rwhite$ places this whole product strictly inside the unperturbed and perturbed sequential products and the joint neighborhood; the constant $c_{\mathrm{tube},J}$ includes every finite sequential embedding. Indeed, the incoming displacement term is still $2\delta_V$ at the preliminary bound $\sigma_0$, and the new free radius is at most $\Ren\le Q$. Then \eqref{ch17:eq:sequential-displacement} implies that $C_{\mathrm{seq},J}(2\delta_V+\Ren)<\varepsilon_{\mathrm{tube}}/2$ at every coarser argument. The separate margin of the larger construction domain in \cref{ch16:lem:strict-two-radius-images} includes the background, the frozen values, and the free radius; it places all reconstructions before the gauge transformation in the larger charts of the Haar densities and of the action.

\paragraph{Step 4: the image increments.}
For either radius, we abbreviate $C_T=C_{T,\bullet}$ and $C_\Delta=C_{\Delta,\bullet}$. By hypothesis, we have $C_T(K_*\rho_\bullet+\Ren)\le1/2$. The rational increment in the statement is therefore at most $(2C_T+C_\Delta)\Ren\le\rho_\bullet/8$. Adding the bound $A_{0,\bullet}\le\rho_\bullet/4$ at zero fluctuation, we obtain the bound $3\rho_\bullet/8$ at exponent zero. If these inequalities hold in a weighted coefficient row, each calculation holds in that row. Otherwise apply the following interpolation to the unperturbed image row at such an exponent $c$, with $S_c^0\le C_J\rho_\bullet$. For nonnegative coefficient weights $b_\eta$ and network lengths $l_\eta$, H\"older's inequality for finite sums gives
\[
 \sum_\eta b_\eta e^{\vartheta cl_\eta}\le\Bigl(\sum_\eta b_\eta\Bigr)^{1-\vartheta}\Bigl(\sum_\eta b_\eta e^{cl_\eta}\Bigr)^\vartheta.
\]
Passing to increasing finite sums proves the stated interpolation. Its right side, divided by $\rho_\bullet$, tends to $3/8$ as $\vartheta\downarrow0$. Hence one positive $\vartheta$ works for both radii and the finitely many stages of the unperturbed construction. Its exponent is fixed before the regional rates are chosen. At allocated working rates below this exponent and satisfying \cref{ch16:thm:controlled-perturbation}, that theorem bounds the coefficient row of the perturbed-minus-unperturbed reconstruction by $C_{P,\bullet}\sigma_*\rho_{\mathrm{cl}}$ on the same enclosing products. Fixing this constant before the final choice of $\sigma_*$ and making that bound smaller than the remaining half-radius margin gives the asserted perturbed image bound. No interpolation at the unperturbed exponent is applied to inherited perturbative networks.

\paragraph{Step 5: curvature.}
For the curvature bound, we fix one derivative at the scalar coordinate $j$ and sum over the second differentiated index. Differentiating a homogeneous monomial of degree $k$ twice gives the factor $k-1$ relative to the norm $\norm{\cdot}_{a',1}$. The matrix weight is bounded by the weight of the original network, since the distance between the two differentiated indices is at most the network length. At the radius $\varrho=\Ren$ in the variables $\zeta$, the two contributions are bounded by
\begin{align*}
 \sum_{k\ge3}(k-1)A^k\varrho^{k-2}&=A^2\bigl[(1-A\varrho)^{-2}-1\bigr]\le\frac{2A^3\varrho}{(1-A\varrho)^2},\\
 t^2\mathcal L_V\sum_{k\ge2}(k-1)A^k\varrho^{k-2}&=\frac{A^2t^2\mathcal L_V}{(1-A\varrho)^2}.
\end{align*}
Here we used that on $\mathcal K_{\mathrm{en}}$ the coordinates satisfy $\abs{w_j}\le\Ren/t$, so that $t^{k-2}(\Ren/t)^{k-2}=\Ren^{k-2}$ for the classical terms and $t^k(\Ren/t)^{k-2}=t^2\Ren^{k-2}$ for the others. We enlarge $C_0\ge2A^3$ to include the fixed conversion between scalar and three-component block norms. The column bound follows from the symmetry of the Hessian. The weaker denominator $(1-A\Ren)^{-3}$ in \eqref{ch17:eq:curvature} applies to both contributions. Since $A\Ren\le1/4$, $C_0\Ren\le1/16$, and $C_0t^2\mathcal L_V\le1/16$, the result is
\[
 \frac{1/16+1/16}{(1-1/4)^3}=\frac8{27}<\frac12.
\]
All inequalities have strict margins in the charts, so they hold on a neighborhood of the real product, including its faces and corners. The constant $C_1$ of \cref{ch17:lem:ambient-white-comparison} is determined from $A$, $a'$, the label density, and the curvature bound just proved, so the second inequality in \eqref{ch17:eq:coupling-budgets} is a separate requirement compatible with the present choices.

\paragraph{Step 6: logarithmic cutoffs.}
A ball of radius $p_i$ lies in $\mathcal K_{\mathrm{en}}$ whenever $tp_i<\Ren$. The finite logarithmic offset in the statement bounds every such $p_i$, since $p(g_i)=p(t)+A_{\mathrm w}\log(1/\lambda_i)$ and $\lambda_i\ge\lambda_*$. This proves the inclusion both for the common cutoff and for the cutoffs at the individual scales. All constants were computed at $\sigma_0$ on fixed larger products, before any final enclosure quantity or width. The proof involves no derivative in the frozen values, no sum over earlier renormalization group steps, and no dependence on the individual radii of the fluctuation balls.
\end{proof}

\begin{remark}\label{ch17:rem:profile-jets}
The incoming boundary function $f$ in \cref{ch17:lem:complete-white-potential} is measured in the norm $\mathcal X^\rho$ with the uniform profile $\rho$, and \eqref{ch17:eq:residual-jets} is stated for this case. For a function measured with a nonuniform profile $\phi$, the analogous estimate holds for the rescaled function $\widetilde R(x)=R((\phi_ix_i)_i)$: if its degree-wise rooted sum at unit profile is at most $H\rho^2$, then at a point whose image under the absolute reconstruction majorant is bounded by $\kappa\phi$, $0\le\kappa<1$,
\[
 \abs{D^j\widetilde R}_{\kappa,\mathrm{root}}\le H\rho^2\sum_{l\ge j}(l)_j\kappa^{l-j}=H\rho^2\frac{j!}{(1-\kappa)^{j+1}},\qquad j\ge1.
\]
Thus each differentiated slot carries its own radius $\phi_i$. The unscaled form \eqref{ch17:eq:residual-jets} does not follow for a spatially varying radius profile. Indeed, fix an unfixed index $i$ at distance $d_i$ from $\Omega^c$. The function $R(v)=H\rho^2e^{-\awall d_i}v_i/\phi_i$ has norm $H$ and derivative $H\rho^2e^{-\awall d_i}/\phi_i$, which is unbounded as the boundary radius $\phi_i=\rho_\partial$ decreases with $d_i$ fixed. In the coefficient sums of \cref{ch17:sec:expansion} every differentiated background index is multiplied by its own radius, so no inverse power of the small boundary radius $\rho_\partial$ enters.
\end{remark}

\begin{lemma}\label{ch17:lem:ambient-white-comparison}
The following statements hold.
\begin{enumerate}
\item Let $\mathcal K$ be a finite product of three-dimensional balls, let $S,\mathcal W\subset\mathcal K$ be Borel sets, let $\chi:\mathcal K\to[0,1]$ be Borel, let $e^{-V}$ be a positive Borel function with $0<Z_{\mathcal K}=\int_{\mathcal K}e^{-V}\,dw<\infty$, and let $P>0$ be a constant prefactor (unrelated to the perturbation $P$ of \cref{ch17:thm:perturbed-hessian-coarea}). Set
\[
 G=P\int_S\chi e^{-V}\,dw,\qquad\widetilde F=P\int_{\mathcal W}e^{-V}\,dw,\qquad d\nu_{\mathcal K}=Z_{\mathcal K}^{-1}\one_{\mathcal K}e^{-V}\,dw.
\]
If $\nu_{\mathcal K}(\mathcal W^c)\le\delta<1$, then $G\le(1-\delta)^{-1}\widetilde F$. Nonnegative functions of the coarse background and external factors may multiply both sides. No inclusion between $\mathcal W$ and $S$ is needed. A product with no factors has mass one.
\item Let $\mathcal K=\mathcal K_{\mathrm{en}}$ be the enclosing product of \cref{ch17:lem:complete-white-potential}, under its hypotheses with $t_{\max}\le t_*$, at an admissible real coarse background and fixed real frozen values, and let $V$ be the potential \eqref{ch17:eq:potential}. There is a constant $C_1$, depending only on $A$, $a'$, and $v_J$, such that
\[
 D_w^2V\ge I/2,\qquad\sup_j\abs{\E_{\nu_{\mathcal K}}w_j}\le C_1t_{\max}(\mathcal L_V+2),
\]
where $j$ runs over the scalar coordinates. For every scalar function $F$ that is $1$-Lipschitz for the Euclidean norm and every $v\ge0$,
\[
 \nu_{\mathcal K}(\abs{F-\E F}\ge v)\le2e^{-v^2/4}.
\]
Let $p(t)=A_{\mathrm w}\log(1/t)$, and choose radii $p_i\ge p(t_{\max})$ with $t_{\max}\max_ip_i<\Ren$, so that the product $\mathcal W=\prod_{i,e}\{\abs{w_{i,e}}\le p_i\}$ is contained in $\mathcal K$. If $C_1t_{\max}(\mathcal L_V+2)\le p(t_{\max})/(2\sqrt3)$, which holds for all sufficiently small $t_{\max}$ and is one of the requirements on the coupling bound $g_*$ in \cref{ch17:thm:controlled-kernel}, then
\[
 \nu_{\mathcal K}(\mathcal W^c)\le6\mathsf b\exp\{-p(t_{\max})^2/48\},\qquad\mathsf b\le C_Jm^3,
\]
where $\mathsf b$ counts all fluctuation blocks of the step, jointly over all components of $\Omega$ and all scales. These are means under $\nu_{\mathcal K}$ at fixed frozen values; no symmetry of the frozen values is assumed. The constants have the dependence described in item (3) of \cref{ch01:conv:constants} and are independent of the later widths.
\item Consider the stopping sides $m_t=mB^t$, $t\ge0$, together with at most one additional side $\eta mB^j$, where $\eta$ ranges over a fixed finite set of positive ratios (arising from a final composition of fewer than $J$ averaging steps). Suppose that nonnegative errors at these sides satisfy
\[
 \delta_t\le C_Jm_t^3e^{-a_{\mathrm w}\log^2\sqrt{Mm_t/\gamma}},\qquad z_t\le C_Jm_t^3e^{-bm_t/\gamma},\qquad a_{\mathrm w},b>0,
\]
and similarly at the additional side. By part (2), the first bound holds with $a_{\mathrm w}=A_{\mathrm w}^2/48$ for the comparison errors of this lemma. Then for every fixed $\Gamma\in(0,\infty)$, in particular for every $\Gamma=\gamma_0\le1$, there is a minimum side $m_0$, depending only on $\Gamma$, the fixed data, and the finite set of ratios, such that, uniformly for $m\ge m_0$ and $0<\gamma\le\Gamma$, all errors are at most $1/2$ and
\[
 \sum_t(\delta_t+z_t)<\infty,\qquad\log\prod_t(1-\delta_t)^{-1}\le2\sum_t\delta_t\le C_\delta.
\]
The sums and the logarithm tend to zero as $m\to\infty$, uniformly in $0<\gamma\le\Gamma$ and in the finite or infinite number of sides. The same conclusion holds for $\epsilon_t=\delta_t/(1-z_t)$ after increasing $m_0$. There is one comparison for the whole product at each stopping side, not a separate factor for each component of $\Omega^c$. The minimum side and $C_\delta$ have the dependence described in item (4) of \cref{ch01:conv:constants}; none of them depends on the number of steps or on the large-field decomposition.
\end{enumerate}
\end{lemma}

\begin{proof}
Step 1 proves part (1). Steps 2 to 6 prove part (2): we first check that the measure is well defined, then interpolate between the Gaussian product and the potential, bound the force of the interaction using fourth moments, integrate the derivative of the means along the interpolation, and finally bound the probability outside the chosen product. Step 7 reduces part (3) to \cref{ch17:lem:finite-ceiling-endpoint-gates}.

\paragraph{Step 1: the abstract comparison.}
The Borel hypotheses make both integrands measurable, and since $0\le\chi\le1$ and $0<Z_{\mathcal K}<\infty$, we have $0\le G,\widetilde F\le P Z_{\mathcal K}<\infty$. Since $S\subset\mathcal K$, we have $G\le P Z_{\mathcal K}$. On the other hand
\[
 \widetilde F=P Z_{\mathcal K}\nu_{\mathcal K}(\mathcal W)\ge P Z_{\mathcal K}(1-\delta).
\]
Dividing by the positive finite number $P Z_{\mathcal K}$ gives part (1); the common prefactor cancels before any estimate. Multiplying by nonnegative functions of the coarse background preserves the inequality pointwise. For a product with no factors we use the unit point mass, and the same proof applies.

\paragraph{Step 2: normalizability on the enclosing product.}
The potential of \cref{ch17:lem:complete-white-potential} is real and smooth on a neighborhood of the fixed compact product $\mathcal K_{\mathrm{en}}$. Its ball radii $\Ren/t_{\max}$ are positive and finite, so the product has positive finite Lebesgue volume. If $V_{\min}$ and $V_{\max}$ are the minimum and maximum of $V$ on $\mathcal K_{\mathrm{en}}$, then
\[
 0<e^{-V_{\max}}\operatorname{vol}(\mathcal K_{\mathrm{en}})\le Z_{\mathcal K_{\mathrm{en}}}\le e^{-V_{\min}}\operatorname{vol}(\mathcal K_{\mathrm{en}})<\infty.
\]
Thus the integrand is positive, Borel, and normalizable. The chosen product $\mathcal W$ is closed, hence Borel; the comparison in part (1) keeps its Borel hypotheses on $S$ and $\chi$.

\paragraph{Step 3: convexity and concentration.}
We write $t=t_{\max}$ and $H(w)=V(w)-\abs w^2/2$, and interpolate on the same product
\[
 V_\theta(w)=\abs w^2/2+\theta H(w),\qquad0\le\theta\le1.
\]
Each $V_\theta$ is real and smooth on the compact product, so the normalization bounds of Step 2 hold throughout the interpolation. By \eqref{ch17:eq:curvature} the Hessian perturbation is at most $\kappa=8/27$, so every $V_\theta$ has Hessian at least $(1-\kappa)I\ge I/2$; for $\theta=1$ this is the first assertion. The absolute coefficient bound in \eqref{ch17:eq:curvature} bounds the exponentially weighted sums over $y$ of the suprema over the product of $\abs{(V_\theta)_{xy}}$, in rows and in columns, by the same fixed constant; the three-component block form follows by including the fixed block conversion in $C_0$. Fix $0<b\le a'$ so small that
\[
 \max\Bigl\{\sup_x\sum_{y\ne x}(e^{b\dist(x,y)}-1)\sup\norm{(V_\theta)_{xy}},\ \sup_y\sum_{x\ne y}(e^{b\dist(x,y)}-1)\sup\norm{(V_\theta)_{xy}}\Bigr\}\le\tfrac14.
\]
Such a choice depends only on the fixed weighted bound, since the convexity of $u\mapsto e^u$ gives $e^{bd}-1\le(b/a')(e^{a'd}-1)$ for $b\le a'$. The off-diagonal blocks of the Gaussian part vanish. Thus all hypotheses of \cref{ch02:thm:reflecting-concentration}, which adapts the Bakry--\'Emery and Helffer--Sj\"ostrand methods \cite{BakryEmery1985,HelfferSjostrand1994} to products of balls with reflecting boundary, hold with convexity constant $1/2$, on a neighborhood of the finite product; its ball faces and product corners are included in that theorem. It yields the scalar tail $2e^{-v^2/4}$ and the covariance bound
\begin{equation}\label{ch17:eq:ambient-covariance}
 \abs{\Cov_\theta(F,G)}\le4\sum_{x,y}e^{-b\dist(x,y)}\norm{\nabla_xF}_2\norm{\nabla_yG}_2,
\end{equation}
where $\Cov_\theta$ and $\norm{\cdot}_2$ refer to the normalized measure with density proportional to $e^{-V_\theta}$ on $\mathcal K_{\mathrm{en}}$. These constants do not depend on the number or the radii of the balls.

\paragraph{Step 4: the force.}
Let $M(\theta)=\max_j\abs{\E_\theta w_j}$, where $j$ is scalar (in this proof $M(\theta)$ is unrelated to the blocking factor $M$). The tail just obtained implies a universal bound on the fourth moment of $w_j-\E_\theta w_j$, by integrating $\int_0^\infty4u^3\Prob(\abs{w_j-\E_\theta w_j}>u)\,du$ against the Gaussian majorant. Hence
\[
 \norm{w_j}_4\le C_4+M(\theta).
\]
Consider a differentiated classical monomial of degree $k\ge3$. After fixing its differentiated coordinate, we place two of its remaining $k-1$ factors in $L^4$ and the other $k-3$ factors in the supremum norm $\Ren/t$ of the product. H\"older's inequality, which applies also when the two factors are the same scalar coordinate, gives the factor
\[
 t^{k-2}(\Ren/t)^{k-3}(C_4+M(\theta))^2=t\Ren^{k-3}(C_4+M(\theta))^2.
\]
The coefficient sum with the differentiated index fixed is bounded by $A^k$, and the geometric sum is at most $C(A)t(1+M(\theta)^2)$ because $A\Ren\le1/4$. For the quantum terms, including degree one, a direct supremum bound gives
\[
 \sum_{k\ge1}\mathcal L_VA^kt^k(\Ren/t)^{k-1}=\frac{\mathcal L_VAt}{1-A\Ren}.
\]
Combining the two bounds, we obtain
\begin{equation}\label{ch17:eq:ambient-force}
 \sup_j\norm{\partial_jH}_2\le C(A)t[\mathcal L_V+1+M(\theta)^2].
\end{equation}

\paragraph{Step 5: the means.}
Differentiating the finite integrals on the fixed compact product, we obtain
\[
 \frac{d}{d\theta}\E_\theta w_j=-\Cov_\theta(w_j,H).
\]
Indeed, the derivative of the numerator is $-\int w_jHe^{-V_\theta}$ and that of the denominator is $-\int He^{-V_\theta}$, and their quotient is the displayed centered expression. The integrands and their derivatives are bounded on this finite product, so the differentiation is justified directly. This is the calculation of \cref{ch10:thm:actual-means-and-scores} on a fixed product, carried out here for the present measure; we do not use the conclusions of that theorem, whose hypotheses concern the torus products. Combining \eqref{ch17:eq:ambient-covariance}, \eqref{ch17:eq:ambient-force}, and \eqref{ch17:eq:density-sum}, we obtain a fixed $C_1$ with
\begin{equation}\label{ch17:eq:mean-interpolation}
 \Bigl\lvert\frac{d}{d\theta}\E_\theta w_j\Bigr\rvert\le C_1t[\mathcal L_V+1+M(\theta)^2].
\end{equation}
It depends only on $A$, $a'$, and the label density; the margin $\kappa<1/2$ has already been imposed. This is the constant in the hypothesis \eqref{ch17:eq:coupling-budgets} of \cref{ch17:lem:complete-white-potential}.

At $\theta=0$ the measure is a product of centered truncated Gaussian measures on balls, so every coordinate has mean zero, for every frozen value. All means are continuous in $\theta$. Suppose that $M$ reached the value one, and let $\theta_1$ be the first time at which it does. Integrating \eqref{ch17:eq:mean-interpolation} for each coordinate on $[0,\theta_1]$, where $M\le1$, would give $M(\theta_1)\le C_1t(\mathcal L_V+2)<1/(4\sqrt3)<1$, a contradiction. Inserting $M<1$ into \eqref{ch17:eq:mean-interpolation} and integrating once more proves $M(1)\le C_1t(\mathcal L_V+2)$. This argument integrates the inequalities for the individual coordinates before taking the maximum; it does not use any differentiability of the maximum.

\paragraph{Step 6: the probability outside the chosen product.}
The condition $t\max_ip_i<\Ren$ puts the chosen product strictly inside $\mathcal K_{\mathrm{en}}$. For the logarithmic radii of \cref{ch17:thm:controlled-kernel} this condition follows from \eqref{ch17:eq:uniform-logarithmic-gates}; for other radii it is the explicit inclusion hypothesis. If a fluctuation block lies outside its ball of radius $p_i\ge p(t)$, then some scalar coordinate has absolute value larger than $p(t)/\sqrt3$. Since the mean of that coordinate is at most $p(t)/(2\sqrt3)$ in absolute value, the coordinate deviates from its mean by at least $p(t)/(2\sqrt3)$. By the scalar tail and a union bound over the three coordinates, the probability that a given block lies outside its ball is at most $6e^{-p(t)^2/48}$. The step has at most $\mathsf b\le C_Jm^3$ blocks: each of the at most $J$ scales has a fixed finite number of free labels per output block, and removing blocks from the region only decreases this number. This proves the bound on $\nu_{\mathcal K}(\mathcal W^c)$, with no factor depending on the number of components. The enclosure and curvature hypotheses are used on the whole product $\mathcal K_{\mathrm{en}}$, which does not depend on the configuration.

\paragraph{Step 7: summation.}
Part (3) is the case $C=C_J$, $c=b$, $a=a_{\mathrm w}$ of \cref{ch17:lem:finite-ceiling-endpoint-gates}, with the ratio $B$ and the finite set of additional ratios, since $\log^2\sqrt{Mm/\gamma}=z^2$ with $z=(1/2)\log(Mm/\gamma)$. With $p(t_{\max})=A_{\mathrm w}\log\sqrt{Mm/\gamma}$, the bound of part (2) is of the stated form with $a_{\mathrm w}=A_{\mathrm w}^2/48$, since $6\mathsf b\le C_Jm^3$.
\end{proof}

\section{Cluster expansion with boundary-localized coefficients}\label{ch17:sec:expansion}

We now expand the logarithm of the fluctuation integral, keeping two structures apart. Probabilistic independence is determined by the set of fluctuation blocks on which a monomial depends: two monomials depending on disjoint sets of blocks are independent under a product reference measure, even if they depend on the same background coordinates. Spatial decay is carried instead by a connected network through all positions of the coefficient, background positions included, which may lie far from the fluctuation blocks. Overlapping monomials share a fluctuation block, so their networks meet there, and the union of their networks is connected without any additional connecting path.

The organization into connected polymer activities and a convergent logarithm follows the constructive cluster expansion; see Ba\l aban \cite{Balaban1988RGII,Balaban1988Convergent} for its implementation in lattice gauge theories, Brydges \cite{Brydges1986} for a general exposition, and Koteck\'y and Preiss \cite{KoteckyPreiss1986} for the abstract polymer criterion. The cancellation of connected graphs against rooted trees used in \cref{ch09:thm:hard-core-log} is Penrose's partition \cite{Penrose1967}; its construction by generations and parents, and the positive tree recursions used below, are described in \cite[Sections~2 and~4.1]{FernandezProcacci2007}. The finite graph inequality is imported in the form of \cite[Appendix~B, (323), p.~42]{Dimock2013v2}, through \cref{ch09:thm:hard-core-log}; no other result of that paper, which concerns a scalar field theory, is used. We prove directly the features needed here: the two kinds of roots, the exact fluctuation supports, and the visits to the fixed complement.

\paragraph{Complex coupling parameter.}
In this section we replace the coupling powers in \eqref{ch17:eq:potential-series} by powers of a complex parameter $z$: the classical powers $t_{\max}^{q-2}$ become $z^{q-2}$ and the other powers $t_{\max}^q$ become $z^q$. The incoming functions, the couplings of earlier steps, the ratios $\lambda_i$, and the normalized symmetric reference measures on balls are kept fixed. The incoming insertion becomes
\[
 A_R(z)=R\bigl(\Log C(v,(z\lambda_iw_i)_i;\Rec)\bigr)-R\bigl(\Log C(v,0;\Rec)\bigr).
\]
We call the interaction without insertion the sum of all terms of $V-\abs w^2/2$ in \eqref{ch17:eq:potential} except the incoming insertion, after this substitution. For $p_x>0$ attached to the fluctuation blocks, with $p_x=p_i$ for a block at scale $i$, we write $p=\max_ip_i$, and we always assume $p\ge1$.

\begin{definition}\label{ch17:def:rooted-norms}
Let $F$ be an interaction, given with a resolution into scalar monomials
\[
 U_\lambda=u_\lambda v^{\nu_\lambda}w^{\eta_\lambda},\qquad\abs{\eta_\lambda}\ge1,
\]
as in \cref{ch17:lem:coefficient-resolution}. A fluctuation index is a whole fluctuation block at its scale and position; a background index is an unfixed scalar background coordinate on the lattice where it was introduced. Write $S_\lambda$ for the set of fluctuation blocks on which $\eta_\lambda$ is positive (the fluctuation support of the monomial) and $k_\lambda=\abs{\nu_\lambda}$ for its background degree. Each monomial carries a connected network through all of its fluctuation and background positions; for a monomial whose coefficient has a visit to $\Omega^c_{L_J}$, the network contains this visit. We write $\ell_\lambda$ for the total length of the network, including the visit. No enlargement of $S_\lambda$ is made. Internal indices are summed, using a fixed part of the exponential decay rate, before the resolution into monomials, and repeated occurrences keep their own network refinements. For the exponent $s_0>0$ set
\[
 b_\lambda(F)=\abs{u_\lambda(F)}\WideProfile^{\nu_\lambda}\prod_xp_x^{\abs{\eta_{\lambda,x}}}e^{s_0\ell_\lambda},
\]
and for $h\ge0$
\begin{align*}
 W_h(F)&=\sup_x\sum_{\lambda:x\in S_\lambda}e^{h\abs{S_\lambda}}b_\lambda(F),\\
 D_h(F)&=\sum_{k\ge1}\sup_i\sum_{\lambda:k_\lambda=k}\nu_{\lambda,i}e^{h\abs{S_\lambda}}b_\lambda(F).
\end{align*}
In $W_h$ the root $x$ is a fluctuation block, and in $D_h$ the root $i$ is a background index. In $D_h$ the supremum over the root is taken before the sum over the background degree, as in the coefficient norm with boundary decay.
\end{definition}

\begin{lemma}\label{ch17:lem:decorated-white-primitives}
Fix the admissible larger products and the preliminary perturbation ball of \cref{ch17:thm:perturbed-hessian-coarea,ch17:lem:complete-white-potential}, including both strict margins for the images under the reconstruction majorant and the inequalities on the conditional forces at every stage. Fix the ratio $K_*$ and the profiles $\InflatedProfile=K_*\phi$ and $\WideProfile=\sqrt{K_*}\phi$. The only restrictions on the rates needed here are a fixed allocation of the exponential decay rate available in the coefficient bounds of the potential and the restrictions on the transport from input to output of \cref{ch16:thm:controlled-perturbation,ch15:lem:wall-networks}; let $s_0=a+\alpha$ lie strictly below these bounds. Consider the following families of interactions, resolved into monomials as in \cref{ch17:def:rooted-norms}: the interaction without insertion, the incoming insertion, the differences between the regional families and the restricted torus families, and the mixed difference of the boundary change and the incoming insertion. Let $H=E+\mathcal X^\rho(f)$ and set
\[
 \mathfrak s(z)=\frac{C_J\abs z(p^3+p+Hp)}{1-C_J\abs zp},
\]
for $z$ in the fixed larger radius with $C_J\abs zp<1$, inside both incoming image margins. Then, uniformly in the fixed real frozen values,
\[
 W_{16}(F)+D_{16}(F)\le\epsilon_F
\]
for each of these families $F$. One may take $\epsilon_F\le C_J\mathfrak s(z)$ for the family without insertion and for the boundary differences, and
\begin{equation}\label{ch17:eq:decorated-inherited-gate}
 \epsilon_A+\epsilon_{\partial A}\le\frac{C_J\abs zHp}{1-C_J\abs zp}
\end{equation}
for the incoming insertion $A$ and its boundary difference $\partial A$. In particular the latter bound holds with a visit to $\Omega^c$ and with an arbitrary background root. The constants have the dependence described in item (3) of \cref{ch01:conv:constants} and are independent of the later widths; the profiles and the input radii are not decreased at later steps.
\end{lemma}

\begin{proof}
Step 1 records the coefficient resolution of the terms of the potential and its basic bounds. Steps 2 and 3 write out, with every contraction, the resolution in the two typical cases of a classical monomial and of a monomial of the Haar density. Step 4 treats the incoming insertion and the boundary differences. Step 5 converts the coefficient bounds into the bounds on $W_{16}$ and $D_{16}$, keeping the order of the supremum and the degree sum.

\paragraph{Step 1: the resolution.}
The tensors of \cref{ch17:thm:perturbed-hessian-coarea}, with bounds at every external index, and the coefficient version of \cref{ch17:lem:complete-white-potential} are available on the outer profile, before the background is evaluated. By \cref{ch17:lem:coefficient-resolution}, every term of the potential has an exact resolution into monomials carrying connected networks. Rooted at an arbitrary free index, this resolution has convergent majorants in the free degree $q$, namely
\[
 C_J^q\abs z^{q-2}p^q\quad(q\ge3)\qquad\text{and}\qquad C_J^q\abs z^qp^q\quad(q\ge1),
\]
for the classical terms and for the other terms, respectively. For the incoming insertion the second majorant is multiplied by $H$. The ratios $\lambda_i$ have modulus at most one, and $J$ is fixed. For the last assertion the input is \eqref{ch17:eq:residual-jets} with $E+\sigma_0$ replaced by $H$: fixing its first differentiated input and summing over the other inputs gives
\[
 \norm{D^jR}_{\mathrm{root}}\le H\rho^{2-j}\frac{j!}{(1-\kappa)^{j+1}},\qquad\kappa\le\tfrac12,
\]
by applying the factor $(l)_j\kappa^{l-j}$ to the input norm at degree $l$ with that input fixed and summing the geometric series. The constants of $R$ disappear in the fluctuation difference; they are not bounded by $H$. The composition with every index rooted in \cref{ch16:thm:controlled-perturbation}, with its volume-normalized input columns, keeps a specified external background index as well as a specified free index. At a background root $i$, the sum of the monomials containing that index is bounded by $\InflatedProfile(i)$ times the absolute coefficient majorant of $\partial_iF$. This produces the factor $\nu_{\lambda,i}$ in $D_h$. No division by the smaller boundary radius occurs.

\paragraph{Step 2: a classical monomial.}
Fix a scale $i$. By the definition of the partial minimization, $F_i(v,c_{i+1},y)$ is the perturbed objective $\mathcal Q_1$ of the proof of \cref{ch17:thm:perturbed-hessian-coarea}, in the free logarithmic coordinates and multiplied by the factor $M^i$, evaluated with the coordinates of the finer scales at their conditional minimizers, at fixed frozen values. We write this as $F_i(v,c_{i+1},y)=\mathcal Q^{(i)}(\Psi_i(v,c_{i+1},y))$, where $\mathcal Q^{(i)}=M^i\mathcal Q_1$, $y$ denotes the free logarithmic coordinates of scale $i$, $c_{i+1}=C_{i+1}(v,\zeta_{>i})$ is the coarser reconstruction, and $\Psi_i$ maps $(v,c_{i+1},y)$ to the full list of coordinates $(v,z^*_{<i},y,z_{>i})$, with the finer coordinates $z^*_{<i}$ at their conditional minimizers. The coefficient bounds for $\mathcal Q^{(i)}$ on the common product are those of \cref{ch16:thm:controlled-perturbation}, and those for $\Psi_i$, which contains the conditional minimizers of the finer stages, are those of \cref{ch17:thm:perturbed-hessian-coarea}. The factorization only displays the chain rule; the general term can equally be bounded through the tensors $D_y^rF_i$ of that theorem. At fixed $(v,c_{i+1})$ the minimizer $z_i^*$ satisfies $F_{i,y}(v,c_{i+1},z_i^*)=0$, and $K_i=F_{i,yy}(v,c_{i+1},z_i^*)$, $S_i=K_i^{-1/2}$, $z_i=z_i^*+g_iS_iw_i$. Since $S_i^{\mathsf T}K_iS_i=I$, Taylor's formula in $y$ on the common coefficient product gives
\[
 \frac{F_i(v,c_{i+1},z_i)-F_i(v,c_{i+1},z_i^*)}{g_i^2}-\frac{\abs{w_i}^2}2=\sum_{r\ge3}\frac{g_i^{r-2}}{r!}\sum_{a_1,\ldots,a_r}\sum_{\beta_1,\ldots,\beta_r}F_{i,a_1\cdots a_r}^*\prod_{l=1}^rS_i(a_l,\beta_l)w_{i,\beta_l}.
\]
The star denotes evaluation at $y=z_i^*$ (not at zero coarser input), and the $\beta_l$ are scalar labels at scale $i$. This expression is analytic at $g_i=0$: stationarity and the exact subtraction of the quadratic term remove the apparent singularity.

We display the term $r=3$. For a multi-index $\tau=(\nu,\eta_{>i})$ in the background variables and the coarser free variables, we write $[\,\cdot\,]_\tau$ for the coefficient of $v^\nu\zeta_{>i}^{\eta_{>i}}$ after the substitution $c_{i+1}=C_{i+1}(v,\zeta_{>i})$, and set
\[
 e^{(m)}_{x_1\cdots x_m;\tau}=[(D^m\mathcal Q^{(i)})_{x_1\cdots x_m}(\Psi_i^*)]_\tau,\qquad
 P_{x;A;\tau}=[(\partial_{y_A}\Psi_{i,x})^*]_\tau,\qquad
 s_{a\beta;\tau}=[S_i(a,\beta)]_\tau,
\]
where $\partial_{y_A}=\partial_{y_a}\partial_{y_b}$ for $A=(a,b)$, and similarly for longer $A$. The chain rule gives, for $T_{abc;\tau}=[F_{i,abc}^*]_\tau$,
\begin{align*}
 T_{abc;\tau}={}&\sum_{x_1,x_2,x_3}\sum_{\tau_0+\tau_1+\tau_2+\tau_3=\tau}e^{(3)}_{x_1x_2x_3;\tau_0}P_{x_1;a;\tau_1}P_{x_2;b;\tau_2}P_{x_3;c;\tau_3}\\
 &+\sum_{x_1,x_2}\sum_{\tau_0+\tau_1+\tau_2=\tau}e^{(2)}_{x_1x_2;\tau_0}\bigl(P_{x_1;ab;\tau_1}P_{x_2;c;\tau_2}+P_{x_1;ac;\tau_1}P_{x_2;b;\tau_2}+P_{x_1;bc;\tau_1}P_{x_2;a;\tau_2}\bigr)\\
 &+\sum_x\sum_{\tau_0+\tau_1=\tau}e^{(1)}_{x;\tau_0}P_{x;abc;\tau_1}.
\end{align*}
In particular, the Hessian is
\[
 F_{i,ab}^*=\sum_{x,x'}\mathcal Q^{(i)}_{xx'}(\Psi_i^*)\Psi_{i,x,a}^*\Psi_{i,x',b}^*+\sum_x\mathcal Q^{(i)}_x(\Psi_i^*)\Psi_{i,x,ab}^*.
\]
Stationarity in $y$ does not imply $(\mathcal Q^{(i)})'(\Psi_i^*)=0$, so the term with $(\mathcal Q^{(i)})'$ and the last line of the formula for $T_{abc;\tau}$ cannot be deleted. Let $\eta_i$ be a multi-index of degree three in the fluctuation variables of scale $i$. After the substitutions $g_l=z\lambda_l$ and $\zeta_{>i}=z\lambda_{>i}w_{>i}$, the coefficient of $v^\nu w_i^{\eta_i}w_{>i}^{\eta_{>i}}$ in this cubic term is
\begin{equation}\label{ch17:eq:cubic-primitive}
 \frac{z^{1+\abs{\eta_{>i}}}\lambda_i}6\prod_{j=i+1}^{s-1}\lambda_j^{\abs{\eta_j}}\sum_{a,b,c}\sum_{\substack{\beta_1,\beta_2,\beta_3\\ \mathbf e_{\beta_1}+\mathbf e_{\beta_2}+\mathbf e_{\beta_3}=\eta_i}}\sum_{\tau_0+\tau_1+\tau_2+\tau_3=(\nu,\eta_{>i})}T_{abc;\tau_0}s_{a\beta_1;\tau_1}s_{b\beta_2;\tau_2}s_{c\beta_3;\tau_3},
\end{equation}
where $\eta_j$ is the part of the multi-index $\eta_{>i}$ at scale $j$. The factor $1/6$ also accounts for repeated $\beta$'s; no further factorial occurs. For each summand we take the networks of its coefficients and join them at their common indices $x$ and $a$, and at the coarser arguments. The network of the resulting monomial is this connected union, and its fluctuation support is
\[
 S_\lambda=\{\text{blocks of }\beta_1,\beta_2,\beta_3\}\cup\{\text{coarser blocks used by }\eta_{>i}\},
\]
with repeated blocks counted once. The background positions in $\nu$ belong to the network even if they lie outside $S_\lambda$. Let $\mathcal Q_m^+$, $P_m^+$, and $S^+$ be the absolute scalar majorants of the jets and maps above on their common product, including the volume ratios of part (4) of \cref{ch17:lem:coefficient-resolution}. The formula for $T_{abc;\tau}$ and \cref{ch17:lem:coefficient-resolution} give the cubic majorant
\[
 \tfrac16\bigl(\mathcal Q_3^+(P_1^+)^3+3\mathcal Q_2^+P_2^+P_1^++\mathcal Q_1^+P_3^+\bigr)(S^+)^3,
\]
and the same expression with a derivative falling on one factor bounds the sum rooted at a fluctuation or background index. Each product here corresponds to a connected contraction in the formula for $T_{abc;\tau}$, so no unrooted product over an extensive set appears. For general $r$ the rule is the partition formula
\[
 \partial_{a_1}\cdots\partial_{a_r}(\mathcal Q^{(i)}\circ\Psi_i)=\sum_{\pi}\sum_{(x_B)_{B\in\pi}}(D^{\abs\pi}\mathcal Q^{(i)})_{(x_B)}(\Psi_i)\prod_{B\in\pi}\partial_{a_B}\Psi_{i,x_B},
\]
where $\pi$ runs over the partitions of $\{1,\ldots,r\}$. Together with \cref{ch17:lem:coefficient-resolution}, it gives a convergent coefficientwise majorant on a common smaller free product. Since $r\ge3$, a term of total free degree $q=r+\abs{\eta_{>i}}$ has the coupling power $z^{q-2}$, and the ratios of the scales have modulus at most one. This proves the cubic starting order before any expectation over the fluctuation variables is taken.

\paragraph{Step 3: a monomial of the Haar density.}
Let $\Gamma(v,c_{i+1},y)$ be the chart of solved links at scale $i$, let $\mathcal E_i^\circ$ be its links outside the block axial gauge trees, and let $\mathfrak c$ run over the set $\mathfrak C_i$ of constraints solved at scale $i$, with eliminated edges $e^*_{\mathfrak c}$. We write
\[
 A_{\mathfrak c}=D_{X_{e^*_{\mathfrak c}}}\mathsf C_{\mathfrak c}(\Gamma),\qquad L_i=-\sum_{e\in\mathcal E_i^\circ}\log j_H(\Gamma_e)+\sum_{\mathfrak c\in\mathfrak C_i}\log\det A_{\mathfrak c},
\]
with $j_H(X)=\sinc^2(\abs X/2)$ as in \cref{ch01:def:block-map}; only field-independent factors have been removed. The fluctuation-dependent part of $L_i$ at scale $i$ is
\[
 L_i(v,c_{i+1},z_i)-L_i(v,c_{i+1},z_i^*)=\sum_{r\ge1}\frac{g_i^r}{r!}\sum_{a_1,\ldots,a_r}\sum_{\beta_1,\ldots,\beta_r}L_{i,a_1\cdots a_r}^*\prod_{l=1}^rS_i(a_l,\beta_l)w_{i,\beta_l}.
\]
Its linear coefficient is
\[
 L_{i,a}^*=-\sum_{e\in\mathcal E_i^\circ}\sum_{\varsigma=1}^3(\partial_\varsigma\log j_H)(\Gamma_e^*)\Gamma_{e\varsigma,a}^*+\sum_{\mathfrak c\in\mathfrak C_i}\sum_{A,B=1}^3(A_{\mathfrak c}^{-1})_{BA}^*(\partial_{y_a}A_{\mathfrak c})_{AB}^*,
\]
where $(\partial_{y_a}A_{\mathfrak c})_{AB}=\sum_{e,\varsigma}(\partial_{X_{e\varsigma}}A_{\mathfrak c})_{AB}\Gamma_{e\varsigma,a}$, and all factors are evaluated at the minimizer and at the coarser reconstruction $C_{i+1}(v,\zeta_{>i})$. Taking the coefficient $[\,\cdot\,]_\tau$ of each factor and summing over all allocations of $\tau$ gives the coefficients $l_{a;\tau}=[L_{i,a}^*]_\tau$, with every link, constraint, color, integrated coordinate, and multi-index sum displayed. For $\abs{\eta_i}=1$ the coefficient after the coupling substitution is
\[
 z^{1+\abs{\eta_{>i}}}\lambda_i\prod_{j=i+1}^{s-1}\lambda_j^{\abs{\eta_j}}\sum_a\sum_{\beta:\mathbf e_\beta=\eta_i}\sum_{\tau_0+\tau_1=(\nu,\eta_{>i})}l_{a;\tau_0}s_{a\beta;\tau_1}.
\]
The trace in $L_{i,a}^*$ is rooted at the index $a$ through the derivative $\partial_{y_a}A_{\mathfrak c}$; its coefficient bound sums over $\mathfrak c,e,A,B,\varsigma$. It is not the unrooted norm of a sum of logarithmic determinants. The network is the connected union of the networks of the Haar factors, the inverse matrices, and the chart, joined to the network of $S_i$ at $a$, and the support is the block of $\beta$ together with the coarser blocks used by $\eta_{>i}$. \Cref{ch17:lem:coefficient-resolution} bounds these coefficients by the products of the majorants of the factors and their derivatives. The series starts at total free degree one, with coupling power equal to the total free degree.

The remaining part of the measure fluctuation in \eqref{ch17:eq:measure-fluctuation} is the difference $B_i(c_{i+1})-B_i(\widehat c_{i+1})$, whose fluctuation support may lie entirely in coarser scales; it is not part of the field-independent prefactor. For a coordinate $h_\gamma$ of the coarser argument, the total derivative is
\[
 \partial_{h_\gamma}B_i=L_{i,h_\gamma}^*+\sum_aL_{i,a}^*\partial_{h_\gamma}z_{i,a}^*+\frac12\sum_{a,b}(K_i^{-1})_{ba}\partial_{h_\gamma}K_{i,ab},
\]
where $\partial_{h_\gamma}z_i^*=-K_i^{-1}F_{i,yh_\gamma}^*$ and $\partial_{h_\gamma}K_{i,ab}=F_{i,abh_\gamma}^*+\sum_{a'}F_{i,aba'}^*\partial_{h_\gamma}z_{i,a'}^*$. The difference equals the convergent coefficient expansion of
\[
 \int_0^1\sum_{j=i+1}^{s-1}\sum_{\beta\in I_j}\zeta_{j,\beta}\sum_\gamma(\partial_{h_\gamma}B_i)\bigl(v,C_{i+1}(v,\tau\zeta_{>i})\bigr)\,(\partial_{\zeta_{j,\beta}}C_{i+1,\gamma})(v,\tau\zeta_{>i})\,d\tau,
\]
since the integrand is the derivative in $\tau$ of $B_i(v,C_{i+1}(v,\tau\zeta_{>i}))$; here $I_j$ is the set of unfixed scalar free labels at scale $j$, and $\partial_{\zeta_{j,\beta}}C_{i+1}$ is evaluated at $\tau\zeta_{>i}$ without an extra factor $\tau$. Each Taylor coefficient of this integral carries the factor $1/(n+1)$ when the remaining integrand has degree $n$ in $\tau$. Likewise, at fixed couplings,
\[
 dz_i=(\partial_vz_i^*)dv+(\partial_cz_i^*)dc_{i+1}+g_i\{(\partial_vS_i)dv+(\partial_cS_i)dc_{i+1}\}w_i+g_iS_i\,dw_i,
\]
and none of these terms may be dropped. The inverse square root $S_i$ is treated as a constituent tensor with the bounds of \cref{ch17:thm:perturbed-hessian-coarea}; if it is expanded, \eqref{ch17:eq:square-root} and the Neumann series of the resolvent around $K_i(0,0)$ give a chain of ordinary matrix products whose coefficients are its network refinements. The comparison of determinant matrices with the same indices also includes the correction \eqref{ch17:eq:restricted-log}, which is a further constituent tensor with a visit to $\Omega^c$.

\paragraph{Step 4: the incoming insertion and the boundary differences.}
We write $\mathcal A(v,\zeta)=\Log C(v,\zeta;\Rec)$. Then
\[
 R(\mathcal A(v,\zeta))-R(\mathcal A(v,0))=\sum_{n\ge1}\frac1{n!}D^nR(\mathcal A(v,0))\bigl[\mathcal A(v,\zeta)-\mathcal A(v,0),\ldots,\mathcal A(v,\zeta)-\mathcal A(v,0)\bigr],
\]
and each factor $\mathcal A(v,\zeta)-\mathcal A(v,0)$ has positive degree in $\zeta$. The coefficient of $v^\nu\zeta^\eta$ is the sum over $n$, over the contracted indices $j_1,\ldots,j_n$ of the old lattice, and over the allocations of the multi-indices to the $n$ factors with positive free degree in each, of the products of the jet coefficients of $R$ and the coefficients of the factors. No additional multinomial factor occurs, since each coefficient is already a multi-index coefficient. Each such occurrence keeps the union of the old network of $R$, scaled by $B^{-1}$ in output units, and the attached reconstruction networks, and its fluctuation support is the set of blocks used by $\eta$. By the jet bound of Step 1 with $j=1$ at an image bounded by $\rho/2$, the derivative series of $R$ at the image is at most $4H\rho$. The column majorant of the reconstruction with a free index fixed vanishes at zero free radius and is linear in the free radius on the fixed product, so it is at most $C_J\abs zp$. Fixing an external input in the old counting convention contributes the fixed volume ratio $\Gamma_J=\max_iv_i/v_{\mathrm{old}}\le M^{3J}$ between the input mesh and the old mesh. We obtain the incoming bound, linear in $H$ and starting at $\abs zp$.

For a difference between the regional and the restricted torus constructions we compare factor by factor, before any expansion of logarithms. In each product of coefficients we use the ordered identity
\[
 \prod_{r=1}^mA_r^1-\prod_{r=1}^mA_r^0=\sum_{r=1}^m\Bigl(\prod_{l<r}A_l^1\Bigr)(A_r^1-A_r^0)\Bigl(\prod_{l>r}A_l^0\Bigr),
\]
and for the incoming function at two reconstructed arguments $C_0,C_1$ in the identity chart (here $C_0,C_1$ denote arguments, not the constants of \cref{ch17:lem:complete-white-potential}),
\[
 R(C_1)-R(C_0)=\int_0^1DR\bigl(C_0+s(C_1-C_0)\bigr)[C_1-C_0]\,ds.
\]
Both images under the absolute reconstruction majorant lie in the common product with half the incoming radius, so the whole segment does. Terms involving the incoming boundary function use its incoming boundary weight; torus terms acquire the visit carried by the difference of the reconstructions. Fixed input slots are visits to $\Omega^c$. Changed formulas, minimizers, inverse square roots, and Haar densities are treated by the same identity, with $-A^{-1}(A-B)B^{-1}$ for inverses and \eqref{ch17:eq:square-root} for inverse square roots. Each term contains one changed factor with its visit to $\Omega^c$. The versions with every index rooted in \cref{ch16:thm:controlled-perturbation} keep every differentiated background or free index, even when the changed factor is not differentiated; the degree-one terms are therefore included.

We join the incoming networks, scaled by $B^{-1}$, to the chart networks at their contracted indices. The transport of the route to the fixed complement is exactly \eqref{ch16:eq:network-transport}, and \cref{ch15:lem:wall-networks} sums the remaining internal indices. We append the resulting route to the network of the monomial. One fixed allocation of the available exponents bounds $e^{s_0\ell_\lambda}$, including this route and the fixed stencil offsets; the offsets contribute a fixed factor in the degree, absorbed on the margin already provided by the construction. The allocation is the same for every order and every contraction, and no exponent is lost at each vertex. The scaled incoming networks and routes are bounded by the inequalities against $Ba$ and $B\alpha$. Every term still contains an incoming factor, so its bound is linear in $H$ and starts at $\abs zp$. The arguments of Step 5 apply to this majorant with a visit to $\Omega^c$ and give \eqref{ch17:eq:decorated-inherited-gate}. Finally, a torus monomial deleted because it uses a missing fluctuation block already has a visit at that block; if it has an unfixed root, its network provides the route to $\Omega^c$.

\paragraph{Step 5: conversion to $W_{16}$ and $D_{16}$.}
By Steps 1 to 4, there exist a free radius $V_1>0$, strictly inside all image and stage margins, and a common bound $A_0$ for the coefficientwise majorants of the diagrams on the outer profile $\InflatedProfile$ and at free radius $V_1$, including each permitted fixed coefficient index and the volume ratios of part (4) of \cref{ch17:lem:coefficient-resolution}. For the incoming family the bound is $A_0H$, and for a difference it is linear in the bound of the changed factor. We write the resolved coefficients of free degree $q$ and background degree $k$ as $u^0_\lambda v^{\nu_\lambda}\zeta^{\eta_\lambda}$, after removing the factor $z^{q-2}$ for the classical terms and $z^q$ for the others; the ratios $\lambda_i$ have modulus at most one. The mixed factorial identity is
\[
 \frac1{k!\,q!}\sum_{\text{other ordered indices}}\bigl\lvert D_v^kD_\zeta^q(u^0v^\nu\zeta^\eta)\bigr\rvert=\begin{cases}(\eta_l/q)\abs{u^0},&\text{one free input }l\text{ fixed},\\(\nu_i/k)\abs{u^0},&\text{one background input }i\text{ fixed},\end{cases}
\]
for $q\ge1$ and $k\ge1$ respectively, and the radii of all inputs multiply these expressions. With $\widetilde C=\max\{1,V_1^{-1}\}$, \cref{ch17:lem:coefficient-resolution} and the finite majorant at $V_1$ give
\[
 \sup_l\sum_{\lambda:(k_\lambda,q_\lambda)=(k,q)}\frac{\eta_{\lambda,l}}q\abs{u^0_\lambda}\InflatedProfile^{\nu_\lambda}e^{s_0\ell_\lambda}\le A_0\widetilde C^q,\qquad
 \sup_i\sum_{\lambda:(k_\lambda,q_\lambda)=(k,q)}\frac{\nu_{\lambda,i}}k\abs{u^0_\lambda}\InflatedProfile^{\nu_\lambda}e^{s_0\ell_\lambda}\le A_0\widetilde C^q.
\]
These bounds concern the total variation of the resolution produced by composition of coefficientwise majorants.

Let $\vartheta=K_*^{-1/2}<1$. Passing from $\InflatedProfile$ to $\WideProfile$ multiplies a term of background degree $k$ by exactly $\vartheta^k$, independently of the positions of its background indices in the spatially varying radius profile. Each term has $\abs{S_\lambda}\le q$, and selecting a fluctuation block is bounded by the sum of its three scalar multiplicities. At fixed $k,q$, the preceding bounds therefore contribute at most $3qA_0\widetilde C^q\vartheta^k$ for a fluctuation root and $kA_0\widetilde C^q\vartheta^k$ for a background root; the latter bound is uniform in the root at each degree. Taking the supremum over the root at each degree and only then summing over $k$, and using $q\le2^q$, we find that the sums in the definitions of $W_h$ and $D_h$, at fixed $q$, are at most
\[
 A_0\widetilde C^qe^{hq}\Bigl(\frac{3q}{1-\vartheta}+\frac\vartheta{(1-\vartheta)^2}\Bigr)\le\mathsf A\mathsf B^q,\qquad\mathsf A=A_0\Bigl(\frac3{1-\vartheta}+\frac\vartheta{(1-\vartheta)^2}\Bigr),\quad\mathsf B=2e^h\widetilde C.
\]
Multiplying by $\abs z^{q-2}p^q$, $q\ge3$, or by $\abs z^qp^q$, $q\ge1$, and summing the geometric series, we obtain on $\mathsf B\abs zp<1$
\[
 W_h^{\mathrm{cl}}+D_h^{\mathrm{cl}}\le\frac{\mathsf A\mathsf B^3\abs zp^3}{1-\mathsf B\abs zp},\qquad
 W_h^{\mathrm{Haar}}+D_h^{\mathrm{Haar}}\le\frac{\mathsf A\mathsf B\abs zp}{1-\mathsf B\abs zp},\qquad
 W_h^{\mathrm{inc}}+D_h^{\mathrm{inc}}\le\frac{\mathsf A\mathsf BH\abs zp}{1-\mathsf B\abs zp}
\]
for the classical, Haar-density, and incoming families. The same bounds hold for the differences with a visit to $\Omega^c$, with their majorants. Note that the classical bound starts with $\abs zp^3$: the two coupling powers removed by the division by $g_i^2$ are not restored. For $h=16$, enlarging the constant to a common $C_J$ and restricting the disc accordingly gives the assertions of the lemma. Choosing the auxiliary disc smaller by this fixed factor keeps it inside the same larger free radius, and neither an input radius nor the ratio $K_*$ changes. Finitely many additional background marks contribute a polynomial factor in $k$, which is summable since $\sum_k(k+1)^m\vartheta^k<\infty$. The argument never interchanges $\sup_i\sum_k$ with $\sum_k\sup_i$. The bounds at fixed real frozen values are proved first and are uniform, so taking the supremum over the frozen values afterwards preserves them. All coefficientwise majorant series converge on the same fixed larger products, and finite truncations and the absolute majorants justify the displayed sums.
\end{proof}

\begin{lemma}\label{ch17:lem:decorated-white-log}
Let $\pi=\bigotimes_x\pi_x$ be a finite product of probability measures on fluctuation blocks, each $\pi_x$ supported in the ball $\{\abs{w_x}\le p_x\}$ with $p_x<\infty$; these are the radii used in $b_\lambda$, $W_{16}$, and $D_{16}$. Use the same working rates $a,\alpha>0$ as in $\mathcal X^\phi$, and assume $s_0\ge a+\alpha$. Consider monomial families as in \cref{ch17:def:rooted-norms}, with nonempty supports $S_\lambda$ and networks through all of their positions, and put $\epsilon=W_{16}(U)+D_{16}(U)$. There is an absolute constant $\epsilon_0>0$ such that $\epsilon<\epsilon_0$ implies the following.
\begin{enumerate}
\item The product expansion followed by the hard-core polymer expansion, with incompatibility $S_\lambda\cap S_\mu\ne\varnothing$, gives a normally convergent logarithm of $\int e^{-\sum_\lambda U_\lambda}\,d\pi$, equal to zero at zero interaction, and this integral does not vanish in the coefficient ball. Neither spatial connectedness of $S_\lambda$ nor membership of the background positions in $S_\lambda$ is needed.
\item Every resolved term of the logarithm carries the connected union of the networks of its monomials, with no additional length. Its absolute sum rooted at a fluctuation block, with the weight $e^{4\abs S}$ where $S$ is the union of the supports of its monomials, is at most $C\epsilon$. After passage from $\WideProfile$ to $\phi$, the sum over the background degrees of its rooted coefficient suprema is at most $C\epsilon$.
\item For one marked family $A$ the same bounds hold with $C\epsilon_A$. Let $U(s,r)=U+sB+rA+srC$, $0\le s,r\le1$, where $B$ carries the boundary change, $A$ the incoming insertion, and $C$ both (here $s,r$ are interpolation parameters and $B,C$ are families of monomials). Assume that every occurrence in $B$ and $C$ has an assigned network visiting $\Omega^c_{L_J}$, with that visit included in $\ell_\lambda$ and in its weighted bound, and that every interpolated family lies in the same fixed small ball. Then the mixed derivative $\partial_s\partial_r$ of the logarithm and its four-corner difference are bounded, in both rooted norms after passage of profile, by
\begin{equation}\label{ch17:eq:exact-white-two-mark}
 C(\epsilon_{\partial A}+\epsilon_\partial\epsilon_A),
\end{equation}
where $\epsilon_\partial,\epsilon_A,\epsilon_{\partial A}$ are the values of $W_{16}+D_{16}$ for $B,A,C$. A term containing the boundary change keeps a visit to $\Omega^c$, and a term containing both changes keeps an incoming factor. These bounds keep the weight of the network length. In particular, for a difference $F$ of logarithms containing the boundary change, on the working profile,
\[
 \mathcal X^\phi(F)\le C_J(\epsilon_{\partial A}+\epsilon_\partial\epsilon_A)
\]
in the case of two changes, and $\mathcal X^\phi(F)\le C_J\epsilon_\partial$ in the case of one change.
\item The logarithm and its derivatives in the changes are holomorphic on this ball.
\end{enumerate}
The combinatorial constants are absolute, with the dependence described in item (1) of \cref{ch01:conv:constants}; the constants of the passage of profile, of the density, and of the spatial conversion have the dependence described in item (3).
\end{lemma}

\begin{proof}
The proof follows the standard two-stage cluster expansion; we give the estimates so that the root and the support weights stay visible. We first truncate to finitely many monomials.

\paragraph{Step 1: the product expansion.}
Let $f_\lambda=e^{-U_\lambda}-1$. Its absolute coefficient majorant is at most $e^{b_\lambda}-1\le2b_\lambda$ for small $\epsilon$. Indeed $\abs{w_x}\le p_x$ holds $\pi_x$-almost surely, so the polynomial bounds defining $b_\lambda$ bound the integrands on the support of the product measure. The normalized truncated Gaussian measures on balls used later satisfy this hypothesis. We expand the exact identity
\[
 e^{-\sum_\lambda U_\lambda}=\prod_\lambda(1+f_\lambda).
\]
We group each selected family of distinct indices by the connected components of the graph in which $\lambda,\mu$ are joined when $S_\lambda\cap S_\mu\ne\varnothing$. Different components depend on disjoint sets of fluctuation blocks, so their expectations factor, even if their background parameters coincide. This is an exact hard-core gas whose polymers are unions of fluctuation supports. Spatially disconnected sets may serve as polymers, because compatibility uses only set intersection and the independence of the reference variables.

\paragraph{Step 2: the tree majorant at the first stage.}
Let $n_\lambda=\abs{S_\lambda}$ and $a_\lambda=2b_\lambda e^{10n_\lambda}$. The majorant by rooted trees is the increasing iteration from zero of
\begin{equation}\label{ch17:eq:exact-white-tree}
 T_\lambda=a_\lambda\exp\Bigl(\sum_{\mu:S_\mu\cap S_\lambda\ne\varnothing}T_\mu\Bigr).
\end{equation}
Every connected selected family has a spanning tree; expanding the exponential supplies the factorials for unordered children, and dropping the distinctness of the indices only increases this sum of nonnegative terms. If $T_\mu\le a_\mu e^{n_\mu}$ for all $\mu$, then the exponent in \eqref{ch17:eq:exact-white-tree} is at most
\[
 n_\lambda\sup_x\sum_{\mu:x\in S_\mu}a_\mu e^{n_\mu}\le2\epsilon n_\lambda.
\]
For $2\epsilon\le1$ the candidate $a_\lambda e^{n_\lambda}$ therefore bounds every iterate. Its sum rooted at a fluctuation block is at most $2\epsilon$. Thus the polymer activities are absolutely summable with the weight $e^{10\abs S}$.

\paragraph{Step 3: the second stage.}
We apply the finite graph inequality of \cref{ch09:thm:hard-core-log} to this hard-core gas; only its bound on alternating sums over connected graphs is imported. At the level of polymers we repeat \eqref{ch17:eq:exact-white-tree}, using the weight $e^{4\abs S}$ for the output and a factor $e^{\abs S}$ of the support weight for the children. The estimate with weight $e^{10\abs S}$ bounds the sum over the children by $C\epsilon$ times the size of the parent support. Decreasing $\epsilon_0$ if necessary, we obtain the bound of part (2) with weight $e^{4\abs S}$. Repeated polymer occurrences are allowed, as in the finite graph theorem. Thus the input weight is $e^{16\abs S}$, the polymer activities have weight $e^{10\abs S}$, and the logarithm has weight $e^{4\abs S}$.

\paragraph{Step 4: the networks.}
We keep the resolution into monomials inside each exponential and each polymer. Within a connected family of the first stage, the networks of adjacent factors share a fluctuation position. At the second stage, adjacent polymers share a fluctuation block used by a factor in each of them, so their networks also meet at that position. By induction along a spanning tree at each stage, the union is connected and
\begin{equation}\label{ch17:eq:exact-anchor-union}
 \ell(\mathcal N_{\mathrm{union}})\le\sum_{\lambda\ \mathrm{occurrences}}\ell_\lambda.
\end{equation}
Coincident positions of different color coordinates have distance zero. Copies of the same monomial in an exponential have the same support and network. A visit to $\Omega^c$ remains a visit in the union. Thus the network weights multiply, with no extra connecting length at either stage.

\paragraph{Step 5: the background root.}
Fix a background coordinate $i$ in a coefficient of a product. The identity $\nu_i=\sum_\lambda\nu_{\lambda,i}$ assigns the derivative in $i$ to one monomial occurrence, with its multiplicity. We root the spanning tree of the first stage at this factor and, at the second stage, at the polymer containing it. The estimates for the children just proved bound the remaining part of the expansion by $Ce^{cn_\lambda}b_\lambda$ with $c<16$, so that the total support weight used is less than $e^{16n_\lambda}$. At the distinguished factor we use the degree-$k$ term in $D_{16}$ instead of its norm rooted at a fluctuation block. This works even if the background position is not a fluctuation position. The other factors are summed through their intersections of fluctuation supports.

Fixing $i$ and summing over all output degrees at the intermediate profile, we obtain the bound $C\sum_kd_k\le C\epsilon$, where $d_k$ is the degree-$k$ supremum in $D_{16}$ for the distinguished monomial. Indeed, the remaining part of the expansion is summed at its fixed fluctuation root, and its total coefficient majorant there is bounded by the estimates for the children. We thus obtain a bound on $\sup_i\sum_K$ at the intermediate profile, which is not yet the coefficient norm with boundary decay. Let $q_*=K_*^{-1/2}<1$. At output degree $K$, passing to the working profile multiplies every coefficient by $q_*^K$. Thus, taking the supremum separately at each degree,
\[
 \sum_{K\ge1}\sup_i\bigl\{\text{output coefficients of degree }K\text{ rooted at }i,\text{ at }\phi\bigr\}\le C\epsilon\sum_{K\ge1}q_*^K.
\]
This is the required order of operations; a bound on $\sup_i\sum_K$ at the same radius would not imply a bound on $\sum_K\sup_i$. The second half of the ratio between the profiles is used here, once for each local map on the fixed profiles. The factor $\nu_i/K$ in the coefficient norm with boundary decay only decreases the multiplicity bound. Repeated copies in an exponential obey the same identity for $\nu_i$, and differentiating the exponential accounts for all possible distinguished occurrences.

\paragraph{Step 6: one and two changes.}
The tree estimates also hold in the rooted norms with changes. Differentiating \eqref{ch17:eq:exact-white-tree} once gives
\[
 T_\lambda'=e^{\sum_{\mu\sim\lambda}T_\mu}a_\lambda'+T_\lambda\sum_{\mu\sim\lambda}T_\mu',
\]
where $\mu\sim\lambda$ means $S_\mu\cap S_\lambda\ne\varnothing$. In the norm rooted at a fluctuation block, the last operator has norm at most $C\epsilon<1/2$, since its incidence factor $n_\lambda\le e^{n_\lambda}$ is bounded by a factor $e^{n_\lambda}$ of the support weight that is not used in the output bound. Its inverse is a Neumann series with nonnegative terms and norm at most two. Hence a first variation is bounded by the norm of the changed source. Writing $h_\lambda=\sum_{\mu\sim\lambda}T_\mu$, the mixed variation satisfies
\[
 T_{\lambda,sr}=e^{h_\lambda}a_{\lambda,sr}+e^{h_\lambda}(a_{\lambda,s}h_{\lambda,r}+a_{\lambda,r}h_{\lambda,s})+T_\lambda h_{\lambda,s}h_{\lambda,r}+T_\lambda\sum_{\mu\sim\lambda}T_{\mu,sr}.
\]
Its inhomogeneous terms are a source with both changes, products of the two first variations, and a source with one change times a first variation. Their bounds are $C(\epsilon_{\partial A}+\epsilon_\partial\epsilon_A)$, and the same inverse bounds the mixed variation. The polynomial incidence factors have degree at most two, and $\sup_{n\ge1}n^2e^{-n}<\infty$. In the bilinear interpolation, terms containing the family $C$ are bounded by a constant times $\epsilon_{\partial A}$, since the other monomial norms are bounded on the fixed small ball. Integrating the mixed derivative over the unit square, we obtain the four-corner difference, with the same bound. These calculations apply at both stages of the expansion. For a background root, we first sum the resulting rooted expressions at the distinguished factor at the intermediate profile, with the sequence $d_k$ described above, and then apply the factor $q_*^K$ at each output degree. If the background root lies on a factor different from both changes, the remaining expansion has at most one additional distinguished branch. Its cubic incidence factor is bounded by a further factor $e^{n_\lambda}$ of the support weight, since $\sup_{n\ge1}n^3e^{-n}<\infty$; the passages from $e^{16\abs S}$ to $e^{10\abs S}$ and from $e^{10\abs S}$ to $e^{4\abs S}$ leave these factors available. The source with both changes keeps its incoming factor, and the product source contains the first variation in the incoming insertion. This proves \eqref{ch17:eq:exact-white-two-mark} in both rooted norms, including the case of a background root on an unchanged factor.

\paragraph{Step 7: spatial conversion.}
We keep the joined networks. By the hypothesis on $B$ and $C$, every term containing the boundary change carries a visit to $\Omega^c_{L_J}$. If a coefficient with unfixed background root $i$ carries such a visit, then by part (5) of \cref{ch17:lem:coefficient-resolution}
\[
 a\Tree(\nu)+\alpha d_\Omega(i)\le(a+\alpha)\ell(\mathcal N_{\mathrm{union}})+C_J\alpha.
\]
Since $s_0\ge a+\alpha$, the product of the weights $e^{s_0\ell_\lambda}$ in $b_\lambda$ dominates $e^{(a+\alpha)\ell(\mathcal N_{\mathrm{union}})}$, by \eqref{ch17:eq:exact-anchor-union}. The fixed buffer is used only at the final visit to $\Omega^c$. The factor $\nu_i/K$ in \eqref{ch01:eq:regional-norm} is at most the multiplicity of the distinguished index. The bound at the intermediate profile and the degree sum $\sum_{K\ge1}q_*^K$ just proved, multiplied by the fixed factor $\rho^{-2}$ and the buffer constant, give the stated bounds on $\mathcal X^\phi$.

\paragraph{Step 8: the logarithm.}
The bounds on the coefficientwise majorants prove the normal convergence of every resolved series and of the displayed derivatives. For a small multiple of the interaction, the two series are the Taylor expansion of the logarithm of the finite product integral. Their exponential agrees with the integral throughout the strict ball, by the identity theorem in that multiplier. Hence the integral does not vanish there. The uniform absolute bounds allow us to remove the truncations to finitely many monomials and to differentiate termwise in the parameters, and they prove holomorphy in the absolute coefficient norms. We have not used the hypothesis on background supports in the holomorphy theorem of \cref{ch09:chapter}.
\end{proof}

\begin{theorem}\label{ch17:thm:carried-networks}
Fix the torus rate $a_{\mathbb T}>0$ and the ratio $K_*$ between the profiles $\InflatedProfile=K_*\phi$ and $\WideProfile=\sqrt{K_*}\phi$. There are geometric constants $C_{\mathrm{alloc}}$ and $C_g$, independent of $J$, with the following property. Choose $C_*^{\mathrm{geom}}\ge C_g$ and choose $J$ once with
\[
 B=M^J\ge B_{\mathrm{geom}},\qquad\text{where }B_{\mathrm{geom}}=256C_{\mathrm{alloc}}(C_*^{\mathrm{geom}}+1).
\]
After $J$ has been fixed, choose positive rates
\[
 a=a_{\mathrm{cl}}=\asp,\qquad\alpha=\alpha_{\mathrm{cl}}=a/16,\qquad m_{\mathrm{rate}}=C_*^{\mathrm{geom}}(a+\alpha),\qquad b_1=8C_*^{\mathrm{geom}}\alpha
\]
so that
\[
 m_{\mathrm{rate}}+b_1<\min(c_{\mathrm{prim},J},Ba)/C_{\mathrm{alloc}},\qquad b_1<B\alpha/C_{\mathrm{alloc}},\qquad \sqrt3\,(a+\alpha)<a_{\mathbb T},
\]
where $c_{\mathrm{prim},J}$ is the fixed positive reserve for the intrinsic tensors and the unperturbed weighted reconstruction images, in the graph metric of \cref{ch17:sec:setting}, chosen no larger than the intrinsic exponent $c_J$ of \cref{ch16:thm:controlled-perturbation}. The inherited perturbative networks use the separate inequalities against $Ba$ and $B\alpha$ from \cref{ch16:thm:controlled-perturbation}; their images satisfy the half-radius bound at the allocated working rates.

Consider a full renormalization group step ($s=J$) on admissible data. Sum the internal indices first, keep networks through all positions of the monomials, and use the exact nonempty fluctuation supports of \cref{ch17:def:rooted-norms} in both expansions. All factors and polymer occurrences that are joined share a fluctuation position; the union of their networks keeps every background index and every visit to $\Omega^c$, and adds no length. Set $p=\max_ip_i\ge1$, $x=\mathcal X^\rho(f)$, and $H=E+x$, and substitute the complex parameter $z$ for the coupling powers as described before \cref{ch17:def:rooted-norms}, with the insertion $A_R(z)$. Let $\Xi_0$ and $\Xi_R$ be the integrals of $e^{-U}$ against the product of normalized symmetric reference measures on balls, for the interaction $U$ without insertion and for $U$ plus the insertion, respectively. For the matched torus potential, with the fixed background slots set to zero, define its restriction by deleting the monomials that use a missing fluctuation block, and adjoin the missing blocks with the identical normalized reference measures.

There are $\epsilon_J,c_{z,J}>0$ and $C_J<\infty$ such that the following holds if
\[
 \mathfrak s(z)=\frac{C_J\abs z(p^3+p+Hp)}{1-C_J\abs zp}\le\epsilon_J,\qquad C_J\abs zp<1.
\]
The integrals without and with insertion, the restricted torus integrals, and the linearly interpolated integrals do not vanish, and their logarithms are the normally convergent expansions of \cref{ch17:lem:decorated-white-log}. The regional and restricted families satisfy $W_{16}+D_{16}\le C_J\mathfrak s(z)$; their boundary-difference families satisfy the same bound with networks visiting $\Omega^c$. The comparisons of the nonconstant parts of the logarithms, including the deleted torus terms with an unfixed background index, satisfy
\[
 \mathcal X^\phi\bigl(\log\Xi_\bullet^\Omega-(\log\Xi_\bullet^{\mathbb T})|\bigr)\le C_J\mathfrak s(z),\qquad\bullet\in\{0,R\},
\]
where the bar denotes zero extension of the fixed background slots; constants supported in $\Omega^c$ remain separate. The mixed difference of the boundary change and the insertion is bounded by $C_J(\epsilon_{\partial A}+\epsilon_\partial\epsilon_A)$ in the norms of \cref{ch17:def:rooted-norms}, and every term keeps an incoming factor. The family extends holomorphically to $z=0$, and its compared insertion logarithm $F$ satisfies $F(0)=F'(0)=0$. In particular, if
\[
 r_z=\frac{c_{z,J}}{p^3+p+Hp},\qquad0<t_{\max}\le r_z/2,
\]
then
\[
 \mathcal X^\phi\bigl(\mathsf D^\Omega R-(\mathsf D^{\mathbb T}r_{d_0})|\bigr)\le C_Jt_{\max}^2H(p^4+Hp^2),\qquad\mathsf DR=-\log\Xi_R+\log\Xi_0.
\]
After the whole boundary layer with reduced radius has been frozen as in \cref{ch15:lem:freezing}, the classical and quantum comparisons return the same rates $a,\alpha$ and the same input radii; the torus rate remains $a_{\mathbb T}$. The constants have the dependence described in item (3) of \cref{ch01:conv:constants} and are independent of the later widths. The return of the same profile is asserted only for full steps with $B\ge B_{\mathrm{geom}}$, not for compositions of fewer than $J$ averaging steps.
\end{theorem}

\begin{proof}
The constants $C_{\mathrm{alloc}}$ and $C_g$ bound only the fixed contraction with every index rooted and the transport of the input boundary weight in \cref{ch16:thm:controlled-perturbation,ch15:lem:wall-networks}; they include the factor $\sqrt3$ of the conversion between Euclidean and graph exponents in \cref{ch17:sec:setting}, and hence do not depend on $J$. The torus coefficients, bounded in $\mathcal B_{\rho,a_{\mathbb T}}$ with Euclidean network lengths, have graph exponent $a_{\mathbb T}/\sqrt3>a+\alpha$ by the last inequality in the statement. With $\alpha=a/16$ we have $m_{\mathrm{rate}}+b_1=(25/16)C_*^{\mathrm{geom}}a$, and both restrictions against $Ba$ and $B\alpha$ follow from the lower bound on $B$. At this fixed $J$, we decrease $a>0$ to satisfy the bounds on the exponents of the coefficients of the potential, of the weighted images, and of the torus, including $\sqrt3\,(a+\alpha)<a_{\mathbb T}$. No constant depending on $J$ is used to choose $J$. The comparison of the classical boundary networks uses the same allocation of the decay rate and no expansion with respect to a reference product. At zero fluctuation it works as follows. Consider a classical comparison, or the comparison of the pullback terms, with the boundary tensors rooted at every index. At background degree $k\ge1$, its rooted absolute sum on $\InflatedProfile$ is bounded by its tensor majorant for fixed $J$, and that allocation already accounts for the weights of its output network and of its distance to $\Omega^c$. Passing to $\phi$ gives the factor $K_*^{-k}$. Summing over the degree, including any fixed polynomial factor for the choice of a distinguished index, we obtain the bound in the coefficient norm with boundary decay. For a pulled-back incoming boundary term, we use its jet majorant at half the radius and its distinguished input columns before this calculation. Thus the classical coefficients are converted without assigning a fluctuation support to a term at zero fluctuation, and these terms never enter either expansion.

We apply \cref{ch17:lem:decorated-white-primitives} at these rates and then \cref{ch17:lem:decorated-white-log}. The hypotheses of the latter on the rates and on the boundary-change families hold, since \cref{ch17:lem:decorated-white-primitives} uses $s_0=a+\alpha$ and gives the boundary differences with their networks visiting $\Omega^c$. Restricting a torus family decreases its norms, which are sums of absolute values. For its difference from the full family, every deleted coefficient with an unfixed background root visits a missing fluctuation position and hence $\Omega^c$, and the same lemma applies. Pure background constants without an unfixed index have zero norm with boundary decay and stay outside this comparison. The estimate with two changes implies the comparison of the logarithms with and without insertion. The logarithms and their derivatives in the parameters are determined by the finite identities and normal convergence proved above, not by a choice of a complex logarithm after integration.

We choose $c_{z,J}$ so that the disc $\abs z\le r_z$ lies strictly inside the margins of \cref{ch17:lem:decorated-white-primitives}, of the free product, and of the logarithm. On this disc, \eqref{ch17:eq:decorated-inherited-gate} and the bounds with one and two changes give
\[
 \norm{F(z)}_{\mathcal X^\phi}\le C_J\abs zHp.
\]
At $z=0$ every fluctuating monomial vanishes. The coefficient of $z$ is cubic in the fluctuation variables for the classical terms and linear for the quantum, Haar density, determinant, and incoming terms. Each of these is odd under the simultaneous reversal $w\mapsto-w$, so its expectation under the symmetric reference measures on balls vanishes. At first order the logarithm contains a single monomial occurrence, hence $F(0)=F'(0)=0$ coefficientwise, at every real frozen value.

The normally convergent series $F(z)/z$ is holomorphic, vanishes at zero, and has norm at most $C_JHp$. Cauchy's coefficient bound on $\abs z=r<r_z$, followed by $r\uparrow r_z$, gives for $t\le r_z/2$
\[
 \norm{F(t)}_{\mathcal X^\phi}\le2C_Jt^2Hp/r_z\le C_J't^2H(p^4+Hp^2),
\]
where in the last step we used $p\ge1$. The integral proof of the coefficient bound applies to the supremum over the fixed index at every background degree, and then to their sum, by the triangle inequality.

Finally, \cref{ch15:lem:freezing} restores the incoming radii when the whole zone with reduced radius is frozen with the stated bounds on the frozen values. All conversions of coefficients were made on the fixed outer and intermediate profiles, and none is iterated at later steps. This proves the return of the same profile and the stated independence of the constants.
\end{proof}

\section{Effective interactions and normalization constants}\label{ch17:sec:normalization}

\begin{theorem}\label{ch17:thm:quantum-ledger}
Consider a full renormalization group step, $s=J$, on admissible data as in \cref{ch17:thm:carried-networks}. Set
\[
 \mathsf P^\Omega R(v)=R\bigl(\Log C^\Omega(v,0;\Rec)\bigr),\qquad
 \mathsf Q^\Omega=B_{\mathrm{ref}}^\Omega-\sum_if_i\log z_{p_i}-\log\Xi_0^\Omega,
\]
\[
 B_{\mathrm{ref}}=\sum_iB_i(\widehat c_{i+1}),
\]
where $f_i$ is the number of free blocks at scale $i$ and $z_p=\int_{\abs w\le p}e^{-\abs w^2/2}\,d^3w$, in the convention of \cref{ch09:thm:sequential-density}. The field-independent factors of the block averaging constraint and of the Haar measure extracted from $L_i$ are kept in $\ScalarLedger$. Up to the exact constant kept in $\ScalarLedger$, the new effective interaction $\mathsf M^\Omega R$ produced by the step satisfies
\[
 \mathsf M^\Omega R=\mathsf Q^\Omega+\mathsf P^\Omega R+\mathsf D^\Omega R.
\]
Here $\mathsf Q^\Omega$ is the newly generated term, $\mathsf P^\Omega R$ is the pullback of the incoming interaction by the zero-fluctuation reconstruction, and $\mathsf D^\Omega R$ is the fluctuation correction of \cref{ch17:thm:carried-networks}. For the matched torus reference with zero extension of the fixed slots, there are constants $A_J$, $K_J$, and $C_J$, independent of the later widths, such that
\begin{align*}
 \mathcal X^\phi\bigl(\mathsf Q^\Omega-\mathsf Q^{\mathbb T}|_{\Omega'}+\mathsf P^\Omega(r_{d_0}|_\Omega)-(\mathsf P^{\mathbb T}r_{d_0})|_{\Omega'}\bigr)&\le A_J(1+E),\\
 \mathcal X^\phi(\mathsf P^\Omega f)&\le K_J\mathcal X^\rho(f).
\end{align*}
Here $\Omega'$ is the set of unfixed scalar output labels. Assume that the column bounds of \cref{ch16:thm:hybrid-columns} and their perturbed versions hold on the product with doubled output analytic radii, and that the image under the reconstruction majorant on this product is at most half the incoming radius, as in \cref{ch16:lem:strict-two-radius-images}. Then one may take $K_J=4(\log2)K_{\mathrm{col},J}$. Here $K_{\mathrm{col},J}$ bounds, for one fixed output, the sum over all old input labels, weighted by their boundary decay, of the distinguished derivative series of the perturbed zero-fluctuation reconstruction; it is chosen uniformly on the preliminary perturbation ball by \cref{ch16:thm:controlled-perturbation}. The degree-one terms and every nonconstant dependence of an omitted torus density on the unfixed labels are included.

At every admissible complex value of the coarse background and every fixed real admissible frozen value, let $I_i$ and $H_i$ be the sets of unfixed and of missing scalar free labels at scale $i$. With matched torus arguments extended by zero,
\[
 \log\det K_i^\Omega=\sum_{x\in I_i}(\log K_i^{\mathbb T})_{xx}+\mathcal B_i,\qquad\abs{(\log K_i^{\mathbb T})_{xx}}\le C_J,\qquad\abs{\mathcal B_i}\le C_J\abs{\partial_{\mathrm{blk}}\Omega}.
\]
These determinant statements also hold in the coefficient norm before evaluation. Scalar determinant labels and fluctuation blocks are counted separately. All constants have the dependence described in item (3) of \cref{ch01:conv:constants} and are independent of the later widths.
\end{theorem}

\begin{proof}
Every step here is a full step, $s=J$; a final composition of fewer than $J$ averaging steps that is evaluated deterministically adds no fluctuation integral.

\paragraph{Step 1: the exact factorization.}
The triangular calculation in the proof of \cref{ch17:lem:complete-white-potential} gives, at a fixed coarse background and fixed frozen values, the value of the integral
\begin{equation}\label{ch17:eq:auxiliary-factorization}
 \Bigl(\prod_i\mathcal J_{i,\mathrm{ref}}g_i^{3f_i}\Bigr)e^{-F_{\mathrm{out}}/g_s^2-B_{\mathrm{ref}}-R(\Log C(v,0;\Rec))}\prod_iz_{p_i}^{f_i}\,\Xi_R(v),
\end{equation}
where $F_{\mathrm{out}}/g_s^2$ is the classical value of the step with the output coupling, and the $\mathcal J_{i,\mathrm{ref}}$ are the extracted constant Haar disintegration factors. Inserting the definitions of $\mathsf Q$, $\mathsf P$, and $\mathsf D$ into the exponent of \eqref{ch17:eq:auxiliary-factorization}, and writing $\log\Xi_R=\log\Xi_0-\mathsf DR$, we obtain $\mathsf MR=\mathsf Q+\mathsf PR+\mathsf DR$. This proves the recursion, including its signs and its Gaussian convention. Centering moves the value at zero into $\ScalarLedger$ and changes no factor of this identity.

\paragraph{Step 2: the generated term and the torus pullback.}
The deterministic part of $\mathsf Q^\Omega-\mathsf Q^{\mathbb T}|$ is the difference of the Haar densities at the minimizers and of the diagonal matrix-logarithm densities along the two zero-fluctuation chains. Their rooted comparisons with boundary decay are part (3) of \cref{ch17:thm:perturbed-hessian-coarea}. A density based at a missing torus label is not discarded: a nonconstant coefficient with an unfixed index has a network to that missing label, and its sum over the missing label is bounded from the unfixed index by \eqref{ch17:eq:density-sum}. Its constant is left in $\ScalarLedger$. The terms $\log z_{p_i}$ and the constant Haar factors have no nonconstant background coefficient. The difference of the logarithms without insertion has the same fixed rates by \cref{ch17:thm:carried-networks}, with a bounded norm in the cluster expansion. We obtain a finite bound $C_J$ for the comparison of the nonconstant generated terms.

For the torus pullback, denote by $a_\Omega$ the zero extension of $\Log C^\Omega(v,0;\Rec)$ and by $a_{\mathbb T}$ the matched torus argument; in this paragraph $a_{\mathbb T}$ is an argument, not the torus rate. We have
\begin{equation}\label{ch17:eq:principal-torus-difference}
 r_{d_0}(a_\Omega)-r_{d_0}(a_{\mathbb T})=\int_0^1Dr_{d_0}\bigl(a_{\mathbb T}+\tau(a_\Omega-a_{\mathbb T})\bigr)[a_\Omega-a_{\mathbb T}]\,d\tau.
\end{equation}
Both absolute images are below half the incoming radius, so the whole segment admits the coefficient expansion of the integrand. On an unfixed old input, the difference of the arguments is the unperturbed defect of the reconstruction followed by the gauge transformation (\cref{ch16:thm:hybrid-columns}) plus the defect of the perturbation (\cref{ch16:thm:controlled-perturbation}). On a fixed old input, the difference is a visit to $\Omega^c$. The action and interaction tensors join every output index to that visit, even if the source does not decay inside its own component of $\Omega^c$. The jet bound \eqref{ch17:eq:residual-jets} with $\sigma_0=0$ bounds the integrand by $C_JE$, with its visit. The conversion of networks and degrees in the proof of \cref{ch17:thm:carried-networks} therefore bounds its $\mathcal X^\phi$ norm by $C_JE$. Together these bounds prove the first inequality, with $A_J$ fixed on the larger products. No first derivative of the output was canceled.

\paragraph{Step 3: the pullback of the boundary function.}
For the zero-fluctuation reconstruction $\Phi_e(v)=\Log C_e(v,0;\Rec)$, we define
\[
 C_{e,x}(2\phi)=\sum_\nu\nu_x\abs{\Phi_{e,\nu}}(2\phi)^\nu e^{a_1\Tree(\{e\}\cup\supp\nu)}.
\]
All constants in this step refer to the perturbed map. Let $D_J^0$ bound the unperturbed distinguished series on the larger product with doubled output analytic radius. The difference of the reconstructions in \cref{ch16:thm:controlled-perturbation}, rooted at every index, with the uniform incoming norm and the inequalities on the conditional forces checked above, bounds the additional distinguished series by $\sigma_*D_J^1$ there. We fix a spatial decay rate $b>\alpha$ and define
\begin{equation}\label{ch17:eq:actual-column-ceiling}
 D_J^{\mathrm{act}}=D_J^0+\sigma_0D_J^1,\qquad C_{e,x}(2\phi)\le D_J^{\mathrm{act}}(2\phi_x)e^{-b\dist(e,x)}.
\end{equation}
The factor $2\phi_x$ appears because the input index is distinguished before the sum of absolute coefficients is taken. The rates are below the common exponent available in the coefficient bounds of the potential, and both $D_J^0$ and $D_J^1$ are computed at the preliminary bound on fixed larger products, before the later widths. Here and below, for an old input label $e$ of $R$ (a label of the incoming set $D_{\mathrm{in}}$, on the input lattice of the step), we write $d_{\mathrm{old}}(e)=d_{D_{\mathrm{in}}}(e)$ for its distance to the incoming fixed complement, measured in units of the input lattice as in the incoming norm. For a current unfixed background input $x$, we write $d_{\mathrm{out}}(x)=d_\Omega(x)$ for its distance to $\Omega^c$, in units of the output lattice. Finally, $\dist(e,x)$ is the distance between $x$ and the position $e/B_s$ of $e$, in units of the output lattice. Applying \eqref{ch17:eq:incoming-wall-gate} on the old input set of $R$, which is the same set of unfixed labels, with quantum radius $\rho$, and dividing by $B_s$, we obtain
\[
 d_{\mathrm{out}}(x)\le d_{\mathrm{old}}(e)/B_s+\dist(e,x)+C_{\mathrm{buf},J}/B_s.
\]
In \cref{ch16:thm:hybrid-columns} the distance $d_{\mathrm{old}}$ is measured in units of the output lattice, and the inequality $d_{\mathrm{out}}(x)\le d_{\mathrm{old}}(e)+\dist(e,x)+C_J$ is an explicit hypothesis of that theorem. With $d_{\mathrm{old}}$ there equal to our $d_{\mathrm{old}}(e)/B_s$, the inequality just displayed verifies this hypothesis, with buffer $C_{\mathrm{buf},J}/B_s$; it holds because the incoming fixed complement lies on the current fixed side. Since $B_s\ge1$, this also gives the weaker comparison with $d_{\mathrm{old}}(e)$ in place of $d_{\mathrm{old}}(e)/B_s$. The density of the old labels then implies
\begin{equation}\label{ch17:eq:pullback-column}
 \sum_{e\ \mathrm{old\ unfixed}}e^{-\alpha d_{\mathrm{old}}(e)}C_{e,x}(2\phi)\le K_{\mathrm{col},J}\,\rho\,e^{-\alpha d_{\mathrm{out}}(x)}.
\end{equation}
We define $K_{\mathrm{col},J}$ as a common bound for this inequality, uniformly on the preliminary perturbation ball. One admissible choice is any number at least $2D_J^{\mathrm{act}}e^{\alpha C_{\mathrm{buf},J}/B_s}S_J(b-\alpha)$, with the fixed volume factors included in $S_J$, and at least the unperturbed column bound; we take the finite maximum over the label types and the stages. The absolute image on the product with doubled output analytic radius is at most $\rho/2$. The construction used a ratio larger than two between the profiles, so no repeated reduction of radii is involved. This is the column calculation of \cref{ch16:thm:hybrid-columns}, with its constant enlarged by \eqref{ch17:eq:actual-column-ceiling}. The hypotheses on the doubled product and on the incoming image at half the radius are used here.

Let $x_f=\mathcal X^\rho(f)$. At an old input root $e$, the distinguished derivative sum evaluated at $\rho/2$ satisfies
\begin{equation}\label{ch17:eq:half-derivative}
 \sum_{\eta:\eta_e>0}\eta_e\abs{f_\eta}(\rho/2)^{\abs\eta-1}e^{a\Tree(\supp\eta)}\le4\rho x_fe^{-\alpha d_{\mathrm{old}}(e)}.
\end{equation}
For degree $l$ this follows by multiplying the normalized coefficient with this input fixed by $l2^{1-l}\rho$, which is at most $4\rho$; the case $l=1$ is included. The chain rule for coefficientwise majorants, the union of networks, and \eqref{ch17:eq:pullback-column} bound the distinguished coefficient sum of the output at $2\phi$ by $4K_{\mathrm{col},J}\rho^2x_fe^{-\alpha d_{\mathrm{out}}(x)}$. At degree $k$ the fixed-index weight of the coefficient norm with boundary decay inserts the factor $1/k$, and evaluation at $\phi$ instead of $2\phi$ inserts $2^{-k}$. We take the supremum over the fixed index at each degree and then sum. Since integrating $\sum_{j\ge0}u^j=(1-u)^{-1}$ over $[0,1/2]$ gives $\sum_{k\ge1}2^{-k}/k=\log2$, we obtain
\[
 \mathcal X^\phi(\mathsf P^\Omega f)\le4(\log2)K_{\mathrm{col},J}x_f.
\]
This also shows that a finite $K_J$ exists when no particular column constant is specified.

\paragraph{Step 4: the logarithmic determinants.}
The diagonal coefficient series of $\log K_i^{\mathbb T}$ is bounded uniformly for each scalar label. To bound the constant as well as the derivatives, we write
\[
 \log K=\int_0^\infty\bigl[(1+t)^{-1}I-(K+t)^{-1}\bigr]\,dt.
\]
Near $t=0$ the gap bounds the integrand. For large $t$, the resolvent identity with $I$ bounds it by $C_J(1+t)^{-2}$, since $K-I$ has a bounded coefficient row on the larger product. This integrable bound and the derivative bounds with every index rooted allow us to integrate coefficient by coefficient. The exact normalization of the scales adds only the same term $i\log M$ for each scalar label, with $i<J$. Thus $\abs{(\log K_i^{\mathbb T})_{xx}}\le C_J$ on the stated complex profile, not only at zero background.

The comparison of a regional logarithm with the restriction of the torus logarithm is the sum of \eqref{ch17:eq:log-difference} and \eqref{ch17:eq:restricted-log}. Its diagonal density at an unfixed root $u$ is bounded in coefficient norm by $C_Je^{-c_Jd_\Omega(u)}$. Summing these diagonal defects and using \eqref{ch17:eq:boundary-sum}, we obtain the bound on $\mathcal B_i$. If $\Omega^c$ is empty, the defect is zero.
\end{proof}

To state the recursion for $\Lambda$ we fix the counts that enter it.

\begin{definition}\label{ch17:def:counts}
Consider a full step on the block region $\Omega$ and the matched torus step. Every scalar label is counted on the lattice on which it is introduced.
\begin{enumerate}
\item For each scale $i$, $\abs{H_i}$ is the number of missing scalar free labels at scale $i$, that is, free labels of the torus step that are not labels of the regional step.
\item $N_{\mathrm{missing}}$ is the number of scalar labels of the torus step (background labels on the output lattice and free labels at the intermediate scales) that are not labels of the regional step. We write $\mathcal M_{\mathrm{block}}$ for the set of missing torus fluctuation blocks; since each such block consists of three missing scalar labels, $\abs{\mathcal M_{\mathrm{block}}}\le N_{\mathrm{missing}}$.
\item $N_{\mathrm{bd}}$ is the number of scalar labels of the regional step whose positions lie within the fixed distance $w_{\mathrm{band}}\ge1$ of $W$, where $w_{\mathrm{band}}$ is the width of the interface band determined by the construction stencils; it does not depend on the later widths. This is $N_{\mathrm{band}}$ of \eqref{ch17:eq:block-boundary-count}, so $\abs{\partial_{\mathrm{blk}}\Omega}\le N_{\mathrm{bd}}\le C_J\abs{\partial_{\mathrm{blk}}\Omega}$, including for isolated final blocks. Both counts are zero when $W$ is empty.
\item $N_{\mathrm{freeze}}$ is the number of analytic scalar labels that are newly frozen at the end of the step, including, at the first step, the labels frozen at initialization.
\end{enumerate}
When the steps are indexed by $t$, these counts are written $\abs{H_i}$ for the scales $i$ of step $t$, and $N_{\mathrm{missing},t}$, $N_{\mathrm{bd},t}$, $N_{\mathrm{freeze},t}$.
\end{definition}

\begin{lemma}\label{ch17:lem:real-cut-integral}
Assume the hypotheses on the products and the small-interaction hypothesis of \cref{ch17:thm:carried-networks}, in a full step, and fix a real admissible coarse background, a real admissible frozen value, and a real coupling $t=t_{\max}$. Use the potential without or with the insertion, from \cref{ch17:lem:complete-white-potential}, and remove its entire zero-fluctuation part before forming the interaction $U_\bullet$, $\bullet\in\{0,R\}$. Let $\pi$ be the product of the identical normalized three-dimensional Gaussian reference measures on balls, including the reference measures of all missing torus blocks, and write
\[
 \Xi_\bullet^{\mathbb T}=\int e^{-U_\bullet^{\mathbb T}}\,d\pi,\qquad\Xi_\bullet^\Omega=\int e^{-U_\bullet^\Omega}\,d\pi,\qquad\bullet\in\{0,R\}.
\]
The regional interaction does not depend on the adjoined variables. With $N_{\mathrm{missing}}$ and $N_{\mathrm{bd}}$ as in \cref{ch17:def:counts}, there is a restricted interaction $U_\bullet^c$ on the same product such that
\begin{align*}
 \norm{U_\bullet^{\mathbb T}-U_\bullet^c}_\infty&\le C_J\mathfrak s(t)(N_{\mathrm{missing}}+N_{\mathrm{bd}}),\\
 \norm{U_\bullet^\Omega-U_\bullet^c}_\infty&\le C_J\mathfrak s(t)N_{\mathrm{bd}}.
\end{align*}
Consequently, for the positive real logarithms,
\begin{equation}\label{ch17:eq:real-vacuum-comparison}
 \abs{\log\Xi_\bullet^{\mathbb T}-\log\Xi_\bullet^\Omega}\le C_J\mathfrak s(t)(N_{\mathrm{missing}}+N_{\mathrm{bd}}).
\end{equation}
This includes evaluation at coarse background zero, with the frozen values unchanged. Empty products have integral one.
\end{lemma}

\begin{proof}
The reference measure on a block $(i,e)$ is $z_{p_i}^{-1}\one_{\abs w\le p_i}e^{-\abs w^2/2}\,d^3w$. We take the product of these measures over the union of the two sets of blocks. Adjoining blocks multiplies the regional integral by one, even for blocks at different scales with different radii.

For this comparison of constants only, we resolve the torus monomials as in \cref{ch17:def:rooted-norms}, and let $\mathcal D_{\mathrm{mon}}=\{\lambda:S_\lambda\cap\mathcal M_{\mathrm{block}}\ne\varnothing\}$ be the set of deleted monomials. We call the remaining interaction $U^c$. The bound on $W_{16}$ gives
\[
 \norm{U^{\mathbb T}-U^c}_\infty\le\sum_{x\in\mathcal M_{\mathrm{block}}}\sum_{\lambda:x\in S_\lambda}b_\lambda\le C_J\mathfrak s\,\abs{\mathcal M_{\mathrm{block}}}.
\]
The set of monomials may be infinite; no bound on its cardinality is used. By \cref{ch17:def:counts}, $\abs{\mathcal M_{\mathrm{block}}}\le N_{\mathrm{missing}}\le N_{\mathrm{missing}}+N_{\mathrm{bd}}$. This proves the first inequality.

For the regional and the restricted monomials, we compare coefficients on the common set of unfixed fluctuation labels, at the same background argument extended by zero. We use the comparison before the logarithm in \cref{ch17:thm:perturbed-hessian-coarea,ch17:lem:complete-white-potential}: a reconstruction formula, a minimizer, an inverse square root, or a Haar density is replaced one factor at a time. A change of the coarser arguments is the integral of their derivative against the difference of the reconstructions, which carries a visit to $\Omega^c$. For an incoming torus term we integrate $Dr_{d_0}$ on the segment between the two absolute images, which lie below half the radius; fixed old input slots are visits to $\Omega^c$, and unfixed slots use the defect of the reconstruction in \cref{ch16:thm:controlled-perturbation}. The boundary function $f$ enters with its own incoming boundary weight. By the versions with every index rooted, these comparisons keep a visit to $\Omega^c$ and every specified unfixed fluctuation index, including all degree-one terms. They are estimates on monomials before either logarithm is expanded, so the present restriction need not coincide with the one used in the earlier expansions.

Every remaining fluctuation monomial has an unfixed fluctuation index. We assign each compared monomial to its first such index $y$. The comparison of coefficientwise majorants with every index rooted bounds its absolute sum there by $C_J\mathfrak s(t)e^{-\alpha d_\Omega(y)}$. Summing over these indices, we obtain
\[
 \norm{U^\Omega-U^c}_\infty\le C_J\mathfrak s(t)\sum_{y\ \mathrm{unfixed}}e^{-\alpha d_\Omega(y)}\le C_J'\mathfrak s(t)\abs{\partial_{\mathrm{blk}}\Omega}\le C_J'\mathfrak s(t)N_{\mathrm{bd}},
\]
by \eqref{ch17:eq:boundary-sum} and \cref{ch17:def:counts}. This proves the second inequality, at every fixed real frozen value.

The zero-fluctuation values of $F_i$, $L_i$, $B_i$, and $R$ in \cref{ch17:lem:complete-white-potential} were subtracted before this sum and stay in the prefactor. In particular, a constant depending on the frozen values does not enter the boundary estimate. The convergent coefficientwise majorants justify all assignments and sums, by finite truncation and monotone convergence of absolute sums.

Finally, if real bounded $U,V$ satisfy $\abs{U-V}\le D$ pointwise, then $e^{-D}e^{-V}\le e^{-U}\le e^De^{-V}$. Integrating against a probability measure and taking logarithms gives
\[
 \Bigl\lvert\log\int e^{-U}\,d\pi-\log\int e^{-V}\,d\pi\Bigr\rvert\le D.
\]
All integrals here are finite and strictly positive because the potentials have bounded majorants. We apply this twice, through $U^c$, which proves \eqref{ch17:eq:real-vacuum-comparison}. For a real interaction, the integral with the interaction multiplied by $0\le\theta\le1$ is positive and equals one at $\theta=0$. The logarithm given by the cluster expansion vanishes at $\theta=0$ and therefore agrees with the positive real logarithm along this path. The comparison therefore also bounds the constant used in the analytic recursion.
\end{proof}

\begin{theorem}\label{ch17:thm:real-scalar-ledger}
Use full steps on admissible data, the exact extraction of field-independent factors, and the centered torus recursion of \cref{ch17:thm:quantum-ledger}. Restrict the coarse backgrounds, the couplings, and the frozen values to real data on the chosen products of fluctuation balls. Index the full steps by $t=0,1,\ldots$, keep the exact incoming value of $\ScalarLedger$, and write $\Delta\ScalarLedger_t=\ScalarLedger_t-\ScalarLedger_t^{\mathbb T}$ for the difference between the regional and the torus normalization constants before step $t$. At the end of step $t$, the newly frozen analytic scalar labels form a set $F$, frozen simultaneously at the values $c^\sharp$. Assume the following.
\begin{enumerate}
\item[(a)] The outgoing centered effective interaction of step $t$ has the form $R^{\mathrm{out}}=r_{\mathbb T}+f$ with $R^{\mathrm{out}}(0)=0$, $Dr_{\mathbb T}(0)=0$, $\Anorm{\rho}{a}{r_{\mathbb T}}\le E$, and $\mathcal X^\phi_\Omega(f)\le Y$, where $Y$ is a fixed bound.
\item[(b)] Every newly frozen value satisfies $\abs{c^\sharp_j}\le\phi(j)/4$ for $j\in F$.
\end{enumerate}
Then at each step, including the simultaneous restriction to fixed values at its end,
\begin{equation}\label{ch17:eq:real-ledger-increment}
 \Delta\ScalarLedger_{t+1}\le\Delta\ScalarLedger_t+\sum_{i\ \text{in step }t}\abs{H_i}\log(1/g_i)+C_J(N_{\mathrm{missing},t}+N_{\mathrm{bd},t}+N_{\mathrm{freeze},t}),
\end{equation}
with the counts of \cref{ch17:def:counts}, where $C_J$ may also depend on $E$ and $Y$. Each Gaussian dimension occurs once, and the normalized auxiliary variables and the Haar measure of the block axial gauge trees have mass one. The restriction to fixed values is exactly
\[
 R_{\mathrm{new}}(v)=R^{\mathrm{out}}(v,c^\sharp)-R^{\mathrm{out}}(0,c^\sharp),\qquad\ScalarLedger_{\mathrm{new}}=\ScalarLedger^{\mathrm{out}}-R^{\mathrm{out}}(0,c^\sharp),
\]
and the frozen constant satisfies
\begin{equation}\label{ch17:eq:freezing-charge}
 \abs{R^{\mathrm{out}}(0,c^\sharp)}\le\rho^2\abs F(E/16+Y/4).
\end{equation}
The classical constants remain in the full classical action. A constant of the incoming interaction is first moved into $\ScalarLedger$ exactly, using $e^{\ScalarLedger-R}=e^{\ScalarLedger-R(0)}e^{-(R-R(0))}$; no nonconstant norm bounds this constant. No additional term proportional to $\abs\Omega\log(1/g)$ appears. The restriction at initialization is included, with $R^{\mathrm{out}}$ the torus reference and $f=0$. The constant $C_J$ has the dependence described in item (3) of \cref{ch01:conv:constants} and is independent of the later widths, of the preceding depth, and of the number of steps. No inequality for complex values and no recursion for compositions of fewer than $J$ averaging steps is asserted.
\end{theorem}

\begin{proof}
We compare the logarithms of the regional and the torus factorizations term by term: first the deterministic factors (Step 1), then the fluctuation integrals (Step 2), then the pullback at zero fluctuation (Step 3), and finally the constant produced by fixing coordinates (Step 4).

\paragraph{Step 1: coupling powers, determinants, and Haar factors.}
At fixed real frozen values, the factorization \eqref{ch17:eq:auxiliary-factorization} contains exactly one factor $g_i$ for each scalar free coordinate and one triangular determinant. Subtracting the logarithms of the regional and the torus factorizations, we find for the coupling powers
\[
 \sum_i\Bigl[\sum_{x\in I_i}\log g_i-\sum_{x\in I_i\cup H_i}\log g_i\Bigr]=\sum_i\abs{H_i}\log(1/g_i).
\]
The determinant statement of \cref{ch17:thm:quantum-ledger} bounds the difference of the logarithmic determinants by $C_J$ for each missing scalar label plus $C_J\abs{\partial_{\mathrm{blk}}\Omega}\le C_JN_{\mathrm{bd}}$, by \eqref{ch17:eq:block-boundary-count}. The same argument applies to the Haar densities at zero fluctuation, by the comparison in part (3) of \cref{ch17:thm:perturbed-hessian-coarea}. Their constant normalizations depend only on $M$ and $J$ and contribute $C_J$ for each omitted or interface coordinate. Since $p_i\ge1$ we have $0<z_1\le z_{p_i}\le(2\pi)^{3/2}$, so the terms $\log z_{p_i}$ contribute a bounded amount for each missing block. The Haar measure of the block axial gauge trees and the adjoined reference measures contribute no dimension.

\paragraph{Step 2: the fluctuation integrals.}
At coarse background zero, \cref{ch17:lem:real-cut-integral} bounds the difference of the fluctuation logarithms without insertion, and of those with insertion whenever the factorization uses them, by $C_J\mathfrak s(N_{\mathrm{missing}}+N_{\mathrm{bd}})$. In that proof, each missing fluctuation block is counted among the missing labels and each comparison at unfixed roots in the interface count; no geometric crossing argument is needed. Using $\mathfrak s\le\epsilon_J$, we find that this step contributes $C_J(N_{\mathrm{missing}}+N_{\mathrm{bd}})$.

\paragraph{Step 3: the pullback at zero fluctuation.}
The zero-fluctuation value of the pullback has not entered Step 2. For its centered torus and boundary parts we assign each monomial of degree $l\ge1$ to its input roots with the shares $\eta_e/l$, whose sum is one. The input norms of \cref{ch01:def:analytic-norms} give $C_JE$ or $C_J\sigma_*$ for each input root, after the common image at half the radius has summed the degree series. Missing roots are counted as missing labels. Differences rooted at unfixed labels carry the visit in \eqref{ch17:eq:principal-torus-difference} and are summed by \eqref{ch17:eq:boundary-sum}. The incoming constant was moved into $\ScalarLedger$ exactly before this calculation.

\paragraph{Step 4: fixing coordinates.}
Finally we apply \cref{ch15:lem:freezing}. To prove \eqref{ch17:eq:freezing-charge}, note that when the unfixed variables are set to zero, only monomials supported in $F$ survive. On such a monomial of degree $l$ we have $\sum_{j\in F}\nu_j/l=1$, and by hypothesis (b) its frozen value has modulus at most $4^{-l}\abs{f_\nu}\phi^\nu$. Summing over $j\in F$ and using the definition of $\mathcal X^\phi_\Omega$ with the weights larger than one dropped, we find that the boundary part contributes at most $\abs F\rho^2Y/4$, since $l\ge1$. The torus part has $l\ge2$ and $\abs{c^\sharp_j}\le\rho/4$, so the seminorm bounds its contribution by $\abs F\rho^2E/16$. This proves \eqref{ch17:eq:freezing-charge}; the bound contributes $C_J$ for each newly frozen analytic label, degree one included. The value is subtracted from $\ScalarLedger$ with the sign in the statement, and its absolute value gives the upper bound. The same identity applies at initialization. Data deep inside the fixed complement stay in the external factors, and the classical constants stay in the classical action. Adding the finitely many contributions proves \eqref{ch17:eq:real-ledger-increment}. Every estimate used a coefficient bound for each label or the boundary sum \eqref{ch17:eq:boundary-sum}, which does not depend on the shape of the region; none of them estimates the incoming constant or introduces a factor depending on the number of components or of steps.
\end{proof}

\section{Preservation of the analytic bounds under iteration}\label{ch17:sec:iteration}

The following theorem is the main result of this section.

\begin{theorem}\label{ch17:thm:controlled-kernel}
Fix the torus analytic radii and rates, the torus support weights and spatial margins, and $M$, and fix $A_{\mathrm w}\ge1$. Choose $J$ once so that the conditions on $J$ of the torus construction used in \cref{ch16:lem:strict-two-radius-images} hold and $B\ge B_{\mathrm{geom}}$, as in \cref{ch17:thm:carried-networks}. There exist positive fixed radii $\rho$, $\rho_{\mathrm{cl}}$, $\Rraw$, $\Rwhite$, and $\Ren$, rates $a,\alpha$ as in \cref{ch17:thm:carried-networks}, a perturbation bound $\sigma_*$, the classical and quantum radius profiles defined in \cref{ch17:sec:setting}, a width $s_{\mathrm{er}}$ of the frozen layer, a bound $\rrecord$ on the frozen values, a support radius $r$, and a coupling bound $g_*<e^{-1}$, with the following property. One may take $\rho=\rho_{\mathrm{cl}}$; other fixed ratios are allowed if both image conditions are imposed. The bound $g_*$ may depend on $A_{\mathrm w}$.

Consider a full step on an admissible block region satisfying the protected-neighborhood conditions of \cref{ch16:lem:strict-hybrid-gauge}, at any preceding depth and after any number of earlier steps, with a real admissible frozen value and an input satisfying
\[
 \mathcal X^{\mathrm{cl}}(\bar P)\le2\sigma_*,\qquad\Anorm{\rho}{a_{\mathbb T}}{r_{d_0}}\le E,\qquad\mathcal X^\rho(f)\le\sigma_*,\qquad\max_ig_i\le g_*.
\]
Assume the uniform incoming classical norm on $D_{\mathrm{in}}$, the boundary comparison \eqref{ch17:eq:incoming-wall-gate}, and the inequality on the conditional forces of \cref{ch16:thm:controlled-perturbation,ch16:lem:admitted-products}. Set $p(t)=A_{\mathrm w}\log(1/t)$ and require at the $t$-th step, with $t_t=\max_ig_{t,i}$,
\[
 p(t_t)\le p_{t,i}\le p(t_t)+\tfrac12A_{\mathrm w}(J-1)\log M.
\]
Assume the inclusion of the products, the positivity of the rational image denominators, the small-interaction condition, and the conditions on the auxiliary half-disc of \cref{ch17:lem:complete-white-potential,ch17:thm:carried-networks}, with the bound chosen for this whole family of radii. The small-field event uses $p_{t,i}=p(t_t)$. Under these hypotheses the perturbed objective has all of the following properties: the coefficient bounds of the potential with each index fixed in turn, the gap of the constrained Hessians at every stage, and the strict bounds on the image under the reconstruction majorant, on the outer profile times the whole enclosing product; the comparisons with boundary decay of \cref{ch17:thm:perturbed-hessian-coarea}; the potential of \cref{ch17:lem:complete-white-potential}; the expansions without and with insertion of \cref{ch17:thm:carried-networks}, with input support weight $e^{16\abs S}$ and output support weight $e^{4\abs S}$ as in \cref{ch17:lem:decorated-white-log}, and with the same returned rates and radii; the bounds on the generated term, the pullback, and the determinants of \cref{ch17:thm:quantum-ledger}; and the recursion for $\ScalarLedger$ of \cref{ch17:thm:real-scalar-ledger}. These properties are conclusions of the earlier results.

More explicitly, let $x_t=\mathcal X(f_t)$ and $x_t^{\mathrm{cl}}=\mathcal X^{\mathrm{cl}}(P_t)$ be the incoming boundary norms of the quantum and the classical parts at step $t$, the boundary-completion adjustment of the classical part being understood. Before coordinates are fixed, the outgoing boundary norms $y_t$ and $y_t^{\mathrm{cl}}$ satisfy
\[
\begin{split}
 y_t&\le A_J(1+E)+K_Jx_t+C_Jt_t^2(E+x_t)\bigl[p_t^4+(E+x_t)p_t^2\bigr],\\
 y_t^{\mathrm{cl}}&\le A_J^{\mathrm{cl}}+K_J^{\mathrm{cl}}x_t^{\mathrm{cl}},\qquad t_t=\max_ig_{t,i},\quad p_t=\max_ip_{t,i}.
\end{split}
\]
Under the hypotheses on the doubled output analytic radius and on the incoming image at half the radius, $K_J=4(\log2)K_{\mathrm{col},J}$, where $K_{\mathrm{col},J}$ is the column bound of \cref{ch17:thm:quantum-ledger} for the perturbed zero-fluctuation reconstruction, uniform on the preliminary perturbation ball. The recursion for $\ScalarLedger$ is \cref{ch17:thm:real-scalar-ledger}.

Whenever coordinates are fixed, fix every label with reduced radius and keep the remaining unfixed labels at distance at least $s_{\mathrm{er}}$ from the previous boundary. For $\bullet\in\{\mathrm{cl},\mathrm q\}$, with $(\rho_{\mathrm q},E_{\mathrm q})=(\rho,E)$ and $(\rho_{\mathrm{cl}},E_{\mathrm{cl}})$ the classical radius and the classical torus bound, require
\[
 C_{\mathrm{bc},\bullet}E_\bullet\rrecord/\rho_\bullet\le\sigma_*/4,\qquad\abs{c_j^\sharp}\le\min\{\rrecord,\phi_\bullet(j)/4\}.
\]
Here $C_{\mathrm{bc},\bullet}\ge8$ is fixed before $\rrecord$ and the later widths, and includes any finite sum over types of boundary-completion differences and over conversions between coordinates; a single simultaneous restriction of the union of all newly frozen labels to their fixed values uses $C_{\mathrm{bc},\bullet}=8$ in the coordinates in which the labels were introduced. Every frozen label, whether in the boundary layer, in a protected neighborhood, or among the boundary-crossing coordinates, and including the labels frozen at initialization, is covered by this bound. Here $c^\sharp$ is the value in the analytic coordinates in which the label was introduced; any conversion from group values or transported coordinates is included in the bound by $\rrecord$ beforehand. Assume $x_0,x_0^{\mathrm{cl}}\le\sigma_*$. Then, including the new boundary-completion differences of the torus parts,
\[
 x_{t+1}\le z_t+2e^{-\alpha(s_{\mathrm{er}}-1)}y_t,\qquad x_{t+1}^{\mathrm{cl}}\le z_t^{\mathrm{cl}}+2e^{-\alpha(s_{\mathrm{er}}-1)}y_t^{\mathrm{cl}},
\]
where $z_t$ and $z_t^{\mathrm{cl}}$ are the norms of the boundary-completion differences of the torus parts created when coordinates are fixed at the end of step $t$. The choices can be made with
\[
 x_0,x_0^{\mathrm{cl}}\le\sigma_*,\qquad z_t,z_t^{\mathrm{cl}}\le\sigma_*/4,
\]
and with the last terms of both recursions at most $\sigma_*/2$, uniformly on the input ball. In particular
\[
 x_t,x_t^{\mathrm{cl}}\le3\sigma_*/4\qquad(t\ge1)
\]
after any finite number of steps, with no further loss of radius or rate. For $\abs{v_i}\le\phi(i)$ the oscillation of the centered effective interaction satisfies
\[
 \abs{R_t(v;\Rec)-R_t(0;\Rec)}\le\rho^2\Bigl[E\abs{\Omega_t}+\sigma_*\sum_{j\in\Omega_t}e^{-\alpha d_{\Omega_t}(j)}\Bigr],
\]
where, for this count and this sum, $\Omega_t$ denotes the set of unfixed scalar labels, directions and colors included. At the output of a step that is not followed by fixing coordinates, $\sigma_*$ in this formula is replaced by the fixed output bound for $y_t$.

All constants have the dependence described in item (3) of \cref{ch01:conv:constants} and exclude the later widths. The output bounds are computed first and $s_{\mathrm{er}}$ afterwards; $J$ is not enlarged by means of these constants for fixed $J$. The theorem concerns admissible data; that the configurations of the large-field decomposition are covered by them is proved separately in \cref{ch18:chapter}. For a composition of $s<J$ averaging steps the estimates at each stage hold on their admissible products, but a conclusion about the two images or about the enclosing product also requires the separate image condition for that shorter composition; the preservation of the profile is asserted only for full steps. An empty set of unfixed labels has nonconstant norm zero, determinant one, logarithmic determinant zero, and fluctuation integral one.
\end{theorem}

\begin{proof}
The proof has three parts. Steps 1 to 3 prove the bounds on $y_t$ and $y_t^{\mathrm{cl}}$ before coordinates are fixed. Steps 4 to 10 make the parameter choices in the order summarized in \cref{ch17:tab:parameters} and check that each choice uses only quantities fixed before it. Steps 11 and 12 prove the iteration and the bound on the oscillation.

\begin{table}[t]
\centering
\small
\begin{tabular}{@{}p{0.55\textwidth}p{0.38\textwidth}@{}}
\toprule
Quantity & Chosen after\\
\midrule
$K_*$, torus support weights, and margins & fixed data\\
$J$ & $K_*$ and the torus data\\
$\rho,\rho_{\mathrm{cl}},\Rraw$ & $J$\\
$a,\alpha$ & $J$ and the exponents of the coefficients of the potential\\
$\sigma_0$, $r_q$ & radii and rates\\
$A,C_0,C_1,C_{\mathrm{en},J},\mathcal L_V$, and all output, tensor, and image constants & $\sigma_0$\\
$\Rwhite$, then $\sigma_*$ & these constants\\
$\rho_{\partial,\bullet}$, $w_{\mathrm{bd}}$ & these constants\\
$Y,Y^{\mathrm{cl}}$, then $s_{\mathrm{er}}$ & $\sigma_*$\\
$C_{\mathrm{bc},\bullet}$, $C_{\mathrm{word},J}$, $C_{\mathrm{conv},J}$ & the larger charts\\
$\rrecord$, then $r$, then $\Ren$ & $\sigma_*$ and the conversion constants\\
$g_*$ & all of the above, and $A_{\mathrm w}$\\
\bottomrule
\end{tabular}
\caption{The order of the parameter choices in the proof of \cref{ch17:thm:controlled-kernel}. The requirements are stated in Steps 4 to 10.}\label{ch17:tab:parameters}
\end{table}

\paragraph{Step 1: the hypotheses at each step.}
We fix $A_{\mathrm w}\ge1$ before the coupling bound. The coefficient constants of the potential and the torus and geometric conditions used to choose $J$ do not depend on $A_{\mathrm w}$; only the final coupling bound may depend on it. At each step, the incoming classical norm is the uniform norm on $D_{\mathrm{in}}$, the boundary comparison \eqref{ch17:eq:incoming-wall-gate} holds by assumption, with all label positions and volume factors on the lattices where the labels were introduced, and the conditional forces satisfy \eqref{ch17:eq:conditional-force-bound} and \eqref{ch17:eq:sequential-gate}. The strict bounds on the two images place all segments of evaluation inside the uniform incoming products. Thus, at the fixed frozen value $b=(\xi,\Rec)$, both the perturbation and its comparison with zero satisfy the hypotheses of \cref{ch16:thm:controlled-perturbation}.

\paragraph{Step 2: the classical comparison before coordinates are fixed.}
Let $a_R$ and $a_T$ be the two admissible zero-fluctuation chains before the gauge transformation, the regional one with its missing slots set to zero. The comparison of values, before the common factor $M^J$, is
\begin{align*}
 F_R(v,a_R)-F_T(v,a_T)={}&F_R(v,a_R)-F_T(v,a_R)\\
 &+\int_0^1D_zF_T\bigl(v,a_T+\tau(a_R-a_T)\bigr)[a_R-a_T]\,d\tau.
\end{align*}
The fixed larger charts contain this segment, since the smaller solution neighborhoods have strict room in them. The first difference has the representation \eqref{ch17:eq:action-word-difference} through changed formulas. On unfixed free slots, the difference of the chains in the integral is the comparison of minimizers in \cref{ch17:thm:perturbed-hessian-coarea}, here with perturbation zero; on missing free slots it is a bounded source in $\Omega^c$. Every nonconstant coefficient with an unfixed output index therefore has a network to $\Omega^c$, by the action tensors rooted at every index and \cref{ch15:lem:wall-networks}. This includes derivatives falling on the action tensor while the source stays undifferentiated. We sum over the indices in $\Omega^c$ from the unfixed index using \eqref{ch17:eq:density-sum}. Pure constants supported in $\Omega^c$ stay in the classical value. The conversion of coefficients and profiles in the proof of \cref{ch17:thm:carried-networks}, without its expansions, now yields a finite bound $A_J^{\mathrm{cl}}$ for the centered comparison. The factor $M^J$ depends only on $J$ and is included in this constant. No interpolation between different fibers is used in this calculation of values.

\paragraph{Step 3: the perturbation and the quantum comparison.}
Next we interpolate in the adjusted input perturbation. By the envelope formula, the derivative of the minimum value is the evaluation of the perturbation at the minimizer. The strict bound on the image under the reconstruction majorant bounds its input jets, and joining its input boundary network to the reconstruction networks proves the output bound $K_J^{\mathrm{cl}}x^{\mathrm{cl}}$. Equivalently, this is the conclusion of \cref{ch16:thm:controlled-perturbation} on the minimum values for two perturbations, with its exact boundary-completion adjustment. The fixed small boundary-completion difference is included with the unperturbed term in $A_J^{\mathrm{cl}}$. This proves $y^{\mathrm{cl}}\le A_J^{\mathrm{cl}}+K_J^{\mathrm{cl}}x^{\mathrm{cl}}$. The coefficient networks in this argument are the classical ones, with the common rates chosen in \cref{ch17:thm:carried-networks}.

We subtract the matched torus recursion from the recursion of \cref{ch17:thm:quantum-ledger}. The difference is the sum of the comparison of the generated terms and torus pullbacks, the term $\mathsf P^\Omega f$, and the compared fluctuation correction. The first two are bounded in \cref{ch17:thm:quantum-ledger} by $A_J(1+E)$ and $K_Jx$, and the third by \cref{ch17:thm:carried-networks}. By the triangle inequality, we obtain the stated bound on $y$.

\paragraph{Step 4: radii, rates, and the perturbation bound.}
We fix $K_*$ before $J$ and choose $J$ for the torus conditions and the geometric condition. We choose the base radii inside the fixed larger construction domains; one may choose $\rho=\rho_{\mathrm{cl}}$, and different fixed ratios simply keep both image conditions. We choose the common rates below the intrinsic coefficient and unperturbed image exponents at fixed $J$, and impose the separate output-rate conditions against $Ba$ and $B\alpha$ for the inherited perturbation. Its controlled coefficient difference preserves the half-radius image at these working rates after the final perturbation bound is decreased. We decrease $\Rraw$ and the base radii if necessary, for the fixed margin of the larger construction domain and for $C_{T,\bullet}K_*\rho_\bullet\le1/4$. The torus conditions are used with their stronger torus rate, not through an estimate involving an uncontrolled regional constant.

We choose a preliminary bound $\sigma_0$ below the admissible inverse, conditional force, and real gap margins, and a free radius $r_q$ with
\[
 C_{\mathrm{seq},J}(2I_0C_P\sigma_0\rho_{\mathrm{cl}}+r_q)<\varepsilon_{\mathrm{tube}}/2.
\]
The constants are those of the larger construction neighborhoods in \eqref{ch17:eq:conditional-force-bound}, whose uniformity in all higher parameters was checked there. We compute all output, tensor, image, sequential, and enclosure constants uniformly for this bound on the fixed larger products. In particular we fix $A$, $C_0$, $C_1$, $C_{\mathrm{en},J}$, and $\mathcal L_V$, and the difference-row constants $C_{P,\bullet}$ at the allocated working rates, and choose $\Rwhite$ from the positive minimum in \cref{ch17:lem:complete-white-potential}. Then we fix $\sigma_*\le\sigma_0$ so that, for both radii,
\[
 \sigma_*\rho_{\mathrm{cl}}\le\lambda_*\Rwhite/(6C_{\mathrm{en},J}),\qquad
 C_{P,\bullet}\sigma_*\rho_{\mathrm{cl}}\le\rho_\bullet/64.
\]
This order is necessary because this summand of the enclosure remains when $r$ and $\rrecord$ are decreased.

\paragraph{Step 5: the boundary radii.}
We choose the two reduced boundary radii so that the bounds on the image at zero fluctuation near the boundary in \cref{ch16:lem:strict-two-radius-images} are at most $\rho_\bullet/16$. Then we choose a fixed width $w_{\mathrm{bd}}$ of the boundary layer so that the bounds on the tails far from the boundary are at most $\rho_\bullet/16$. The remaining bounds at zero fluctuation are the torus bulk bound and the bound for the frozen values. The first is a torus condition, and the second will be imposed on $\rrecord$. All constants in these choices were computed independently of these widths and reduced radii.

\paragraph{Step 6: output bounds and the width of the frozen layer.}
We set
\begin{equation}\label{ch17:eq:logarithmic-envelope}
 c_J^{\log}=\tfrac12A_{\mathrm w}(J-1)\log M,\qquad p^+(t)=A_{\mathrm w}\log(1/t)+c_J^{\log},\qquad H_*=E+\sigma_*.
\end{equation}
By the hypothesis on the radii, we have $1<p(t)\le p_{t,i}\le p^+(t)$ when $t<e^{-1}$. On the input ball we impose provisionally the requirement
\[
 C_Jt^2H_*\bigl[(p^+(t))^4+H_*(p^+(t))^2\bigr]\le1,
\]
which holds for small couplings. We then take the fixed output bounds
\[
 Y=A_J(1+E)+K_J\sigma_*+1,\qquad Y^{\mathrm{cl}}=A_J^{\mathrm{cl}}+K_J^{\mathrm{cl}}\sigma_*.
\]
We choose the integer width of the frozen layer with
\[
 s_{\mathrm{er}}\ge\max\Bigl\{2,\ w_{\mathrm{bd}},\ 1+\alpha^{-1}\log\bigl(4\max\{Y,Y^{\mathrm{cl}},1\}/\sigma_*\bigr)\Bigr\}.
\]
The bound may be enlarged to account for the fixed buffer between label positions in \cref{ch15:lem:freezing}. Each output restricted to the interior of the frozen layer is then at most $\sigma_*/2$, since $2e^{-\alpha(s_{\mathrm{er}}-1)}\max\{Y,Y^{\mathrm{cl}}\}\le\sigma_*/2$.

\paragraph{Step 7: the frozen values and the new boundary-completion differences.}
Whenever coordinates are fixed, let $F$ be the union of all newly frozen analytic labels: the whole zone with reduced radius, the protected neighborhoods, and the boundary-crossing coordinates. The same convention covers the restriction at initialization. In the coordinates in which the labels were introduced, a single application of \cref{ch15:lem:freezing} to this union bounds the torus boundary-completion difference by $8E_\bullet\rrecord/\rho_\bullet$. The rate hypothesis of part (2) of that lemma holds with the torus rate $a_{\mathbb T}$ and the regional rates $a,\alpha$, since $\sqrt3\,(a+\alpha)<a_{\mathbb T}$ by \cref{ch17:thm:carried-networks}. No separate factor arises for each component of $\Omega^c$. If a fixed finite list of types of boundary-completion differences is kept separately, with conversion factors $L_{k,\bullet}$, the triangle inequality yields instead the bound with $C_{\mathrm{bc},\bullet}=8\sum_kL_{k,\bullet}$. We include the classical boundary-completion adjustment among these types whenever it is recorded separately. These constants are fixed on the larger charts, before $\rrecord$ and all later widths.

The bound concerns the analytic values $c_j^\sharp$ in the coordinates in which the labels were introduced. For a frozen group element in the identity chart, $\abs{\Log U}=\groupdist(U)$ in the Euclidean norm of the Lie algebra, so this conversion has constant one. A longer fixed chart formula, in the coordinates before the gauge transformation or in transported coordinates, has a finite bound $C_{\mathrm{word},J}$ on its derivative rows on the larger construction domain. Integrating along the input segment, we bound its logarithm in the identity chart by $C_{\mathrm{word},J}$ times the maximal input norm. These conversion constants are fixed before the bounds on the frozen values are chosen. If a separate estimate for the initial frozen values is provided with a constant $C_{\mathrm{ad},J}$, we choose $C_{\mathrm{conv},J}\ge\max\{1,C_{\mathrm{ad},J},C_{\mathrm{word},J}\}$ before using a condition relating the support radius to the frozen values. The theorem assumes the bound on every $c_j^\sharp$. It does not deduce this bound from a small support radius in an arbitrary gauge representative, and it transports a norm bound between gauge frames for the frozen coordinates only when the change of frame is the norm-preserving adjoint action of a block.

We choose $\rrecord>0$ below the bound for the frozen values in the image, below the enclosure bound $\lambda_*\Rwhite/(6C_{\mathrm{en},J})$, and, for both $\bullet\in\{\mathrm{cl},\mathrm q\}$, below
\begin{equation}\label{ch17:eq:saved-straddle-budget}
 \min\Bigl\{\rho_{\partial,\bullet}/4,\ \rho_\bullet/4,\ \frac{\sigma_*\rho_\bullet}{4C_{\mathrm{bc},\bullet}(E_\bullet+1)}\Bigr\}.
\end{equation}
These are upper bounds on a quantity chosen after their constants, and they remain positive when a bound $E_\bullet$ is zero. They imply the two inequalities $C_{\mathrm{bc},\bullet}E_\bullet\rrecord/\rho_\bullet\le\sigma_*/4$ in the statement. The hypothesis $\abs{c_j^\sharp}\le\min\{\rrecord,\phi_\bullet(j)/4\}$ is imposed on all of $F$, not only on the old boundary zone.

\paragraph{Step 8: initialization.}
The assumption $x_0,x_0^{\mathrm{cl}}\le\sigma_*$ of the iteration can be met: starting from the torus reference, both boundary parts are zero before initialization, and the same restriction of the union to fixed values, together with \eqref{ch17:eq:saved-straddle-budget}, gives initial boundary parts of norm at most $\sigma_*/4$. For other prescribed inputs the bounds on the initial norms are kept as hypotheses. The coverage of the configurations and the estimate for the initial frozen values are separate inputs of the later applications.

\paragraph{Step 9: the support radius.}
We choose $r>0$ inside all margins of the small-link domain and of the larger construction domain, below $\lambda_*\Rwhite/(6C_{\mathrm{en},J})$, and, when an estimate for the frozen values in terms of the support radius is provided, with $C_{\mathrm{conv},J}r\le\rrecord$, imposed after the conversion constant has been fixed. Then $\Ren=2C_{\mathrm{en},J}(r+\rrecord+\sigma_*\rho_{\mathrm{cl}})/\lambda_*$ is positive and at most $\Rwhite$. Both image denominators and the curvature bound on the whole product now hold by \cref{ch17:lem:complete-white-potential}.

\paragraph{Step 10: the coupling bound.}
Finally, we choose $g_*<e^{-1}$ so that the following inequalities hold for every $0<t\le g_*$, with the largest constant $C_J$ of \cref{ch17:lem:decorated-white-primitives} and the fixed $\epsilon_J,c_{z,J}$ of \cref{ch17:thm:carried-networks}:
\begin{equation}\label{ch17:eq:uniform-logarithmic-gates}
\begin{gathered}
 tp^+(t)<\Ren,\qquad C_Jtp^+(t)\le\tfrac12,\\
 2C_JtD^+(t)\le\epsilon_J,\qquad tD^+(t)\le c_{z,J}/2,\qquad D^+(t)=(p^+(t))^3+(1+H_*)p^+(t),\\
 C_0t^2\mathcal L_V\le1/16,\qquad C_1t(\mathcal L_V+2)<1/(4\sqrt3),\\
 C_Jt^2H_*\bigl[(p^+(t))^4+H_*(p^+(t))^2\bigr]\le1.
\end{gathered}
\end{equation}
The image denominators are at least $1/2$ on the whole enclosing product by the earlier image bounds, which do not involve the coupling. For any radii in the stated logarithmic range we have $p_t\le p^+(t)$ and $E+x_t\le H_*$. The first line implies the inclusion of the products. The next line implies $\mathfrak s(t)\le\epsilon_J$ and $t\le c_{z,J}/(2[p_t^3+p_t+(E+x_t)p_t])=r_z/2$. The third line implies \eqref{ch17:eq:coupling-budgets}; since $p(t)>1$, it also implies the condition $C_1t(\mathcal L_V+2)\le p(t)/(2\sqrt3)$ on the means in \cref{ch17:lem:ambient-white-comparison}. The last line is the provisional requirement of Step 6. The choice is simultaneous, because every left side that has to be small is a fixed combination of $t\log^q(1/t)$ and $t^2\log^q(1/t)$ for finitely many integers $q$. Writing $t=e^{-u}$, we have $e^{-u}u^q\to0$; for example the power series of $e^{u/2}$ bounds $u^qe^{-u}$ by a constant times $e^{-u/2}$. These upper bounds hold for all sufficiently small positive $t$, so a single bound works for all scales of every full step in the prescribed family of radii. The common cutoff has $p_{t,i}=p(t)$, and the cutoffs at the individual scales, $p(g_i)=p(t)+A_{\mathrm w}\log(1/\lambda_i)$, lie in the same range. The expansions are used on these chosen products, not on the enclosing product $\mathcal K_{\mathrm{en}}$. This final choice does not change the output bounds or the width $s_{\mathrm{er}}$, which were computed before.

\paragraph{Step 11: the iteration.}
By \cref{ch15:lem:freezing} with $\eta=1/4$, fixing the coordinates of the union $F$ bounds the old boundary part by $2e^{-\alpha(s_{\mathrm{er}}-1)}y_t$, since $\abs{c_j^\sharp}\le\phi_\bullet(j)/4$ and all labels with reduced radius are frozen. Independently, the boundary-completion differences of the torus parts satisfy, by the hypothesis on the frozen values,
\begin{equation}\label{ch17:eq:creation-straddle-bound}
 z_t^\bullet\le C_{\mathrm{bc},\bullet}E_\bullet\rrecord/\rho_\bullet\le\sigma_*/4,
\end{equation}
where $z_t^{\mathrm q}=z_t$. Fixing coordinates at an additional interior component of $\Omega^c$ satisfies the same estimate, because its values are also bounded by $\rrecord$. Combining the two bounds, we obtain
\[
 x_{t+1}\le z_t+2e^{-\alpha(s_{\mathrm{er}}-1)}y_t,\qquad x_{t+1}^{\mathrm{cl}}\le z_t^{\mathrm{cl}}+2e^{-\alpha(s_{\mathrm{er}}-1)}y_t^{\mathrm{cl}}.
\]
These estimates hold for arbitrary additional components of $\Omega^c$, by \cref{ch15:lem:freezing}. If both input norms are at most $\sigma_*$, then $y_t\le Y$ and $y_t^{\mathrm{cl}}\le Y^{\mathrm{cl}}$, and by Step 6 the right sides are at most $3\sigma_*/4$. Starting from the initial bounds, finite induction proves the bounds on $x_t,x_t^{\mathrm{cl}}$ after every number of steps. The older frozen values are simply fixed in the next application. The remaining labels keep their input radii because the whole zone with reduced radius was frozen: the conclusion of \cref{ch15:lem:freezing} is the norm with the uniform radius on the remaining labels. This verifies the hypothesis of a uniform incoming profile at the next step. The classical boundary-completion adjustment satisfies
\[
 \mathcal X^{\mathrm{cl}}(\bar P)\le\mathcal X^{\mathrm{cl}}(P)+\mathcal X^{\mathrm{cl}}(S_\xi^{\mathrm{cl}})\le\sigma_*+\sigma_*/4<2\sigma_*,
\]
by the included bound on the boundary-completion differences. The inclusion of the old fixed complement and the positions of the labels in the next geometry are hypotheses of that step, not conclusions of this norm inequality. The inequality on the conditional forces then holds at every coarser parameter of the next step, by \eqref{ch17:eq:conditional-force-bound}. Every operation preserves the Borel dependence on the frozen values, as proved in \cref{ch17:thm:perturbed-hessian-coarea}. No analytic loss is summed over earlier steps.

\paragraph{Step 12: the oscillation.}
We write the centered effective interaction as its torus part plus its boundary part. At each degree we assign each monomial to its roots with the shares $\nu_i/\abs\nu$, whose sum is one. Evaluating at $\abs{v_i}\le\phi(i)$ and dropping the network weight, we obtain for the boundary part
\[
 \abs{f_t(v)-f_t(0)}\le\rho^2\mathcal X(f_t)\sum_{i\in\Omega_t}e^{-\alpha d_{\Omega_t}(i)}.
\]
The degree-one term is included in this assignment. For the torus part the same assignment and its seminorm give $\rho^2E\abs{\Omega_t}$. These are the stated bounds on the oscillation; before coordinates are fixed we use $Y$ instead of $\sigma_*$.

The preceding results, with the choices just made, imply every assertion of the theorem. All local constants were evaluated before the later widths, and none of them was used to choose $J$. A composition of fewer than $J$ averaging steps keeps its estimates at each stage on its admissible products, but its conclusions about the full radius require the separate image condition. The iteration for $\ScalarLedger$ and for the norms uses only full steps. A final composition of fewer than $J$ averaging steps may remain as a deterministic function; no recursion for $\ScalarLedger$ is asserted for it here. The assertions for an empty set of unfixed labels are the usual conventions for empty sums, products, and determinants. This completes the proof.
\end{proof}

\section{Summation of truncation errors}\label{ch17:sec:endpoint}

\begin{lemma}\label{ch17:lem:finite-ceiling-endpoint-gates}
Fix $0<\Gamma<\infty$, $M>1$, a finite set $\mathcal R$ of ratios $R>1$, a finite set of positive numbers $\eta$, and finitely many positive constants $a$, $c$, $C$, and $\varepsilon$.
\begin{enumerate}
\item For every finite list of nonnegative integers $k$, there is a minimum side $m_0$, independent of $0<\gamma\le\Gamma$, such that
\[
 C[1+\log(m/\gamma)]^k\le\varepsilon m/\gamma\qquad(m\ge m_0).
\]
\item With $z=(1/2)\log(Mm/\gamma)$, the functions
\[
 X_c(m,\gamma)=m^3e^{-cm/\gamma},\qquad W_a(m,\gamma)=m^3e^{-az^2}
\]
are decreasing in $m$ and increasing in $\gamma$ for $m\ge m_0$ and $0<\gamma\le\Gamma$, after increasing $m_0$, and they satisfy $X_c(Rm,\gamma)\le X_c(m,\gamma)/2$ and $W_a(Rm,\gamma)\le W_a(m,\gamma)/2$ for $R\in\mathcal R$. Their sums over the sides $mR^k$, $k\ge0$, are at most twice their first terms, and tend to zero as $m\to\infty$, uniformly on $(0,\Gamma]$.
\item Let $m_k=mR^k$, $k\ge0$, for some $R\in\mathcal R$, together with at most one additional side $\eta mR^j$, and let nonnegative errors $\delta$ and $z_{\mathrm{ext}}$ at these sides satisfy $\delta\le CW_a$ and $z_{\mathrm{ext}}\le CX_c$. After increasing $m_0$, depending only on $\Gamma$, $M$, $\mathcal R$, the set of $\eta$, and the constants, the following holds uniformly for $m\ge m_0$, $0<\gamma\le\Gamma$, and the number of sides, finite or infinite. Every $\delta$ is at most $1/2$, every $z_{\mathrm{ext}}$ is at most $1/4$, and $\epsilon=\delta/(1-z_{\mathrm{ext}})$ satisfies $\epsilon\le(4/3)\delta\le1/2$. The sums of the $\delta$, of the $z_{\mathrm{ext}}$, and of the $\epsilon$ are finite, and
\[
 \log\prod(1-\delta)^{-1}\le2\sum\delta,\qquad\log\prod(1-\epsilon)^{-1}\le2\sum\epsilon,
\]
also for every finite subproduct. These sums are bounded by any prescribed positive number after increasing $m_0$, and they tend to zero as $m\to\infty$, uniformly on $(0,\Gamma]$.
\end{enumerate}
In particular, part (3) of \cref{ch17:lem:ambient-white-comparison} holds. All constants in this lemma have the dependence described in item (4) of \cref{ch01:conv:constants}; the local constants of \cref{ch17:lem:ambient-white-comparison} keep their own dependence.
\end{lemma}

\begin{proof}
Let $x=m/\gamma\ge m_0/\Gamma$. For every fixed $k$ we have $(1+\log x)^k/x\to0$ as $x\to\infty$: with $y=\log x$, the exponential series bounds $y^ke^{-y/2}$ on a half-line $y\ge y_0$. This proves part (1), for finitely many $k$ simultaneously, and allows us to assume $x\ge e$ throughout.

For part (2), let $\Delta_R=(1/2)\log R>0$. Since $z(Rm)=z(m)+\Delta_R$, the ratios are
\begin{align}
 \frac{W_a(Rm,\gamma)}{W_a(m,\gamma)}&=R^3e^{-a(2\Delta_Rz+\Delta_R^2)},\label{ch17:eq:white-ratio}\\
 \frac{X_c(Rm,\gamma)}{X_c(m,\gamma)}&=R^3e^{-c(R-1)m/\gamma}.\label{ch17:eq:exterior-ratio}
\end{align}
Their maxima over $m\ge m_0$ and $0<\gamma\le\Gamma$ occur at the smallest $m$ and the largest $\gamma$, and they tend to zero as $m_0$ increases, since $z\ge(1/2)\log(Mm_0/\Gamma)\to\infty$ and $m/\gamma\ge m_0/\Gamma\to\infty$. We choose $m_0$ so that both ratios are at most $1/2$ for every $R\in\mathcal R$. The logarithmic derivatives in $\log m$ are $3-cm/\gamma$ for $X_c$ and $3-az$ for $W_a$, and those in $\log\gamma$ are $cm/\gamma$ and $az$. Hence $X_c$ and $W_a$ are increasing in $\gamma$, and they are decreasing in $m$ once $cm_0/\Gamma>3$ and $(a/2)\log(Mm_0/\Gamma)>3$. The logarithms of the first terms at $\gamma=\Gamma$ are $3\log m-cm/\Gamma$ and $3\log m-a\log^2(Mm/\Gamma)/4$, and both tend to $-\infty$. By monotonicity in $\gamma$, the first terms tend to zero uniformly on $(0,\Gamma]$. Summing the series with ratio at most $1/2$ shows that the sums over $mR^k$ are at most twice their first terms, which proves part (2).

For part (3), the additional side $\eta mR^j$ is handled in the same way. Let $m_{\mathrm{low}}=\min\{1,\min\eta\}m$. For sufficiently large $m$, monotonicity in $m$ bounds the term at the additional side by the corresponding function at $m_{\mathrm{low}}$, whether or not $\eta<1$. Since $m_{\mathrm{low}}$ is a fixed positive multiple of $m$, the sums over all sides are at most three times the first terms evaluated at $m_{\mathrm{low}}$, and these tend to zero uniformly on $(0,\Gamma]$ as $m\to\infty$, for finitely or infinitely many sides. We take the maximum of the finitely many minimum sides so obtained. Increasing $m_0$ once more, we make every $\delta$ at most $1/2$, every $z_{\mathrm{ext}}$ at most $1/4$, and the sums smaller than any prescribed positive number. Then $\epsilon=\delta/(1-z_{\mathrm{ext}})\le(4/3)\delta$, so the sum and the maximum of the $\epsilon$ can be made as small as desired as well, in particular every $\epsilon$ at most $1/2$. For $0\le u\le1/2$,
\[
 -\log(1-u)=\int_0^u(1-v)^{-1}\,dv\le2u.
\]
Summing this inequality bounds the logarithms of all reciprocal products, including every finite subproduct; the infinite products are defined by the increasing finite sums of logarithms. This proves part (3).

Part (3) of \cref{ch17:lem:ambient-white-comparison} is the case $R=B$, with $C=C_J$, $a=a_{\mathrm w}$, $c=b$, and the finite set of additional ratios, since $\log^2\sqrt{Mm/\gamma}=z^2$; its bound $\epsilon_t\le2\delta_t$ follows from $\epsilon_t\le(4/3)\delta_t$. The local concentration statements of \cref{ch17:lem:ambient-white-comparison} use only its hypotheses on the products and on the coupling at the given step; the restriction to a coupling bound enters only in the summation just proved, so those local statements and their prefactors are unchanged.
\end{proof}
